\documentclass[14pt,a4paper,oneside]{book} % Dùng extbook để nhận cỡ chữ 14pt chuẩn
\usepackage[utf8]{vietnam}
\usepackage{amsmath,amssymb,mathrsfs,fancyhdr,enumerate,multirow,makecell,venndiagram}
\usepackage{tkz-euclide,tikz-3dplot,tikz,tkz-tab}
\usetikzlibrary{arrows,calc,intersections,angles,quotes,backgrounds,shapes.geometric,patterns,through}
\usepackage{yhmath}
\usepackage{pgfplots}
\pgfplotsset{compat=1.9}
\usepgfplotslibrary{fillbetween}
\usepackage[tikz]{bclogo}
\usepackage[top=1.5cm, bottom=1.5cm, left=1.5cm, right=1.5cm]{geometry}
\usepackage[hidelinks,unicode]{hyperref}

% Gói ex_test (Chọn [dethi] để in đề hoặc [loigiai] để in cả lời giải)
\usepackage[loigiai]{ex_test}
%\usepackage[dethi]{ex_test}

\def\colorEX{\color{blue}}

\usepackage{fontawesome}
\usepackage{framed}
\usepackage[many]{tcolorbox}
\usepackage{polynom}

\colorlet{tcbcol@back}{tcbcolback}
\colorlet{tcbcol@frame}{tcbcolframe}
\newtcolorbox[auto counter]{danggg}[1]{
	fonttitle=\bfseries\upshape,
	colframe=blue,colback=blue!10,
	sharp corners, breakable, 
	title={\color{white}\faFolderO \ DẠNG~\thetcbcounter.\ #1}
	\addcontentsline{toc}{subsection}{Dạng~\thetcbcounter. #1}
}
\newenvironment{dang}[1]{
	\begin{danggg}{#1}
		\textit{Phương pháp giải:~}
	}{\end{danggg}}
\newenvironment{tomtat}{
	\begin{tcolorbox}[boxrule=1pt,arc=0mm,breakable]
	}{\end{tcolorbox}}

\newtheorem{bt}{Bài}
\newtheorem{vd}{Ví dụ}
\theorembodyfont{\it}
\newtheorem{dl}{ Định lí}
\theoremseparator{.}
\theorembodyfont{\rm}
\newtheorem{hq}{ Hệ quả}
\newtheorem{tc}{ Tính chất}
\newtheorem{dn}{ Định nghĩa}
\theoremstyle{nonumberplain}
\newtheorem{kn}{ Khái niệm}
\theoremseparator{.}
\theorembodyfont{\it}
\usepackage{stackengine}
\usepackage{scalerel}

\newcommand\dangersign[1][2ex]{%
	\renewcommand\stacktype{L}%
	\scaleto{\stackon[1.3pt]{\color{red}$\triangle$}{\fontsize{4}{0}\selectfont !}}{#1}%
}
\renewcommand{\parallel}{\mathbin{\!\raisebox{0.25pt}/\mkern-2mu\raisebox{0.25pt}/\!}}
\renewcommand{\backsim}{\ \mathord{\raisebox{0.025\baselineskip}{\rotatebox{90}{\fontfamily{qag}\selectfont S}}}\ }
\theoremseparator{}
\newtheorem{note}{{\dangersign[3ex]}}
\newtheorem{nx}{Nhận xét.}
\newtcolorbox{luuy}[1][]{colback=yellow!20,title={\bfseries Chú ý},#1}

% Tiêu đề chương
\setcounter{secnumdepth}{4}
\setcounter{tocdepth}{2}
\usepackage{titlesec}
\usepackage{titletoc}

\makeatletter
\renewcommand\thesection{\@arabic\c@section}
\renewcommand\thesubsection{\@Alph\c@subsection}
\renewcommand\thesubsubsection{\@arabic\c@subsubsection.}
\makeatother

\titleformat{\section}{\centering\Large\bfseries}{Chủ đề~\thesection.}{-5ex}{
	\setcounter{vd}{0}
	\setcounter{ex}{0}
	\setcounter{bt}{0}
	\setcounter{dl}{0}
	\setcounter{dn}{0}
	\setcounter{hq}{0}
	\setcounter{tc}{0}
	\MakeUppercase
}
\titleformat{\subsection}{\Large\bfseries}{{\thesubsection.}}{2ex}{\MakeUppercase}
\usepackage[final]{pdfpages}

% Cài đặt riêng
\def\hideboxEX{}
\renewcommand{\boxEX}[2][15pt]{%
	\setbox1=\hbox{#2\strut}
	\def\sizeboxEX{#1}%
	\ \EXbox{\hideboxEX#2\strut}\ %
}

\renewcommand\arraystretch{2.4} 
\renewcommand{\tabcolsep}{4mm}

% Dòng kẻ làm bài
\renewcommand{\dotlineEX}[1]{
	\setlength{\parskip}{0.8\baselineskip}%
	\setlength{\parindent}{0pt}%
	\def\numlinedot{#1}
	\par
	\foreach \dotline in{1,...,\numlinedot}
	{
		\noindent\color{gray}{\tiny \dotfill}
		\par
	}\vspace*{-0.8\baselineskip}
	\begin{flushright}
		\hfill{\color{violet!33}\ensuremath{\blacksquare}}
	\end{flushright}
}

% Nút Casio
\renewcommand\boxed[1]{%
	\tikz[baseline=(key.base)]
	\node[%
	draw,
	fill=black!40,
	drop shadow={shadow xshift=0.25ex,shadow yshift=-0.25ex,fill=black!90,opacity=0.75},
	rectangle,minimum width=0.6cm,minimum height=0.6cm,
	rounded corners=2pt,
	inner sep=1pt,
	line width=0.5pt,
	font=\sffamily
	](key) {\color{white}#1\strut}
	;
}

\newcommand{\hoac}[1]{\left[\begin{aligned}#1\end{aligned}\right.}
\newcommand{\heva}[1]{\left\{\begin{aligned}#1\end{aligned}\right.}

\pagestyle{fancy}
\fancyhf{} % Xóa sạch toàn bộ header/footer cũ
\renewcommand{\headrulewidth}{0pt}

\newcommand{\caidat}{2}
\newcommand{\indapan}[2]{
	\begin{center}
		\color{blue!50!green} \fbox{\fontfamily{qag}\bfseries\selectfont BẢNG ĐÁP ÁN}
	\end{center}
	{\fontfamily{lmtt}\bfseries\selectfont \inputansbox[\caidat]{#1}{#2}}
}

\everymath{\color{blue}}

\begin{document}
	% ==============================================================================
% FILE ĐƯỢC XUẤT TỰ ĐỘNG CHO: 1. MAIN\data\Cau1.tex
% Lớp: LỚP 10 | Buổi học: Buổi 2 - Kiểm tra định kỳ
% Tổng số học sinh: 52 | Số mã đề sử dụng: 10
% In 2 mặt (Duplex) để mỗi học sinh là 1 tờ A4 (Trang 1: Mặt trước, Trang 2: Mặt sau)
% ==============================================================================

\dotlinefull{ex}

% ======================================================================
% HỌC SINH: An Nguyên | MÃ: HS026 | MÃ ĐỀ: 004
% ======================================================================
\noindent
\begin{tcolorbox}[colback=white,colframe=black,arc=0mm,boxrule=0.8pt,top=4pt,bottom=4pt,left=6pt,right=6pt]
	\begin{minipage}[c]{0.44\textwidth}
		{\bfseries\large LỚP TOÁN THẦY THÔNG}\\[3pt]
		\textbf{Họ và tên:} \textbf{\large An Nguyên}\\[3pt]
		\textbf{Lớp:} LỚP 10 \quad -- \quad \textbf{Mã HS:} \texttt{HS026}
	\end{minipage}
	\hfill
	\begin{minipage}[c]{0.22\textwidth}
		\begin{tcolorbox}[colback=white,colframe=black,arc=0mm,boxrule=0.6pt,height=1.8cm,halign=center,valign=top,top=2pt]
			{\bfseries\normalsize Điểm}
		\end{tcolorbox}
	\end{minipage}
	\hfill
	\begin{minipage}[c]{0.30\textwidth}
		\raggedleft
		\textit{Thời gian làm bài: 30 phút}\\[4pt]
		\textbf{Mã đề thi:} \fbox{\textbf{\Large\strut\ 004\ }}
	\end{minipage}
\end{tcolorbox}
\vspace{0.25cm}

\noindent\textbf{\large ĐỀ BÀI (MÃ ĐỀ: 004)}
\vspace{0.1cm}

% --- CÂU 1 (TRANG 1) ---
\begin{ex}
	Thực hiện phép tính:
	\begin{enumerate}[a)]
		\item $A = 3x^3y \cdot 4xy^2$
		\item $B = 3x^2y^3 \cdot 2x^2y - 11x^4y^4$
		\item $C = 18x^6y^5 : 3x^2y^2 - 12x^4y^3$
		\item $D = 4x^3y^2 + 8x^5y^3 : 2x^2y$
	\end{enumerate}
	\loigiai{
		\begin{enumerate}[a)]
			\item Ta có:
			\begin{align*}
				A &= 3x^3y \cdot 4xy^2 \\
				&= (3 \cdot 4) \cdot (x^3 \cdot x) \cdot (y \cdot y^2) \\
				&= 12x^4y^3
			\end{align*}
			\item Ta có:
			\begin{align*}
				B &= 3x^2y^3 \cdot 2x^2y - 11x^4y^4 \\
				&= 6x^4y^4 - 11x^4y^4 \\
				&= (6 - 11)x^4y^4 \\
				&= -5x^4y^4
			\end{align*}
			\item Ta có:
			\begin{align*}
				C &= 18x^6y^5 : 3x^2y^2 - 12x^4y^3 \\
				&= 6x^4y^3 - 12x^4y^3 \\
				&= (6 - 12)x^4y^3 \\
				&= -6x^4y^3
			\end{align*}
			\item Ta có:
			\begin{align*}
				D &= 4x^3y^2 + 8x^5y^3 : 2x^2y \\
				&= 4x^3y^2 + 4x^3y^2 \\
				&= 8x^3y^2
			\end{align*}
		\end{enumerate}
	}
\end{ex}

% --- CÂU 2 ---
\begin{ex}
	Thực hiện phép tính:
	\begin{enumerate}[a)]
		\item $A = -5a^2b^3 \cdot 2ab^2$
		\item $B = 4a^3b \cdot 3ab^2 - 17a^4b^3$
		\item $C = 24a^7b^4 : 6a^4b^2 - 9a^3b^2$
		\item $D = 3a^2b^4 + 14a^4b^5 : 2a^2b$
	\end{enumerate}
	\loigiai{
		\begin{enumerate}[a)]
			\item Ta có:
			\begin{align*}
				A &= -5a^2b^3 \cdot 2ab^2 \\
				&= (-5 \cdot 2) \cdot (a^2 \cdot a) \cdot (b^3 \cdot b^2) \\
				&= -10a^3b^5
			\end{align*}
			\item Ta có:
			\begin{align*}
				B &= 4a^3b \cdot 3ab^2 - 17a^4b^3 \\
				&= 12a^4b^3 - 17a^4b^3 \\
				&= (12 - 17)a^4b^3 \\
				&= -5a^4b^3
			\end{align*}
			\item Ta có:
			\begin{align*}
				C &= 24a^7b^4 : 6a^4b^2 - 9a^3b^2 \\
				&= 4a^3b^2 - 9a^3b^2 \\
				&= (4 - 9)a^3b^2 \\
				&= -5a^3b^2
			\end{align*}
			\item Ta có:
			\begin{align*}
				D &= 3a^2b^4 + 14a^4b^5 : 2a^2b \\
				&= 3a^2b^4 + 7a^2b^4 \\
				&= 10a^2b^4
			\end{align*}
		\end{enumerate}
	}
\end{ex}

\setcounterpageref{ex}{0}
\newpage

% ======================================================================
% HỌC SINH: Ánh | MÃ: HS012 | MÃ ĐỀ: 001
% ======================================================================
\noindent
\begin{tcolorbox}[colback=white,colframe=black,arc=0mm,boxrule=0.8pt,top=4pt,bottom=4pt,left=6pt,right=6pt]
	\begin{minipage}[c]{0.44\textwidth}
		{\bfseries\large LỚP TOÁN THẦY THÔNG}\\[3pt]
		\textbf{Họ và tên:} \textbf{\large Ánh}\\[3pt]
		\textbf{Lớp:} LỚP 10 \quad -- \quad \textbf{Mã HS:} \texttt{HS012}
	\end{minipage}
	\hfill
	\begin{minipage}[c]{0.22\textwidth}
		\begin{tcolorbox}[colback=white,colframe=black,arc=0mm,boxrule=0.6pt,height=1.8cm,halign=center,valign=top,top=2pt]
			{\bfseries\normalsize Điểm}
		\end{tcolorbox}
	\end{minipage}
	\hfill
	\begin{minipage}[c]{0.30\textwidth}
		\raggedleft
		\textit{Thời gian làm bài: 30 phút}\\[4pt]
		\textbf{Mã đề thi:} \fbox{\textbf{\Large\strut\ 001\ }}
	\end{minipage}
\end{tcolorbox}
\vspace{0.25cm}

\noindent\textbf{\large ĐỀ BÀI (MÃ ĐỀ: 001)}
\vspace{0.1cm}

% --- CÂU 1 (TRANG 1) ---
\begin{ex}
	Thực hiện phép tính:
	\begin{enumerate}[a)]
		\item $A = 3x^3y \cdot 4xy^2$
		\item $B = 3x^2y^3 \cdot 2x^2y - 11x^4y^4$
		\item $C = 18x^6y^5 : 3x^2y^2 - 12x^4y^3$
		\item $D = 4x^3y^2 + 8x^5y^3 : 2x^2y$
	\end{enumerate}
	\loigiai{
		\begin{enumerate}[a)]
			\item Ta có:
			\begin{align*}
				A &= 3x^3y \cdot 4xy^2 \\
				&= (3 \cdot 4) \cdot (x^3 \cdot x) \cdot (y \cdot y^2) \\
				&= 12x^4y^3
			\end{align*}
			\item Ta có:
			\begin{align*}
				B &= 3x^2y^3 \cdot 2x^2y - 11x^4y^4 \\
				&= 6x^4y^4 - 11x^4y^4 \\
				&= (6 - 11)x^4y^4 \\
				&= -5x^4y^4
			\end{align*}
			\item Ta có:
			\begin{align*}
				C &= 18x^6y^5 : 3x^2y^2 - 12x^4y^3 \\
				&= 6x^4y^3 - 12x^4y^3 \\
				&= (6 - 12)x^4y^3 \\
				&= -6x^4y^3
			\end{align*}
			\item Ta có:
			\begin{align*}
				D &= 4x^3y^2 + 8x^5y^3 : 2x^2y \\
				&= 4x^3y^2 + 4x^3y^2 \\
				&= 8x^3y^2
			\end{align*}
		\end{enumerate}
	}
\end{ex}

\newpage
% ----------------------------------------------------------------------
% TRANG 2 (MẶT SAU) -- DÒNG KẺ TIẾP THEO
% ----------------------------------------------------------------------
\noindent
\textbf{Học sinh:} Ánh \quad -- \quad \textbf{Mã đề:} \fbox{\textbf{001}} \hfill \textit{Trang 2/2}\par
\vspace{-0.1cm}
\noindent\hrulefill
\vspace{0.3cm}

% --- CÂU 2 (TRANG 2) ---
\begin{ex}
	Thực hiện phép tính:
	\begin{enumerate}[a)]
		\item $A = -5a^2b^3 \cdot 2ab^2$
		\item $B = 4a^3b \cdot 3ab^2 - 17a^4b^3$
		\item $C = 24a^7b^4 : 6a^4b^2 - 9a^3b^2$
		\item $D = 3a^2b^4 + 14a^4b^5 : 2a^2b$
	\end{enumerate}
	\loigiai{
		\begin{enumerate}[a)]
			\item Ta có:
			\begin{align*}
				A &= -5a^2b^3 \cdot 2ab^2 \\
				&= (-5 \cdot 2) \cdot (a^2 \cdot a) \cdot (b^3 \cdot b^2) \\
				&= -10a^3b^5
			\end{align*}
			\item Ta có:
			\begin{align*}
				B &= 4a^3b \cdot 3ab^2 - 17a^4b^3 \\
				&= 12a^4b^3 - 17a^4b^3 \\
				&= (12 - 17)a^4b^3 \\
				&= -5a^4b^3
			\end{align*}
			\item Ta có:
			\begin{align*}
				C &= 24a^7b^4 : 6a^4b^2 - 9a^3b^2 \\
				&= 4a^3b^2 - 9a^3b^2 \\
				&= (4 - 9)a^3b^2 \\
				&= -5a^3b^2
			\end{align*}
			\item Ta có:
			\begin{align*}
				D &= 3a^2b^4 + 14a^4b^5 : 2a^2b \\
				&= 3a^2b^4 + 7a^2b^4 \\
				&= 10a^2b^4
			\end{align*}
		\end{enumerate}
	}
\end{ex}

\newpage

% ======================================================================
% HỌC SINH: Anh Quân | MÃ: HS034 | MÃ ĐỀ: 009
% ======================================================================
\noindent
\begin{tcolorbox}[colback=white,colframe=black,arc=0mm,boxrule=0.8pt,top=4pt,bottom=4pt,left=6pt,right=6pt]
	\begin{minipage}[c]{0.44\textwidth}
		{\bfseries\large LỚP TOÁN THẦY THÔNG}\\[3pt]
		\textbf{Họ và tên:} \textbf{\large Anh Quân}\\[3pt]
		\textbf{Lớp:} LỚP 10 \quad -- \quad \textbf{Mã HS:} \texttt{HS034}
	\end{minipage}
	\hfill
	\begin{minipage}[c]{0.22\textwidth}
		\begin{tcolorbox}[colback=white,colframe=black,arc=0mm,boxrule=0.6pt,height=1.8cm,halign=center,valign=top,top=2pt]
			{\bfseries\normalsize Điểm}
		\end{tcolorbox}
	\end{minipage}
	\hfill
	\begin{minipage}[c]{0.30\textwidth}
		\raggedleft
		\textit{Thời gian làm bài: 30 phút}\\[4pt]
		\textbf{Mã đề thi:} \fbox{\textbf{\Large\strut\ 009\ }}
	\end{minipage}
\end{tcolorbox}
\vspace{0.25cm}

\noindent\textbf{\large ĐỀ BÀI (MÃ ĐỀ: 009)}
\vspace{0.1cm}

% --- CÂU 1 (TRANG 1) ---
\begin{ex}
	Thực hiện phép tính:
	\begin{enumerate}[a)]
		\item $A = 3x^3y \cdot 4xy^2$
		\item $B = 3x^2y^3 \cdot 2x^2y - 11x^4y^4$
		\item $C = 18x^6y^5 : 3x^2y^2 - 12x^4y^3$
		\item $D = 4x^3y^2 + 8x^5y^3 : 2x^2y$
	\end{enumerate}
	\loigiai{
		\begin{enumerate}[a)]
			\item Ta có:
			\begin{align*}
				A &= 3x^3y \cdot 4xy^2 \\
				&= (3 \cdot 4) \cdot (x^3 \cdot x) \cdot (y \cdot y^2) \\
				&= 12x^4y^3
			\end{align*}
			\item Ta có:
			\begin{align*}
				B &= 3x^2y^3 \cdot 2x^2y - 11x^4y^4 \\
				&= 6x^4y^4 - 11x^4y^4 \\
				&= (6 - 11)x^4y^4 \\
				&= -5x^4y^4
			\end{align*}
			\item Ta có:
			\begin{align*}
				C &= 18x^6y^5 : 3x^2y^2 - 12x^4y^3 \\
				&= 6x^4y^3 - 12x^4y^3 \\
				&= (6 - 12)x^4y^3 \\
				&= -6x^4y^3
			\end{align*}
			\item Ta có:
			\begin{align*}
				D &= 4x^3y^2 + 8x^5y^3 : 2x^2y \\
				&= 4x^3y^2 + 4x^3y^2 \\
				&= 8x^3y^2
			\end{align*}
		\end{enumerate}
	}
\end{ex}

\newpage
% ----------------------------------------------------------------------
% TRANG 2 (MẶT SAU) -- DÒNG KẺ TIẾP THEO
% ----------------------------------------------------------------------
\noindent
\textbf{Học sinh:} Anh Quân \quad -- \quad \textbf{Mã đề:} \fbox{\textbf{009}} \hfill \textit{Trang 2/2}\par
\vspace{-0.1cm}
\noindent\hrulefill
\vspace{0.3cm}

% --- CÂU 2 (TRANG 2) ---
\begin{ex}
	Thực hiện phép tính:
	\begin{enumerate}[a)]
		\item $A = -5a^2b^3 \cdot 2ab^2$
		\item $B = 4a^3b \cdot 3ab^2 - 17a^4b^3$
		\item $C = 24a^7b^4 : 6a^4b^2 - 9a^3b^2$
		\item $D = 3a^2b^4 + 14a^4b^5 : 2a^2b$
	\end{enumerate}
	\loigiai{
		\begin{enumerate}[a)]
			\item Ta có:
			\begin{align*}
				A &= -5a^2b^3 \cdot 2ab^2 \\
				&= (-5 \cdot 2) \cdot (a^2 \cdot a) \cdot (b^3 \cdot b^2) \\
				&= -10a^3b^5
			\end{align*}
			\item Ta có:
			\begin{align*}
				B &= 4a^3b \cdot 3ab^2 - 17a^4b^3 \\
				&= 12a^4b^3 - 17a^4b^3 \\
				&= (12 - 17)a^4b^3 \\
				&= -5a^4b^3
			\end{align*}
			\item Ta có:
			\begin{align*}
				C &= 24a^7b^4 : 6a^4b^2 - 9a^3b^2 \\
				&= 4a^3b^2 - 9a^3b^2 \\
				&= (4 - 9)a^3b^2 \\
				&= -5a^3b^2
			\end{align*}
			\item Ta có:
			\begin{align*}
				D &= 3a^2b^4 + 14a^4b^5 : 2a^2b \\
				&= 3a^2b^4 + 7a^2b^4 \\
				&= 10a^2b^4
			\end{align*}
		\end{enumerate}
	}
\end{ex}

\newpage

% ======================================================================
% HỌC SINH: Bao Nam A | MÃ: HS035 | MÃ ĐỀ: 010
% ======================================================================
\noindent
\begin{tcolorbox}[colback=white,colframe=black,arc=0mm,boxrule=0.8pt,top=4pt,bottom=4pt,left=6pt,right=6pt]
	\begin{minipage}[c]{0.44\textwidth}
		{\bfseries\large LỚP TOÁN THẦY THÔNG}\\[3pt]
		\textbf{Họ và tên:} \textbf{\large Bao Nam A}\\[3pt]
		\textbf{Lớp:} LỚP 10 \quad -- \quad \textbf{Mã HS:} \texttt{HS035}
	\end{minipage}
	\hfill
	\begin{minipage}[c]{0.22\textwidth}
		\begin{tcolorbox}[colback=white,colframe=black,arc=0mm,boxrule=0.6pt,height=1.8cm,halign=center,valign=top,top=2pt]
			{\bfseries\normalsize Điểm}
		\end{tcolorbox}
	\end{minipage}
	\hfill
	\begin{minipage}[c]{0.30\textwidth}
		\raggedleft
		\textit{Thời gian làm bài: 30 phút}\\[4pt]
		\textbf{Mã đề thi:} \fbox{\textbf{\Large\strut\ 010\ }}
	\end{minipage}
\end{tcolorbox}
\vspace{0.25cm}

\noindent\textbf{\large ĐỀ BÀI (MÃ ĐỀ: 010)}
\vspace{0.1cm}

% --- CÂU 1 (TRANG 1) ---
\begin{ex}
	Thực hiện phép tính:
	\begin{enumerate}[a)]
		\item $A = 3x^3y \cdot 4xy^2$
		\item $B = 3x^2y^3 \cdot 2x^2y - 11x^4y^4$
		\item $C = 18x^6y^5 : 3x^2y^2 - 12x^4y^3$
		\item $D = 4x^3y^2 + 8x^5y^3 : 2x^2y$
	\end{enumerate}
	\loigiai{
		\begin{enumerate}[a)]
			\item Ta có:
			\begin{align*}
				A &= 3x^3y \cdot 4xy^2 \\
				&= (3 \cdot 4) \cdot (x^3 \cdot x) \cdot (y \cdot y^2) \\
				&= 12x^4y^3
			\end{align*}
			\item Ta có:
			\begin{align*}
				B &= 3x^2y^3 \cdot 2x^2y - 11x^4y^4 \\
				&= 6x^4y^4 - 11x^4y^4 \\
				&= (6 - 11)x^4y^4 \\
				&= -5x^4y^4
			\end{align*}
			\item Ta có:
			\begin{align*}
				C &= 18x^6y^5 : 3x^2y^2 - 12x^4y^3 \\
				&= 6x^4y^3 - 12x^4y^3 \\
				&= (6 - 12)x^4y^3 \\
				&= -6x^4y^3
			\end{align*}
			\item Ta có:
			\begin{align*}
				D &= 4x^3y^2 + 8x^5y^3 : 2x^2y \\
				&= 4x^3y^2 + 4x^3y^2 \\
				&= 8x^3y^2
			\end{align*}
		\end{enumerate}
	}
\end{ex}

\newpage
% ----------------------------------------------------------------------
% TRANG 2 (MẶT SAU) -- DÒNG KẺ TIẾP THEO
% ----------------------------------------------------------------------
\noindent
\textbf{Học sinh:} Bao Nam A \quad -- \quad \textbf{Mã đề:} \fbox{\textbf{010}} \hfill \textit{Trang 2/2}\par
\vspace{-0.1cm}
\noindent\hrulefill
\vspace{0.3cm}

% --- CÂU 2 (TRANG 2) ---
\begin{ex}
	Thực hiện phép tính:
	\begin{enumerate}[a)]
		\item $A = -5a^2b^3 \cdot 2ab^2$
		\item $B = 4a^3b \cdot 3ab^2 - 17a^4b^3$
		\item $C = 24a^7b^4 : 6a^4b^2 - 9a^3b^2$
		\item $D = 3a^2b^4 + 14a^4b^5 : 2a^2b$
	\end{enumerate}
	\loigiai{
		\begin{enumerate}[a)]
			\item Ta có:
			\begin{align*}
				A &= -5a^2b^3 \cdot 2ab^2 \\
				&= (-5 \cdot 2) \cdot (a^2 \cdot a) \cdot (b^3 \cdot b^2) \\
				&= -10a^3b^5
			\end{align*}
			\item Ta có:
			\begin{align*}
				B &= 4a^3b \cdot 3ab^2 - 17a^4b^3 \\
				&= 12a^4b^3 - 17a^4b^3 \\
				&= (12 - 17)a^4b^3 \\
				&= -5a^4b^3
			\end{align*}
			\item Ta có:
			\begin{align*}
				C &= 24a^7b^4 : 6a^4b^2 - 9a^3b^2 \\
				&= 4a^3b^2 - 9a^3b^2 \\
				&= (4 - 9)a^3b^2 \\
				&= -5a^3b^2
			\end{align*}
			\item Ta có:
			\begin{align*}
				D &= 3a^2b^4 + 14a^4b^5 : 2a^2b \\
				&= 3a^2b^4 + 7a^2b^4 \\
				&= 10a^2b^4
			\end{align*}
		\end{enumerate}
	}
\end{ex}

\newpage

% ======================================================================
% HỌC SINH: Bảo Nam B | MÃ: HS041 | MÃ ĐỀ: 002
% ======================================================================
\noindent
\begin{tcolorbox}[colback=white,colframe=black,arc=0mm,boxrule=0.8pt,top=4pt,bottom=4pt,left=6pt,right=6pt]
	\begin{minipage}[c]{0.44\textwidth}
		{\bfseries\large LỚP TOÁN THẦY THÔNG}\\[3pt]
		\textbf{Họ và tên:} \textbf{\large Bảo Nam B}\\[3pt]
		\textbf{Lớp:} LỚP 10 \quad -- \quad \textbf{Mã HS:} \texttt{HS041}
	\end{minipage}
	\hfill
	\begin{minipage}[c]{0.22\textwidth}
		\begin{tcolorbox}[colback=white,colframe=black,arc=0mm,boxrule=0.6pt,height=1.8cm,halign=center,valign=top,top=2pt]
			{\bfseries\normalsize Điểm}
		\end{tcolorbox}
	\end{minipage}
	\hfill
	\begin{minipage}[c]{0.30\textwidth}
		\raggedleft
		\textit{Thời gian làm bài: 30 phút}\\[4pt]
		\textbf{Mã đề thi:} \fbox{\textbf{\Large\strut\ 002\ }}
	\end{minipage}
\end{tcolorbox}
\vspace{0.25cm}

\noindent\textbf{\large ĐỀ BÀI (MÃ ĐỀ: 002)}
\vspace{0.1cm}

% --- CÂU 1 (TRANG 1) ---
\begin{ex}
	Thực hiện phép tính:
	\begin{enumerate}[a)]
		\item $A = 3x^3y \cdot 4xy^2$
		\item $B = 3x^2y^3 \cdot 2x^2y - 11x^4y^4$
		\item $C = 18x^6y^5 : 3x^2y^2 - 12x^4y^3$
		\item $D = 4x^3y^2 + 8x^5y^3 : 2x^2y$
	\end{enumerate}
	\loigiai{
		\begin{enumerate}[a)]
			\item Ta có:
			\begin{align*}
				A &= 3x^3y \cdot 4xy^2 \\
				&= (3 \cdot 4) \cdot (x^3 \cdot x) \cdot (y \cdot y^2) \\
				&= 12x^4y^3
			\end{align*}
			\item Ta có:
			\begin{align*}
				B &= 3x^2y^3 \cdot 2x^2y - 11x^4y^4 \\
				&= 6x^4y^4 - 11x^4y^4 \\
				&= (6 - 11)x^4y^4 \\
				&= -5x^4y^4
			\end{align*}
			\item Ta có:
			\begin{align*}
				C &= 18x^6y^5 : 3x^2y^2 - 12x^4y^3 \\
				&= 6x^4y^3 - 12x^4y^3 \\
				&= (6 - 12)x^4y^3 \\
				&= -6x^4y^3
			\end{align*}
			\item Ta có:
			\begin{align*}
				D &= 4x^3y^2 + 8x^5y^3 : 2x^2y \\
				&= 4x^3y^2 + 4x^3y^2 \\
				&= 8x^3y^2
			\end{align*}
		\end{enumerate}
	}
\end{ex}

\newpage
% ----------------------------------------------------------------------
% TRANG 2 (MẶT SAU) -- DÒNG KẺ TIẾP THEO
% ----------------------------------------------------------------------
\noindent
\textbf{Học sinh:} Bảo Nam B \quad -- \quad \textbf{Mã đề:} \fbox{\textbf{002}} \hfill \textit{Trang 2/2}\par
\vspace{-0.1cm}
\noindent\hrulefill
\vspace{0.3cm}

% --- CÂU 2 (TRANG 2) ---
\begin{ex}
	Thực hiện phép tính:
	\begin{enumerate}[a)]
		\item $A = -5a^2b^3 \cdot 2ab^2$
		\item $B = 4a^3b \cdot 3ab^2 - 17a^4b^3$
		\item $C = 24a^7b^4 : 6a^4b^2 - 9a^3b^2$
		\item $D = 3a^2b^4 + 14a^4b^5 : 2a^2b$
	\end{enumerate}
	\loigiai{
		\begin{enumerate}[a)]
			\item Ta có:
			\begin{align*}
				A &= -5a^2b^3 \cdot 2ab^2 \\
				&= (-5 \cdot 2) \cdot (a^2 \cdot a) \cdot (b^3 \cdot b^2) \\
				&= -10a^3b^5
			\end{align*}
			\item Ta có:
			\begin{align*}
				B &= 4a^3b \cdot 3ab^2 - 17a^4b^3 \\
				&= 12a^4b^3 - 17a^4b^3 \\
				&= (12 - 17)a^4b^3 \\
				&= -5a^4b^3
			\end{align*}
			\item Ta có:
			\begin{align*}
				C &= 24a^7b^4 : 6a^4b^2 - 9a^3b^2 \\
				&= 4a^3b^2 - 9a^3b^2 \\
				&= (4 - 9)a^3b^2 \\
				&= -5a^3b^2
			\end{align*}
			\item Ta có:
			\begin{align*}
				D &= 3a^2b^4 + 14a^4b^5 : 2a^2b \\
				&= 3a^2b^4 + 7a^2b^4 \\
				&= 10a^2b^4
			\end{align*}
		\end{enumerate}
	}
\end{ex}

\newpage

% ======================================================================
% HỌC SINH: Duy Khánh | MÃ: HS005 | MÃ ĐỀ: 007
% ======================================================================
\noindent
\begin{tcolorbox}[colback=white,colframe=black,arc=0mm,boxrule=0.8pt,top=4pt,bottom=4pt,left=6pt,right=6pt]
	\begin{minipage}[c]{0.44\textwidth}
		{\bfseries\large LỚP TOÁN THẦY THÔNG}\\[3pt]
		\textbf{Họ và tên:} \textbf{\large Duy Khánh}\\[3pt]
		\textbf{Lớp:} LỚP 10 \quad -- \quad \textbf{Mã HS:} \texttt{HS005}
	\end{minipage}
	\hfill
	\begin{minipage}[c]{0.22\textwidth}
		\begin{tcolorbox}[colback=white,colframe=black,arc=0mm,boxrule=0.6pt,height=1.8cm,halign=center,valign=top,top=2pt]
			{\bfseries\normalsize Điểm}
		\end{tcolorbox}
	\end{minipage}
	\hfill
	\begin{minipage}[c]{0.30\textwidth}
		\raggedleft
		\textit{Thời gian làm bài: 30 phút}\\[4pt]
		\textbf{Mã đề thi:} \fbox{\textbf{\Large\strut\ 007\ }}
	\end{minipage}
\end{tcolorbox}
\vspace{0.25cm}

\noindent\textbf{\large ĐỀ BÀI (MÃ ĐỀ: 007)}
\vspace{0.1cm}

% --- CÂU 1 (TRANG 1) ---
\begin{ex}
	Thực hiện phép tính:
	\begin{enumerate}[a)]
		\item $A = 3x^3y \cdot 4xy^2$
		\item $B = 3x^2y^3 \cdot 2x^2y - 11x^4y^4$
		\item $C = 18x^6y^5 : 3x^2y^2 - 12x^4y^3$
		\item $D = 4x^3y^2 + 8x^5y^3 : 2x^2y$
	\end{enumerate}
	\loigiai{
		\begin{enumerate}[a)]
			\item Ta có:
			\begin{align*}
				A &= 3x^3y \cdot 4xy^2 \\
				&= (3 \cdot 4) \cdot (x^3 \cdot x) \cdot (y \cdot y^2) \\
				&= 12x^4y^3
			\end{align*}
			\item Ta có:
			\begin{align*}
				B &= 3x^2y^3 \cdot 2x^2y - 11x^4y^4 \\
				&= 6x^4y^4 - 11x^4y^4 \\
				&= (6 - 11)x^4y^4 \\
				&= -5x^4y^4
			\end{align*}
			\item Ta có:
			\begin{align*}
				C &= 18x^6y^5 : 3x^2y^2 - 12x^4y^3 \\
				&= 6x^4y^3 - 12x^4y^3 \\
				&= (6 - 12)x^4y^3 \\
				&= -6x^4y^3
			\end{align*}
			\item Ta có:
			\begin{align*}
				D &= 4x^3y^2 + 8x^5y^3 : 2x^2y \\
				&= 4x^3y^2 + 4x^3y^2 \\
				&= 8x^3y^2
			\end{align*}
		\end{enumerate}
	}
\end{ex}

\newpage
% ----------------------------------------------------------------------
% TRANG 2 (MẶT SAU) -- DÒNG KẺ TIẾP THEO
% ----------------------------------------------------------------------
\noindent
\textbf{Học sinh:} Duy Khánh \quad -- \quad \textbf{Mã đề:} \fbox{\textbf{007}} \hfill \textit{Trang 2/2}\par
\vspace{-0.1cm}
\noindent\hrulefill
\vspace{0.3cm}

% --- CÂU 2 (TRANG 2) ---
\begin{ex}
	Thực hiện phép tính:
	\begin{enumerate}[a)]
		\item $A = -5a^2b^3 \cdot 2ab^2$
		\item $B = 4a^3b \cdot 3ab^2 - 17a^4b^3$
		\item $C = 24a^7b^4 : 6a^4b^2 - 9a^3b^2$
		\item $D = 3a^2b^4 + 14a^4b^5 : 2a^2b$
	\end{enumerate}
	\loigiai{
		\begin{enumerate}[a)]
			\item Ta có:
			\begin{align*}
				A &= -5a^2b^3 \cdot 2ab^2 \\
				&= (-5 \cdot 2) \cdot (a^2 \cdot a) \cdot (b^3 \cdot b^2) \\
				&= -10a^3b^5
			\end{align*}
			\item Ta có:
			\begin{align*}
				B &= 4a^3b \cdot 3ab^2 - 17a^4b^3 \\
				&= 12a^4b^3 - 17a^4b^3 \\
				&= (12 - 17)a^4b^3 \\
				&= -5a^4b^3
			\end{align*}
			\item Ta có:
			\begin{align*}
				C &= 24a^7b^4 : 6a^4b^2 - 9a^3b^2 \\
				&= 4a^3b^2 - 9a^3b^2 \\
				&= (4 - 9)a^3b^2 \\
				&= -5a^3b^2
			\end{align*}
			\item Ta có:
			\begin{align*}
				D &= 3a^2b^4 + 14a^4b^5 : 2a^2b \\
				&= 3a^2b^4 + 7a^2b^4 \\
				&= 10a^2b^4
			\end{align*}
		\end{enumerate}
	}
\end{ex}

\newpage

% ======================================================================
% HỌC SINH: Gia Bảo A | MÃ: HS003 | MÃ ĐỀ: 008
% ======================================================================
\noindent
\begin{tcolorbox}[colback=white,colframe=black,arc=0mm,boxrule=0.8pt,top=4pt,bottom=4pt,left=6pt,right=6pt]
	\begin{minipage}[c]{0.44\textwidth}
		{\bfseries\large LỚP TOÁN THẦY THÔNG}\\[3pt]
		\textbf{Họ và tên:} \textbf{\large Gia Bảo A}\\[3pt]
		\textbf{Lớp:} LỚP 10 \quad -- \quad \textbf{Mã HS:} \texttt{HS003}
	\end{minipage}
	\hfill
	\begin{minipage}[c]{0.22\textwidth}
		\begin{tcolorbox}[colback=white,colframe=black,arc=0mm,boxrule=0.6pt,height=1.8cm,halign=center,valign=top,top=2pt]
			{\bfseries\normalsize Điểm}
		\end{tcolorbox}
	\end{minipage}
	\hfill
	\begin{minipage}[c]{0.30\textwidth}
		\raggedleft
		\textit{Thời gian làm bài: 30 phút}\\[4pt]
		\textbf{Mã đề thi:} \fbox{\textbf{\Large\strut\ 008\ }}
	\end{minipage}
\end{tcolorbox}
\vspace{0.25cm}

\noindent\textbf{\large ĐỀ BÀI (MÃ ĐỀ: 008)}
\vspace{0.1cm}

% --- CÂU 1 (TRANG 1) ---
\begin{ex}
	Thực hiện phép tính:
	\begin{enumerate}[a)]
		\item $A = 3x^3y \cdot 4xy^2$
		\item $B = 3x^2y^3 \cdot 2x^2y - 11x^4y^4$
		\item $C = 18x^6y^5 : 3x^2y^2 - 12x^4y^3$
		\item $D = 4x^3y^2 + 8x^5y^3 : 2x^2y$
	\end{enumerate}
	\loigiai{
		\begin{enumerate}[a)]
			\item Ta có:
			\begin{align*}
				A &= 3x^3y \cdot 4xy^2 \\
				&= (3 \cdot 4) \cdot (x^3 \cdot x) \cdot (y \cdot y^2) \\
				&= 12x^4y^3
			\end{align*}
			\item Ta có:
			\begin{align*}
				B &= 3x^2y^3 \cdot 2x^2y - 11x^4y^4 \\
				&= 6x^4y^4 - 11x^4y^4 \\
				&= (6 - 11)x^4y^4 \\
				&= -5x^4y^4
			\end{align*}
			\item Ta có:
			\begin{align*}
				C &= 18x^6y^5 : 3x^2y^2 - 12x^4y^3 \\
				&= 6x^4y^3 - 12x^4y^3 \\
				&= (6 - 12)x^4y^3 \\
				&= -6x^4y^3
			\end{align*}
			\item Ta có:
			\begin{align*}
				D &= 4x^3y^2 + 8x^5y^3 : 2x^2y \\
				&= 4x^3y^2 + 4x^3y^2 \\
				&= 8x^3y^2
			\end{align*}
		\end{enumerate}
	}
\end{ex}

\newpage
% ----------------------------------------------------------------------
% TRANG 2 (MẶT SAU) -- DÒNG KẺ TIẾP THEO
% ----------------------------------------------------------------------
\noindent
\textbf{Học sinh:} Gia Bảo A \quad -- \quad \textbf{Mã đề:} \fbox{\textbf{008}} \hfill \textit{Trang 2/2}\par
\vspace{-0.1cm}
\noindent\hrulefill
\vspace{0.3cm}

% --- CÂU 2 (TRANG 2) ---
\begin{ex}
	Thực hiện phép tính:
	\begin{enumerate}[a)]
		\item $A = -5a^2b^3 \cdot 2ab^2$
		\item $B = 4a^3b \cdot 3ab^2 - 17a^4b^3$
		\item $C = 24a^7b^4 : 6a^4b^2 - 9a^3b^2$
		\item $D = 3a^2b^4 + 14a^4b^5 : 2a^2b$
	\end{enumerate}
	\loigiai{
		\begin{enumerate}[a)]
			\item Ta có:
			\begin{align*}
				A &= -5a^2b^3 \cdot 2ab^2 \\
				&= (-5 \cdot 2) \cdot (a^2 \cdot a) \cdot (b^3 \cdot b^2) \\
				&= -10a^3b^5
			\end{align*}
			\item Ta có:
			\begin{align*}
				B &= 4a^3b \cdot 3ab^2 - 17a^4b^3 \\
				&= 12a^4b^3 - 17a^4b^3 \\
				&= (12 - 17)a^4b^3 \\
				&= -5a^4b^3
			\end{align*}
			\item Ta có:
			\begin{align*}
				C &= 24a^7b^4 : 6a^4b^2 - 9a^3b^2 \\
				&= 4a^3b^2 - 9a^3b^2 \\
				&= (4 - 9)a^3b^2 \\
				&= -5a^3b^2
			\end{align*}
			\item Ta có:
			\begin{align*}
				D &= 3a^2b^4 + 14a^4b^5 : 2a^2b \\
				&= 3a^2b^4 + 7a^2b^4 \\
				&= 10a^2b^4
			\end{align*}
		\end{enumerate}
	}
\end{ex}

\newpage

% ======================================================================
% HỌC SINH: Gia Bảo B | MÃ: HS006 | MÃ ĐỀ: 005
% ======================================================================
\noindent
\begin{tcolorbox}[colback=white,colframe=black,arc=0mm,boxrule=0.8pt,top=4pt,bottom=4pt,left=6pt,right=6pt]
	\begin{minipage}[c]{0.44\textwidth}
		{\bfseries\large LỚP TOÁN THẦY THÔNG}\\[3pt]
		\textbf{Họ và tên:} \textbf{\large Gia Bảo B}\\[3pt]
		\textbf{Lớp:} LỚP 10 \quad -- \quad \textbf{Mã HS:} \texttt{HS006}
	\end{minipage}
	\hfill
	\begin{minipage}[c]{0.22\textwidth}
		\begin{tcolorbox}[colback=white,colframe=black,arc=0mm,boxrule=0.6pt,height=1.8cm,halign=center,valign=top,top=2pt]
			{\bfseries\normalsize Điểm}
		\end{tcolorbox}
	\end{minipage}
	\hfill
	\begin{minipage}[c]{0.30\textwidth}
		\raggedleft
		\textit{Thời gian làm bài: 30 phút}\\[4pt]
		\textbf{Mã đề thi:} \fbox{\textbf{\Large\strut\ 005\ }}
	\end{minipage}
\end{tcolorbox}
\vspace{0.25cm}

\noindent\textbf{\large ĐỀ BÀI (MÃ ĐỀ: 005)}
\vspace{0.1cm}

% --- CÂU 1 (TRANG 1) ---
\begin{ex}
	Thực hiện phép tính:
	\begin{enumerate}[a)]
		\item $A = 3x^3y \cdot 4xy^2$
		\item $B = 3x^2y^3 \cdot 2x^2y - 11x^4y^4$
		\item $C = 18x^6y^5 : 3x^2y^2 - 12x^4y^3$
		\item $D = 4x^3y^2 + 8x^5y^3 : 2x^2y$
	\end{enumerate}
	\loigiai{
		\begin{enumerate}[a)]
			\item Ta có:
			\begin{align*}
				A &= 3x^3y \cdot 4xy^2 \\
				&= (3 \cdot 4) \cdot (x^3 \cdot x) \cdot (y \cdot y^2) \\
				&= 12x^4y^3
			\end{align*}
			\item Ta có:
			\begin{align*}
				B &= 3x^2y^3 \cdot 2x^2y - 11x^4y^4 \\
				&= 6x^4y^4 - 11x^4y^4 \\
				&= (6 - 11)x^4y^4 \\
				&= -5x^4y^4
			\end{align*}
			\item Ta có:
			\begin{align*}
				C &= 18x^6y^5 : 3x^2y^2 - 12x^4y^3 \\
				&= 6x^4y^3 - 12x^4y^3 \\
				&= (6 - 12)x^4y^3 \\
				&= -6x^4y^3
			\end{align*}
			\item Ta có:
			\begin{align*}
				D &= 4x^3y^2 + 8x^5y^3 : 2x^2y \\
				&= 4x^3y^2 + 4x^3y^2 \\
				&= 8x^3y^2
			\end{align*}
		\end{enumerate}
	}
\end{ex}

\newpage
% ----------------------------------------------------------------------
% TRANG 2 (MẶT SAU) -- DÒNG KẺ TIẾP THEO
% ----------------------------------------------------------------------
\noindent
\textbf{Học sinh:} Gia Bảo B \quad -- \quad \textbf{Mã đề:} \fbox{\textbf{005}} \hfill \textit{Trang 2/2}\par
\vspace{-0.1cm}
\noindent\hrulefill
\vspace{0.3cm}

% --- CÂU 2 (TRANG 2) ---
\begin{ex}
	Thực hiện phép tính:
	\begin{enumerate}[a)]
		\item $A = -5a^2b^3 \cdot 2ab^2$
		\item $B = 4a^3b \cdot 3ab^2 - 17a^4b^3$
		\item $C = 24a^7b^4 : 6a^4b^2 - 9a^3b^2$
		\item $D = 3a^2b^4 + 14a^4b^5 : 2a^2b$
	\end{enumerate}
	\loigiai{
		\begin{enumerate}[a)]
			\item Ta có:
			\begin{align*}
				A &= -5a^2b^3 \cdot 2ab^2 \\
				&= (-5 \cdot 2) \cdot (a^2 \cdot a) \cdot (b^3 \cdot b^2) \\
				&= -10a^3b^5
			\end{align*}
			\item Ta có:
			\begin{align*}
				B &= 4a^3b \cdot 3ab^2 - 17a^4b^3 \\
				&= 12a^4b^3 - 17a^4b^3 \\
				&= (12 - 17)a^4b^3 \\
				&= -5a^4b^3
			\end{align*}
			\item Ta có:
			\begin{align*}
				C &= 24a^7b^4 : 6a^4b^2 - 9a^3b^2 \\
				&= 4a^3b^2 - 9a^3b^2 \\
				&= (4 - 9)a^3b^2 \\
				&= -5a^3b^2
			\end{align*}
			\item Ta có:
			\begin{align*}
				D &= 3a^2b^4 + 14a^4b^5 : 2a^2b \\
				&= 3a^2b^4 + 7a^2b^4 \\
				&= 10a^2b^4
			\end{align*}
		\end{enumerate}
	}
\end{ex}

\newpage

% ======================================================================
% HỌC SINH: Gia Hân(on) | MÃ: HS043 | MÃ ĐỀ: 006
% ======================================================================
\noindent
\begin{tcolorbox}[colback=white,colframe=black,arc=0mm,boxrule=0.8pt,top=4pt,bottom=4pt,left=6pt,right=6pt]
	\begin{minipage}[c]{0.44\textwidth}
		{\bfseries\large LỚP TOÁN THẦY THÔNG}\\[3pt]
		\textbf{Họ và tên:} \textbf{\large Gia Hân(on)}\\[3pt]
		\textbf{Lớp:} LỚP 10 \quad -- \quad \textbf{Mã HS:} \texttt{HS043}
	\end{minipage}
	\hfill
	\begin{minipage}[c]{0.22\textwidth}
		\begin{tcolorbox}[colback=white,colframe=black,arc=0mm,boxrule=0.6pt,height=1.8cm,halign=center,valign=top,top=2pt]
			{\bfseries\normalsize Điểm}
		\end{tcolorbox}
	\end{minipage}
	\hfill
	\begin{minipage}[c]{0.30\textwidth}
		\raggedleft
		\textit{Thời gian làm bài: 30 phút}\\[4pt]
		\textbf{Mã đề thi:} \fbox{\textbf{\Large\strut\ 006\ }}
	\end{minipage}
\end{tcolorbox}
\vspace{0.25cm}

\noindent\textbf{\large ĐỀ BÀI (MÃ ĐỀ: 006)}
\vspace{0.1cm}

% --- CÂU 1 (TRANG 1) ---
\begin{ex}
	Thực hiện phép tính:
	\begin{enumerate}[a)]
		\item $A = 3x^3y \cdot 4xy^2$
		\item $B = 3x^2y^3 \cdot 2x^2y - 11x^4y^4$
		\item $C = 18x^6y^5 : 3x^2y^2 - 12x^4y^3$
		\item $D = 4x^3y^2 + 8x^5y^3 : 2x^2y$
	\end{enumerate}
	\loigiai{
		\begin{enumerate}[a)]
			\item Ta có:
			\begin{align*}
				A &= 3x^3y \cdot 4xy^2 \\
				&= (3 \cdot 4) \cdot (x^3 \cdot x) \cdot (y \cdot y^2) \\
				&= 12x^4y^3
			\end{align*}
			\item Ta có:
			\begin{align*}
				B &= 3x^2y^3 \cdot 2x^2y - 11x^4y^4 \\
				&= 6x^4y^4 - 11x^4y^4 \\
				&= (6 - 11)x^4y^4 \\
				&= -5x^4y^4
			\end{align*}
			\item Ta có:
			\begin{align*}
				C &= 18x^6y^5 : 3x^2y^2 - 12x^4y^3 \\
				&= 6x^4y^3 - 12x^4y^3 \\
				&= (6 - 12)x^4y^3 \\
				&= -6x^4y^3
			\end{align*}
			\item Ta có:
			\begin{align*}
				D &= 4x^3y^2 + 8x^5y^3 : 2x^2y \\
				&= 4x^3y^2 + 4x^3y^2 \\
				&= 8x^3y^2
			\end{align*}
		\end{enumerate}
	}
\end{ex}

\newpage
% ----------------------------------------------------------------------
% TRANG 2 (MẶT SAU) -- DÒNG KẺ TIẾP THEO
% ----------------------------------------------------------------------
\noindent
\textbf{Học sinh:} Gia Hân(on) \quad -- \quad \textbf{Mã đề:} \fbox{\textbf{006}} \hfill \textit{Trang 2/2}\par
\vspace{-0.1cm}
\noindent\hrulefill
\vspace{0.3cm}

% --- CÂU 2 (TRANG 2) ---
\begin{ex}
	Thực hiện phép tính:
	\begin{enumerate}[a)]
		\item $A = -5a^2b^3 \cdot 2ab^2$
		\item $B = 4a^3b \cdot 3ab^2 - 17a^4b^3$
		\item $C = 24a^7b^4 : 6a^4b^2 - 9a^3b^2$
		\item $D = 3a^2b^4 + 14a^4b^5 : 2a^2b$
	\end{enumerate}
	\loigiai{
		\begin{enumerate}[a)]
			\item Ta có:
			\begin{align*}
				A &= -5a^2b^3 \cdot 2ab^2 \\
				&= (-5 \cdot 2) \cdot (a^2 \cdot a) \cdot (b^3 \cdot b^2) \\
				&= -10a^3b^5
			\end{align*}
			\item Ta có:
			\begin{align*}
				B &= 4a^3b \cdot 3ab^2 - 17a^4b^3 \\
				&= 12a^4b^3 - 17a^4b^3 \\
				&= (12 - 17)a^4b^3 \\
				&= -5a^4b^3
			\end{align*}
			\item Ta có:
			\begin{align*}
				C &= 24a^7b^4 : 6a^4b^2 - 9a^3b^2 \\
				&= 4a^3b^2 - 9a^3b^2 \\
				&= (4 - 9)a^3b^2 \\
				&= -5a^3b^2
			\end{align*}
			\item Ta có:
			\begin{align*}
				D &= 3a^2b^4 + 14a^4b^5 : 2a^2b \\
				&= 3a^2b^4 + 7a^2b^4 \\
				&= 10a^2b^4
			\end{align*}
		\end{enumerate}
	}
\end{ex}

\newpage

% ======================================================================
% HỌC SINH: Gia Hưng A | MÃ: HS014 | MÃ ĐỀ: 003
% ======================================================================
\noindent
\begin{tcolorbox}[colback=white,colframe=black,arc=0mm,boxrule=0.8pt,top=4pt,bottom=4pt,left=6pt,right=6pt]
	\begin{minipage}[c]{0.44\textwidth}
		{\bfseries\large LỚP TOÁN THẦY THÔNG}\\[3pt]
		\textbf{Họ và tên:} \textbf{\large Gia Hưng A}\\[3pt]
		\textbf{Lớp:} LỚP 10 \quad -- \quad \textbf{Mã HS:} \texttt{HS014}
	\end{minipage}
	\hfill
	\begin{minipage}[c]{0.22\textwidth}
		\begin{tcolorbox}[colback=white,colframe=black,arc=0mm,boxrule=0.6pt,height=1.8cm,halign=center,valign=top,top=2pt]
			{\bfseries\normalsize Điểm}
		\end{tcolorbox}
	\end{minipage}
	\hfill
	\begin{minipage}[c]{0.30\textwidth}
		\raggedleft
		\textit{Thời gian làm bài: 30 phút}\\[4pt]
		\textbf{Mã đề thi:} \fbox{\textbf{\Large\strut\ 003\ }}
	\end{minipage}
\end{tcolorbox}
\vspace{0.25cm}

\noindent\textbf{\large ĐỀ BÀI (MÃ ĐỀ: 003)}
\vspace{0.1cm}

% --- CÂU 1 (TRANG 1) ---
\begin{ex}
	Thực hiện phép tính:
	\begin{enumerate}[a)]
		\item $A = 3x^3y \cdot 4xy^2$
		\item $B = 3x^2y^3 \cdot 2x^2y - 11x^4y^4$
		\item $C = 18x^6y^5 : 3x^2y^2 - 12x^4y^3$
		\item $D = 4x^3y^2 + 8x^5y^3 : 2x^2y$
	\end{enumerate}
	\loigiai{
		\begin{enumerate}[a)]
			\item Ta có:
			\begin{align*}
				A &= 3x^3y \cdot 4xy^2 \\
				&= (3 \cdot 4) \cdot (x^3 \cdot x) \cdot (y \cdot y^2) \\
				&= 12x^4y^3
			\end{align*}
			\item Ta có:
			\begin{align*}
				B &= 3x^2y^3 \cdot 2x^2y - 11x^4y^4 \\
				&= 6x^4y^4 - 11x^4y^4 \\
				&= (6 - 11)x^4y^4 \\
				&= -5x^4y^4
			\end{align*}
			\item Ta có:
			\begin{align*}
				C &= 18x^6y^5 : 3x^2y^2 - 12x^4y^3 \\
				&= 6x^4y^3 - 12x^4y^3 \\
				&= (6 - 12)x^4y^3 \\
				&= -6x^4y^3
			\end{align*}
			\item Ta có:
			\begin{align*}
				D &= 4x^3y^2 + 8x^5y^3 : 2x^2y \\
				&= 4x^3y^2 + 4x^3y^2 \\
				&= 8x^3y^2
			\end{align*}
		\end{enumerate}
	}
\end{ex}

\newpage
% ----------------------------------------------------------------------
% TRANG 2 (MẶT SAU) -- DÒNG KẺ TIẾP THEO
% ----------------------------------------------------------------------
\noindent
\textbf{Học sinh:} Gia Hưng A \quad -- \quad \textbf{Mã đề:} \fbox{\textbf{003}} \hfill \textit{Trang 2/2}\par
\vspace{-0.1cm}
\noindent\hrulefill
\vspace{0.3cm}

% --- CÂU 2 (TRANG 2) ---
\begin{ex}
	Thực hiện phép tính:
	\begin{enumerate}[a)]
		\item $A = -5a^2b^3 \cdot 2ab^2$
		\item $B = 4a^3b \cdot 3ab^2 - 17a^4b^3$
		\item $C = 24a^7b^4 : 6a^4b^2 - 9a^3b^2$
		\item $D = 3a^2b^4 + 14a^4b^5 : 2a^2b$
	\end{enumerate}
	\loigiai{
		\begin{enumerate}[a)]
			\item Ta có:
			\begin{align*}
				A &= -5a^2b^3 \cdot 2ab^2 \\
				&= (-5 \cdot 2) \cdot (a^2 \cdot a) \cdot (b^3 \cdot b^2) \\
				&= -10a^3b^5
			\end{align*}
			\item Ta có:
			\begin{align*}
				B &= 4a^3b \cdot 3ab^2 - 17a^4b^3 \\
				&= 12a^4b^3 - 17a^4b^3 \\
				&= (12 - 17)a^4b^3 \\
				&= -5a^4b^3
			\end{align*}
			\item Ta có:
			\begin{align*}
				C &= 24a^7b^4 : 6a^4b^2 - 9a^3b^2 \\
				&= 4a^3b^2 - 9a^3b^2 \\
				&= (4 - 9)a^3b^2 \\
				&= -5a^3b^2
			\end{align*}
			\item Ta có:
			\begin{align*}
				D &= 3a^2b^4 + 14a^4b^5 : 2a^2b \\
				&= 3a^2b^4 + 7a^2b^4 \\
				&= 10a^2b^4
			\end{align*}
		\end{enumerate}
	}
\end{ex}

\newpage

% ======================================================================
% HỌC SINH: Gia Hưng B | MÃ: HS036 | MÃ ĐỀ: 008
% ======================================================================
\noindent
\begin{tcolorbox}[colback=white,colframe=black,arc=0mm,boxrule=0.8pt,top=4pt,bottom=4pt,left=6pt,right=6pt]
	\begin{minipage}[c]{0.44\textwidth}
		{\bfseries\large LỚP TOÁN THẦY THÔNG}\\[3pt]
		\textbf{Họ và tên:} \textbf{\large Gia Hưng B}\\[3pt]
		\textbf{Lớp:} LỚP 10 \quad -- \quad \textbf{Mã HS:} \texttt{HS036}
	\end{minipage}
	\hfill
	\begin{minipage}[c]{0.22\textwidth}
		\begin{tcolorbox}[colback=white,colframe=black,arc=0mm,boxrule=0.6pt,height=1.8cm,halign=center,valign=top,top=2pt]
			{\bfseries\normalsize Điểm}
		\end{tcolorbox}
	\end{minipage}
	\hfill
	\begin{minipage}[c]{0.30\textwidth}
		\raggedleft
		\textit{Thời gian làm bài: 30 phút}\\[4pt]
		\textbf{Mã đề thi:} \fbox{\textbf{\Large\strut\ 008\ }}
	\end{minipage}
\end{tcolorbox}
\vspace{0.25cm}

\noindent\textbf{\large ĐỀ BÀI (MÃ ĐỀ: 008)}
\vspace{0.1cm}

% --- CÂU 1 (TRANG 1) ---
\begin{ex}
	Thực hiện phép tính:
	\begin{enumerate}[a)]
		\item $A = 3x^3y \cdot 4xy^2$
		\item $B = 3x^2y^3 \cdot 2x^2y - 11x^4y^4$
		\item $C = 18x^6y^5 : 3x^2y^2 - 12x^4y^3$
		\item $D = 4x^3y^2 + 8x^5y^3 : 2x^2y$
	\end{enumerate}
	\loigiai{
		\begin{enumerate}[a)]
			\item Ta có:
			\begin{align*}
				A &= 3x^3y \cdot 4xy^2 \\
				&= (3 \cdot 4) \cdot (x^3 \cdot x) \cdot (y \cdot y^2) \\
				&= 12x^4y^3
			\end{align*}
			\item Ta có:
			\begin{align*}
				B &= 3x^2y^3 \cdot 2x^2y - 11x^4y^4 \\
				&= 6x^4y^4 - 11x^4y^4 \\
				&= (6 - 11)x^4y^4 \\
				&= -5x^4y^4
			\end{align*}
			\item Ta có:
			\begin{align*}
				C &= 18x^6y^5 : 3x^2y^2 - 12x^4y^3 \\
				&= 6x^4y^3 - 12x^4y^3 \\
				&= (6 - 12)x^4y^3 \\
				&= -6x^4y^3
			\end{align*}
			\item Ta có:
			\begin{align*}
				D &= 4x^3y^2 + 8x^5y^3 : 2x^2y \\
				&= 4x^3y^2 + 4x^3y^2 \\
				&= 8x^3y^2
			\end{align*}
		\end{enumerate}
	}
\end{ex}

\newpage
% ----------------------------------------------------------------------
% TRANG 2 (MẶT SAU) -- DÒNG KẺ TIẾP THEO
% ----------------------------------------------------------------------
\noindent
\textbf{Học sinh:} Gia Hưng B \quad -- \quad \textbf{Mã đề:} \fbox{\textbf{008}} \hfill \textit{Trang 2/2}\par
\vspace{-0.1cm}
\noindent\hrulefill
\vspace{0.3cm}

% --- CÂU 2 (TRANG 2) ---
\begin{ex}
	Thực hiện phép tính:
	\begin{enumerate}[a)]
		\item $A = -5a^2b^3 \cdot 2ab^2$
		\item $B = 4a^3b \cdot 3ab^2 - 17a^4b^3$
		\item $C = 24a^7b^4 : 6a^4b^2 - 9a^3b^2$
		\item $D = 3a^2b^4 + 14a^4b^5 : 2a^2b$
	\end{enumerate}
	\loigiai{
		\begin{enumerate}[a)]
			\item Ta có:
			\begin{align*}
				A &= -5a^2b^3 \cdot 2ab^2 \\
				&= (-5 \cdot 2) \cdot (a^2 \cdot a) \cdot (b^3 \cdot b^2) \\
				&= -10a^3b^5
			\end{align*}
			\item Ta có:
			\begin{align*}
				B &= 4a^3b \cdot 3ab^2 - 17a^4b^3 \\
				&= 12a^4b^3 - 17a^4b^3 \\
				&= (12 - 17)a^4b^3 \\
				&= -5a^4b^3
			\end{align*}
			\item Ta có:
			\begin{align*}
				C &= 24a^7b^4 : 6a^4b^2 - 9a^3b^2 \\
				&= 4a^3b^2 - 9a^3b^2 \\
				&= (4 - 9)a^3b^2 \\
				&= -5a^3b^2
			\end{align*}
			\item Ta có:
			\begin{align*}
				D &= 3a^2b^4 + 14a^4b^5 : 2a^2b \\
				&= 3a^2b^4 + 7a^2b^4 \\
				&= 10a^2b^4
			\end{align*}
		\end{enumerate}
	}
\end{ex}

\newpage

% ======================================================================
% HỌC SINH: Hà Vũ (on) | MÃ: HS042 | MÃ ĐỀ: 010
% ======================================================================
\noindent
\begin{tcolorbox}[colback=white,colframe=black,arc=0mm,boxrule=0.8pt,top=4pt,bottom=4pt,left=6pt,right=6pt]
	\begin{minipage}[c]{0.44\textwidth}
		{\bfseries\large LỚP TOÁN THẦY THÔNG}\\[3pt]
		\textbf{Họ và tên:} \textbf{\large Hà Vũ (on)}\\[3pt]
		\textbf{Lớp:} LỚP 10 \quad -- \quad \textbf{Mã HS:} \texttt{HS042}
	\end{minipage}
	\hfill
	\begin{minipage}[c]{0.22\textwidth}
		\begin{tcolorbox}[colback=white,colframe=black,arc=0mm,boxrule=0.6pt,height=1.8cm,halign=center,valign=top,top=2pt]
			{\bfseries\normalsize Điểm}
		\end{tcolorbox}
	\end{minipage}
	\hfill
	\begin{minipage}[c]{0.30\textwidth}
		\raggedleft
		\textit{Thời gian làm bài: 30 phút}\\[4pt]
		\textbf{Mã đề thi:} \fbox{\textbf{\Large\strut\ 010\ }}
	\end{minipage}
\end{tcolorbox}
\vspace{0.25cm}

\noindent\textbf{\large ĐỀ BÀI (MÃ ĐỀ: 010)}
\vspace{0.1cm}

% --- CÂU 1 (TRANG 1) ---
\begin{ex}
	Thực hiện phép tính:
	\begin{enumerate}[a)]
		\item $A = 3x^3y \cdot 4xy^2$
		\item $B = 3x^2y^3 \cdot 2x^2y - 11x^4y^4$
		\item $C = 18x^6y^5 : 3x^2y^2 - 12x^4y^3$
		\item $D = 4x^3y^2 + 8x^5y^3 : 2x^2y$
	\end{enumerate}
	\loigiai{
		\begin{enumerate}[a)]
			\item Ta có:
			\begin{align*}
				A &= 3x^3y \cdot 4xy^2 \\
				&= (3 \cdot 4) \cdot (x^3 \cdot x) \cdot (y \cdot y^2) \\
				&= 12x^4y^3
			\end{align*}
			\item Ta có:
			\begin{align*}
				B &= 3x^2y^3 \cdot 2x^2y - 11x^4y^4 \\
				&= 6x^4y^4 - 11x^4y^4 \\
				&= (6 - 11)x^4y^4 \\
				&= -5x^4y^4
			\end{align*}
			\item Ta có:
			\begin{align*}
				C &= 18x^6y^5 : 3x^2y^2 - 12x^4y^3 \\
				&= 6x^4y^3 - 12x^4y^3 \\
				&= (6 - 12)x^4y^3 \\
				&= -6x^4y^3
			\end{align*}
			\item Ta có:
			\begin{align*}
				D &= 4x^3y^2 + 8x^5y^3 : 2x^2y \\
				&= 4x^3y^2 + 4x^3y^2 \\
				&= 8x^3y^2
			\end{align*}
		\end{enumerate}
	}
\end{ex}

\newpage
% ----------------------------------------------------------------------
% TRANG 2 (MẶT SAU) -- DÒNG KẺ TIẾP THEO
% ----------------------------------------------------------------------
\noindent
\textbf{Học sinh:} Hà Vũ (on) \quad -- \quad \textbf{Mã đề:} \fbox{\textbf{010}} \hfill \textit{Trang 2/2}\par
\vspace{-0.1cm}
\noindent\hrulefill
\vspace{0.3cm}

% --- CÂU 2 (TRANG 2) ---
\begin{ex}
	Thực hiện phép tính:
	\begin{enumerate}[a)]
		\item $A = -5a^2b^3 \cdot 2ab^2$
		\item $B = 4a^3b \cdot 3ab^2 - 17a^4b^3$
		\item $C = 24a^7b^4 : 6a^4b^2 - 9a^3b^2$
		\item $D = 3a^2b^4 + 14a^4b^5 : 2a^2b$
	\end{enumerate}
	\loigiai{
		\begin{enumerate}[a)]
			\item Ta có:
			\begin{align*}
				A &= -5a^2b^3 \cdot 2ab^2 \\
				&= (-5 \cdot 2) \cdot (a^2 \cdot a) \cdot (b^3 \cdot b^2) \\
				&= -10a^3b^5
			\end{align*}
			\item Ta có:
			\begin{align*}
				B &= 4a^3b \cdot 3ab^2 - 17a^4b^3 \\
				&= 12a^4b^3 - 17a^4b^3 \\
				&= (12 - 17)a^4b^3 \\
				&= -5a^4b^3
			\end{align*}
			\item Ta có:
			\begin{align*}
				C &= 24a^7b^4 : 6a^4b^2 - 9a^3b^2 \\
				&= 4a^3b^2 - 9a^3b^2 \\
				&= (4 - 9)a^3b^2 \\
				&= -5a^3b^2
			\end{align*}
			\item Ta có:
			\begin{align*}
				D &= 3a^2b^4 + 14a^4b^5 : 2a^2b \\
				&= 3a^2b^4 + 7a^2b^4 \\
				&= 10a^2b^4
			\end{align*}
		\end{enumerate}
	}
\end{ex}

\newpage

% ======================================================================
% HỌC SINH: Hà Yên | MÃ: HS033 | MÃ ĐỀ: 001
% ======================================================================
\noindent
\begin{tcolorbox}[colback=white,colframe=black,arc=0mm,boxrule=0.8pt,top=4pt,bottom=4pt,left=6pt,right=6pt]
	\begin{minipage}[c]{0.44\textwidth}
		{\bfseries\large LỚP TOÁN THẦY THÔNG}\\[3pt]
		\textbf{Họ và tên:} \textbf{\large Hà Yên}\\[3pt]
		\textbf{Lớp:} LỚP 10 \quad -- \quad \textbf{Mã HS:} \texttt{HS033}
	\end{minipage}
	\hfill
	\begin{minipage}[c]{0.22\textwidth}
		\begin{tcolorbox}[colback=white,colframe=black,arc=0mm,boxrule=0.6pt,height=1.8cm,halign=center,valign=top,top=2pt]
			{\bfseries\normalsize Điểm}
		\end{tcolorbox}
	\end{minipage}
	\hfill
	\begin{minipage}[c]{0.30\textwidth}
		\raggedleft
		\textit{Thời gian làm bài: 30 phút}\\[4pt]
		\textbf{Mã đề thi:} \fbox{\textbf{\Large\strut\ 001\ }}
	\end{minipage}
\end{tcolorbox}
\vspace{0.25cm}

\noindent\textbf{\large ĐỀ BÀI (MÃ ĐỀ: 001)}
\vspace{0.1cm}

% --- CÂU 1 (TRANG 1) ---
\begin{ex}
	Thực hiện phép tính:
	\begin{enumerate}[a)]
		\item $A = 3x^3y \cdot 4xy^2$
		\item $B = 3x^2y^3 \cdot 2x^2y - 11x^4y^4$
		\item $C = 18x^6y^5 : 3x^2y^2 - 12x^4y^3$
		\item $D = 4x^3y^2 + 8x^5y^3 : 2x^2y$
	\end{enumerate}
	\loigiai{
		\begin{enumerate}[a)]
			\item Ta có:
			\begin{align*}
				A &= 3x^3y \cdot 4xy^2 \\
				&= (3 \cdot 4) \cdot (x^3 \cdot x) \cdot (y \cdot y^2) \\
				&= 12x^4y^3
			\end{align*}
			\item Ta có:
			\begin{align*}
				B &= 3x^2y^3 \cdot 2x^2y - 11x^4y^4 \\
				&= 6x^4y^4 - 11x^4y^4 \\
				&= (6 - 11)x^4y^4 \\
				&= -5x^4y^4
			\end{align*}
			\item Ta có:
			\begin{align*}
				C &= 18x^6y^5 : 3x^2y^2 - 12x^4y^3 \\
				&= 6x^4y^3 - 12x^4y^3 \\
				&= (6 - 12)x^4y^3 \\
				&= -6x^4y^3
			\end{align*}
			\item Ta có:
			\begin{align*}
				D &= 4x^3y^2 + 8x^5y^3 : 2x^2y \\
				&= 4x^3y^2 + 4x^3y^2 \\
				&= 8x^3y^2
			\end{align*}
		\end{enumerate}
	}
\end{ex}

\newpage
% ----------------------------------------------------------------------
% TRANG 2 (MẶT SAU) -- DÒNG KẺ TIẾP THEO
% ----------------------------------------------------------------------
\noindent
\textbf{Học sinh:} Hà Yên \quad -- \quad \textbf{Mã đề:} \fbox{\textbf{001}} \hfill \textit{Trang 2/2}\par
\vspace{-0.1cm}
\noindent\hrulefill
\vspace{0.3cm}

% --- CÂU 2 (TRANG 2) ---
\begin{ex}
	Thực hiện phép tính:
	\begin{enumerate}[a)]
		\item $A = -5a^2b^3 \cdot 2ab^2$
		\item $B = 4a^3b \cdot 3ab^2 - 17a^4b^3$
		\item $C = 24a^7b^4 : 6a^4b^2 - 9a^3b^2$
		\item $D = 3a^2b^4 + 14a^4b^5 : 2a^2b$
	\end{enumerate}
	\loigiai{
		\begin{enumerate}[a)]
			\item Ta có:
			\begin{align*}
				A &= -5a^2b^3 \cdot 2ab^2 \\
				&= (-5 \cdot 2) \cdot (a^2 \cdot a) \cdot (b^3 \cdot b^2) \\
				&= -10a^3b^5
			\end{align*}
			\item Ta có:
			\begin{align*}
				B &= 4a^3b \cdot 3ab^2 - 17a^4b^3 \\
				&= 12a^4b^3 - 17a^4b^3 \\
				&= (12 - 17)a^4b^3 \\
				&= -5a^4b^3
			\end{align*}
			\item Ta có:
			\begin{align*}
				C &= 24a^7b^4 : 6a^4b^2 - 9a^3b^2 \\
				&= 4a^3b^2 - 9a^3b^2 \\
				&= (4 - 9)a^3b^2 \\
				&= -5a^3b^2
			\end{align*}
			\item Ta có:
			\begin{align*}
				D &= 3a^2b^4 + 14a^4b^5 : 2a^2b \\
				&= 3a^2b^4 + 7a^2b^4 \\
				&= 10a^2b^4
			\end{align*}
		\end{enumerate}
	}
\end{ex}

\newpage

% ======================================================================
% HỌC SINH: Hoàng Nguyên | MÃ: HS008 | MÃ ĐỀ: 009
% ======================================================================
\noindent
\begin{tcolorbox}[colback=white,colframe=black,arc=0mm,boxrule=0.8pt,top=4pt,bottom=4pt,left=6pt,right=6pt]
	\begin{minipage}[c]{0.44\textwidth}
		{\bfseries\large LỚP TOÁN THẦY THÔNG}\\[3pt]
		\textbf{Họ và tên:} \textbf{\large Hoàng Nguyên}\\[3pt]
		\textbf{Lớp:} LỚP 10 \quad -- \quad \textbf{Mã HS:} \texttt{HS008}
	\end{minipage}
	\hfill
	\begin{minipage}[c]{0.22\textwidth}
		\begin{tcolorbox}[colback=white,colframe=black,arc=0mm,boxrule=0.6pt,height=1.8cm,halign=center,valign=top,top=2pt]
			{\bfseries\normalsize Điểm}
		\end{tcolorbox}
	\end{minipage}
	\hfill
	\begin{minipage}[c]{0.30\textwidth}
		\raggedleft
		\textit{Thời gian làm bài: 30 phút}\\[4pt]
		\textbf{Mã đề thi:} \fbox{\textbf{\Large\strut\ 009\ }}
	\end{minipage}
\end{tcolorbox}
\vspace{0.25cm}

\noindent\textbf{\large ĐỀ BÀI (MÃ ĐỀ: 009)}
\vspace{0.1cm}

% --- CÂU 1 (TRANG 1) ---
\begin{ex}
	Thực hiện phép tính:
	\begin{enumerate}[a)]
		\item $A = 3x^3y \cdot 4xy^2$
		\item $B = 3x^2y^3 \cdot 2x^2y - 11x^4y^4$
		\item $C = 18x^6y^5 : 3x^2y^2 - 12x^4y^3$
		\item $D = 4x^3y^2 + 8x^5y^3 : 2x^2y$
	\end{enumerate}
	\loigiai{
		\begin{enumerate}[a)]
			\item Ta có:
			\begin{align*}
				A &= 3x^3y \cdot 4xy^2 \\
				&= (3 \cdot 4) \cdot (x^3 \cdot x) \cdot (y \cdot y^2) \\
				&= 12x^4y^3
			\end{align*}
			\item Ta có:
			\begin{align*}
				B &= 3x^2y^3 \cdot 2x^2y - 11x^4y^4 \\
				&= 6x^4y^4 - 11x^4y^4 \\
				&= (6 - 11)x^4y^4 \\
				&= -5x^4y^4
			\end{align*}
			\item Ta có:
			\begin{align*}
				C &= 18x^6y^5 : 3x^2y^2 - 12x^4y^3 \\
				&= 6x^4y^3 - 12x^4y^3 \\
				&= (6 - 12)x^4y^3 \\
				&= -6x^4y^3
			\end{align*}
			\item Ta có:
			\begin{align*}
				D &= 4x^3y^2 + 8x^5y^3 : 2x^2y \\
				&= 4x^3y^2 + 4x^3y^2 \\
				&= 8x^3y^2
			\end{align*}
		\end{enumerate}
	}
\end{ex}

\newpage
% ----------------------------------------------------------------------
% TRANG 2 (MẶT SAU) -- DÒNG KẺ TIẾP THEO
% ----------------------------------------------------------------------
\noindent
\textbf{Học sinh:} Hoàng Nguyên \quad -- \quad \textbf{Mã đề:} \fbox{\textbf{009}} \hfill \textit{Trang 2/2}\par
\vspace{-0.1cm}
\noindent\hrulefill
\vspace{0.3cm}

% --- CÂU 2 (TRANG 2) ---
\begin{ex}
	Thực hiện phép tính:
	\begin{enumerate}[a)]
		\item $A = -5a^2b^3 \cdot 2ab^2$
		\item $B = 4a^3b \cdot 3ab^2 - 17a^4b^3$
		\item $C = 24a^7b^4 : 6a^4b^2 - 9a^3b^2$
		\item $D = 3a^2b^4 + 14a^4b^5 : 2a^2b$
	\end{enumerate}
	\loigiai{
		\begin{enumerate}[a)]
			\item Ta có:
			\begin{align*}
				A &= -5a^2b^3 \cdot 2ab^2 \\
				&= (-5 \cdot 2) \cdot (a^2 \cdot a) \cdot (b^3 \cdot b^2) \\
				&= -10a^3b^5
			\end{align*}
			\item Ta có:
			\begin{align*}
				B &= 4a^3b \cdot 3ab^2 - 17a^4b^3 \\
				&= 12a^4b^3 - 17a^4b^3 \\
				&= (12 - 17)a^4b^3 \\
				&= -5a^4b^3
			\end{align*}
			\item Ta có:
			\begin{align*}
				C &= 24a^7b^4 : 6a^4b^2 - 9a^3b^2 \\
				&= 4a^3b^2 - 9a^3b^2 \\
				&= (4 - 9)a^3b^2 \\
				&= -5a^3b^2
			\end{align*}
			\item Ta có:
			\begin{align*}
				D &= 3a^2b^4 + 14a^4b^5 : 2a^2b \\
				&= 3a^2b^4 + 7a^2b^4 \\
				&= 10a^2b^4
			\end{align*}
		\end{enumerate}
	}
\end{ex}

\newpage

% ======================================================================
% HỌC SINH: Hoàng Nhật | MÃ: HS007 | MÃ ĐỀ: 002
% ======================================================================
\noindent
\begin{tcolorbox}[colback=white,colframe=black,arc=0mm,boxrule=0.8pt,top=4pt,bottom=4pt,left=6pt,right=6pt]
	\begin{minipage}[c]{0.44\textwidth}
		{\bfseries\large LỚP TOÁN THẦY THÔNG}\\[3pt]
		\textbf{Họ và tên:} \textbf{\large Hoàng Nhật}\\[3pt]
		\textbf{Lớp:} LỚP 10 \quad -- \quad \textbf{Mã HS:} \texttt{HS007}
	\end{minipage}
	\hfill
	\begin{minipage}[c]{0.22\textwidth}
		\begin{tcolorbox}[colback=white,colframe=black,arc=0mm,boxrule=0.6pt,height=1.8cm,halign=center,valign=top,top=2pt]
			{\bfseries\normalsize Điểm}
		\end{tcolorbox}
	\end{minipage}
	\hfill
	\begin{minipage}[c]{0.30\textwidth}
		\raggedleft
		\textit{Thời gian làm bài: 30 phút}\\[4pt]
		\textbf{Mã đề thi:} \fbox{\textbf{\Large\strut\ 002\ }}
	\end{minipage}
\end{tcolorbox}
\vspace{0.25cm}

\noindent\textbf{\large ĐỀ BÀI (MÃ ĐỀ: 002)}
\vspace{0.1cm}

% --- CÂU 1 (TRANG 1) ---
\begin{ex}
	Thực hiện phép tính:
	\begin{enumerate}[a)]
		\item $A = 3x^3y \cdot 4xy^2$
		\item $B = 3x^2y^3 \cdot 2x^2y - 11x^4y^4$
		\item $C = 18x^6y^5 : 3x^2y^2 - 12x^4y^3$
		\item $D = 4x^3y^2 + 8x^5y^3 : 2x^2y$
	\end{enumerate}
	\loigiai{
		\begin{enumerate}[a)]
			\item Ta có:
			\begin{align*}
				A &= 3x^3y \cdot 4xy^2 \\
				&= (3 \cdot 4) \cdot (x^3 \cdot x) \cdot (y \cdot y^2) \\
				&= 12x^4y^3
			\end{align*}
			\item Ta có:
			\begin{align*}
				B &= 3x^2y^3 \cdot 2x^2y - 11x^4y^4 \\
				&= 6x^4y^4 - 11x^4y^4 \\
				&= (6 - 11)x^4y^4 \\
				&= -5x^4y^4
			\end{align*}
			\item Ta có:
			\begin{align*}
				C &= 18x^6y^5 : 3x^2y^2 - 12x^4y^3 \\
				&= 6x^4y^3 - 12x^4y^3 \\
				&= (6 - 12)x^4y^3 \\
				&= -6x^4y^3
			\end{align*}
			\item Ta có:
			\begin{align*}
				D &= 4x^3y^2 + 8x^5y^3 : 2x^2y \\
				&= 4x^3y^2 + 4x^3y^2 \\
				&= 8x^3y^2
			\end{align*}
		\end{enumerate}
	}
\end{ex}

\newpage
% ----------------------------------------------------------------------
% TRANG 2 (MẶT SAU) -- DÒNG KẺ TIẾP THEO
% ----------------------------------------------------------------------
\noindent
\textbf{Học sinh:} Hoàng Nhật \quad -- \quad \textbf{Mã đề:} \fbox{\textbf{002}} \hfill \textit{Trang 2/2}\par
\vspace{-0.1cm}
\noindent\hrulefill
\vspace{0.3cm}

% --- CÂU 2 (TRANG 2) ---
\begin{ex}
	Thực hiện phép tính:
	\begin{enumerate}[a)]
		\item $A = -5a^2b^3 \cdot 2ab^2$
		\item $B = 4a^3b \cdot 3ab^2 - 17a^4b^3$
		\item $C = 24a^7b^4 : 6a^4b^2 - 9a^3b^2$
		\item $D = 3a^2b^4 + 14a^4b^5 : 2a^2b$
	\end{enumerate}
	\loigiai{
		\begin{enumerate}[a)]
			\item Ta có:
			\begin{align*}
				A &= -5a^2b^3 \cdot 2ab^2 \\
				&= (-5 \cdot 2) \cdot (a^2 \cdot a) \cdot (b^3 \cdot b^2) \\
				&= -10a^3b^5
			\end{align*}
			\item Ta có:
			\begin{align*}
				B &= 4a^3b \cdot 3ab^2 - 17a^4b^3 \\
				&= 12a^4b^3 - 17a^4b^3 \\
				&= (12 - 17)a^4b^3 \\
				&= -5a^4b^3
			\end{align*}
			\item Ta có:
			\begin{align*}
				C &= 24a^7b^4 : 6a^4b^2 - 9a^3b^2 \\
				&= 4a^3b^2 - 9a^3b^2 \\
				&= (4 - 9)a^3b^2 \\
				&= -5a^3b^2
			\end{align*}
			\item Ta có:
			\begin{align*}
				D &= 3a^2b^4 + 14a^4b^5 : 2a^2b \\
				&= 3a^2b^4 + 7a^2b^4 \\
				&= 10a^2b^4
			\end{align*}
		\end{enumerate}
	}
\end{ex}

\newpage

% ======================================================================
% HỌC SINH: Hồ Dũng | MÃ: HS027 | MÃ ĐỀ: 005
% ======================================================================
\noindent
\begin{tcolorbox}[colback=white,colframe=black,arc=0mm,boxrule=0.8pt,top=4pt,bottom=4pt,left=6pt,right=6pt]
	\begin{minipage}[c]{0.44\textwidth}
		{\bfseries\large LỚP TOÁN THẦY THÔNG}\\[3pt]
		\textbf{Họ và tên:} \textbf{\large Hồ Dũng}\\[3pt]
		\textbf{Lớp:} LỚP 10 \quad -- \quad \textbf{Mã HS:} \texttt{HS027}
	\end{minipage}
	\hfill
	\begin{minipage}[c]{0.22\textwidth}
		\begin{tcolorbox}[colback=white,colframe=black,arc=0mm,boxrule=0.6pt,height=1.8cm,halign=center,valign=top,top=2pt]
			{\bfseries\normalsize Điểm}
		\end{tcolorbox}
	\end{minipage}
	\hfill
	\begin{minipage}[c]{0.30\textwidth}
		\raggedleft
		\textit{Thời gian làm bài: 30 phút}\\[4pt]
		\textbf{Mã đề thi:} \fbox{\textbf{\Large\strut\ 005\ }}
	\end{minipage}
\end{tcolorbox}
\vspace{0.25cm}

\noindent\textbf{\large ĐỀ BÀI (MÃ ĐỀ: 005)}
\vspace{0.1cm}

% --- CÂU 1 (TRANG 1) ---
\begin{ex}
	Thực hiện phép tính:
	\begin{enumerate}[a)]
		\item $A = 3x^3y \cdot 4xy^2$
		\item $B = 3x^2y^3 \cdot 2x^2y - 11x^4y^4$
		\item $C = 18x^6y^5 : 3x^2y^2 - 12x^4y^3$
		\item $D = 4x^3y^2 + 8x^5y^3 : 2x^2y$
	\end{enumerate}
	\loigiai{
		\begin{enumerate}[a)]
			\item Ta có:
			\begin{align*}
				A &= 3x^3y \cdot 4xy^2 \\
				&= (3 \cdot 4) \cdot (x^3 \cdot x) \cdot (y \cdot y^2) \\
				&= 12x^4y^3
			\end{align*}
			\item Ta có:
			\begin{align*}
				B &= 3x^2y^3 \cdot 2x^2y - 11x^4y^4 \\
				&= 6x^4y^4 - 11x^4y^4 \\
				&= (6 - 11)x^4y^4 \\
				&= -5x^4y^4
			\end{align*}
			\item Ta có:
			\begin{align*}
				C &= 18x^6y^5 : 3x^2y^2 - 12x^4y^3 \\
				&= 6x^4y^3 - 12x^4y^3 \\
				&= (6 - 12)x^4y^3 \\
				&= -6x^4y^3
			\end{align*}
			\item Ta có:
			\begin{align*}
				D &= 4x^3y^2 + 8x^5y^3 : 2x^2y \\
				&= 4x^3y^2 + 4x^3y^2 \\
				&= 8x^3y^2
			\end{align*}
		\end{enumerate}
	}
\end{ex}

\newpage
% ----------------------------------------------------------------------
% TRANG 2 (MẶT SAU) -- DÒNG KẺ TIẾP THEO
% ----------------------------------------------------------------------
\noindent
\textbf{Học sinh:} Hồ Dũng \quad -- \quad \textbf{Mã đề:} \fbox{\textbf{005}} \hfill \textit{Trang 2/2}\par
\vspace{-0.1cm}
\noindent\hrulefill
\vspace{0.3cm}

% --- CÂU 2 (TRANG 2) ---
\begin{ex}
	Thực hiện phép tính:
	\begin{enumerate}[a)]
		\item $A = -5a^2b^3 \cdot 2ab^2$
		\item $B = 4a^3b \cdot 3ab^2 - 17a^4b^3$
		\item $C = 24a^7b^4 : 6a^4b^2 - 9a^3b^2$
		\item $D = 3a^2b^4 + 14a^4b^5 : 2a^2b$
	\end{enumerate}
	\loigiai{
		\begin{enumerate}[a)]
			\item Ta có:
			\begin{align*}
				A &= -5a^2b^3 \cdot 2ab^2 \\
				&= (-5 \cdot 2) \cdot (a^2 \cdot a) \cdot (b^3 \cdot b^2) \\
				&= -10a^3b^5
			\end{align*}
			\item Ta có:
			\begin{align*}
				B &= 4a^3b \cdot 3ab^2 - 17a^4b^3 \\
				&= 12a^4b^3 - 17a^4b^3 \\
				&= (12 - 17)a^4b^3 \\
				&= -5a^4b^3
			\end{align*}
			\item Ta có:
			\begin{align*}
				C &= 24a^7b^4 : 6a^4b^2 - 9a^3b^2 \\
				&= 4a^3b^2 - 9a^3b^2 \\
				&= (4 - 9)a^3b^2 \\
				&= -5a^3b^2
			\end{align*}
			\item Ta có:
			\begin{align*}
				D &= 3a^2b^4 + 14a^4b^5 : 2a^2b \\
				&= 3a^2b^4 + 7a^2b^4 \\
				&= 10a^2b^4
			\end{align*}
		\end{enumerate}
	}
\end{ex}

\newpage

% ======================================================================
% HỌC SINH: Huyền Thương | MÃ: HS001 | MÃ ĐỀ: 007
% ======================================================================
\noindent
\begin{tcolorbox}[colback=white,colframe=black,arc=0mm,boxrule=0.8pt,top=4pt,bottom=4pt,left=6pt,right=6pt]
	\begin{minipage}[c]{0.44\textwidth}
		{\bfseries\large LỚP TOÁN THẦY THÔNG}\\[3pt]
		\textbf{Họ và tên:} \textbf{\large Huyền Thương}\\[3pt]
		\textbf{Lớp:} LỚP 10 \quad -- \quad \textbf{Mã HS:} \texttt{HS001}
	\end{minipage}
	\hfill
	\begin{minipage}[c]{0.22\textwidth}
		\begin{tcolorbox}[colback=white,colframe=black,arc=0mm,boxrule=0.6pt,height=1.8cm,halign=center,valign=top,top=2pt]
			{\bfseries\normalsize Điểm}
		\end{tcolorbox}
	\end{minipage}
	\hfill
	\begin{minipage}[c]{0.30\textwidth}
		\raggedleft
		\textit{Thời gian làm bài: 30 phút}\\[4pt]
		\textbf{Mã đề thi:} \fbox{\textbf{\Large\strut\ 007\ }}
	\end{minipage}
\end{tcolorbox}
\vspace{0.25cm}

\noindent\textbf{\large ĐỀ BÀI (MÃ ĐỀ: 007)}
\vspace{0.1cm}

% --- CÂU 1 (TRANG 1) ---
\begin{ex}
	Thực hiện phép tính:
	\begin{enumerate}[a)]
		\item $A = 3x^3y \cdot 4xy^2$
		\item $B = 3x^2y^3 \cdot 2x^2y - 11x^4y^4$
		\item $C = 18x^6y^5 : 3x^2y^2 - 12x^4y^3$
		\item $D = 4x^3y^2 + 8x^5y^3 : 2x^2y$
	\end{enumerate}
	\loigiai{
		\begin{enumerate}[a)]
			\item Ta có:
			\begin{align*}
				A &= 3x^3y \cdot 4xy^2 \\
				&= (3 \cdot 4) \cdot (x^3 \cdot x) \cdot (y \cdot y^2) \\
				&= 12x^4y^3
			\end{align*}
			\item Ta có:
			\begin{align*}
				B &= 3x^2y^3 \cdot 2x^2y - 11x^4y^4 \\
				&= 6x^4y^4 - 11x^4y^4 \\
				&= (6 - 11)x^4y^4 \\
				&= -5x^4y^4
			\end{align*}
			\item Ta có:
			\begin{align*}
				C &= 18x^6y^5 : 3x^2y^2 - 12x^4y^3 \\
				&= 6x^4y^3 - 12x^4y^3 \\
				&= (6 - 12)x^4y^3 \\
				&= -6x^4y^3
			\end{align*}
			\item Ta có:
			\begin{align*}
				D &= 4x^3y^2 + 8x^5y^3 : 2x^2y \\
				&= 4x^3y^2 + 4x^3y^2 \\
				&= 8x^3y^2
			\end{align*}
		\end{enumerate}
	}
\end{ex}

\newpage
% ----------------------------------------------------------------------
% TRANG 2 (MẶT SAU) -- DÒNG KẺ TIẾP THEO
% ----------------------------------------------------------------------
\noindent
\textbf{Học sinh:} Huyền Thương \quad -- \quad \textbf{Mã đề:} \fbox{\textbf{007}} \hfill \textit{Trang 2/2}\par
\vspace{-0.1cm}
\noindent\hrulefill
\vspace{0.3cm}

% --- CÂU 2 (TRANG 2) ---
\begin{ex}
	Thực hiện phép tính:
	\begin{enumerate}[a)]
		\item $A = -5a^2b^3 \cdot 2ab^2$
		\item $B = 4a^3b \cdot 3ab^2 - 17a^4b^3$
		\item $C = 24a^7b^4 : 6a^4b^2 - 9a^3b^2$
		\item $D = 3a^2b^4 + 14a^4b^5 : 2a^2b$
	\end{enumerate}
	\loigiai{
		\begin{enumerate}[a)]
			\item Ta có:
			\begin{align*}
				A &= -5a^2b^3 \cdot 2ab^2 \\
				&= (-5 \cdot 2) \cdot (a^2 \cdot a) \cdot (b^3 \cdot b^2) \\
				&= -10a^3b^5
			\end{align*}
			\item Ta có:
			\begin{align*}
				B &= 4a^3b \cdot 3ab^2 - 17a^4b^3 \\
				&= 12a^4b^3 - 17a^4b^3 \\
				&= (12 - 17)a^4b^3 \\
				&= -5a^4b^3
			\end{align*}
			\item Ta có:
			\begin{align*}
				C &= 24a^7b^4 : 6a^4b^2 - 9a^3b^2 \\
				&= 4a^3b^2 - 9a^3b^2 \\
				&= (4 - 9)a^3b^2 \\
				&= -5a^3b^2
			\end{align*}
			\item Ta có:
			\begin{align*}
				D &= 3a^2b^4 + 14a^4b^5 : 2a^2b \\
				&= 3a^2b^4 + 7a^2b^4 \\
				&= 10a^2b^4
			\end{align*}
		\end{enumerate}
	}
\end{ex}

\newpage

% ======================================================================
% HỌC SINH: Khánh Ly | MÃ: HS021 | MÃ ĐỀ: 006
% ======================================================================
\noindent
\begin{tcolorbox}[colback=white,colframe=black,arc=0mm,boxrule=0.8pt,top=4pt,bottom=4pt,left=6pt,right=6pt]
	\begin{minipage}[c]{0.44\textwidth}
		{\bfseries\large LỚP TOÁN THẦY THÔNG}\\[3pt]
		\textbf{Họ và tên:} \textbf{\large Khánh Ly}\\[3pt]
		\textbf{Lớp:} LỚP 10 \quad -- \quad \textbf{Mã HS:} \texttt{HS021}
	\end{minipage}
	\hfill
	\begin{minipage}[c]{0.22\textwidth}
		\begin{tcolorbox}[colback=white,colframe=black,arc=0mm,boxrule=0.6pt,height=1.8cm,halign=center,valign=top,top=2pt]
			{\bfseries\normalsize Điểm}
		\end{tcolorbox}
	\end{minipage}
	\hfill
	\begin{minipage}[c]{0.30\textwidth}
		\raggedleft
		\textit{Thời gian làm bài: 30 phút}\\[4pt]
		\textbf{Mã đề thi:} \fbox{\textbf{\Large\strut\ 006\ }}
	\end{minipage}
\end{tcolorbox}
\vspace{0.25cm}

\noindent\textbf{\large ĐỀ BÀI (MÃ ĐỀ: 006)}
\vspace{0.1cm}

% --- CÂU 1 (TRANG 1) ---
\begin{ex}
	Thực hiện phép tính:
	\begin{enumerate}[a)]
		\item $A = 3x^3y \cdot 4xy^2$
		\item $B = 3x^2y^3 \cdot 2x^2y - 11x^4y^4$
		\item $C = 18x^6y^5 : 3x^2y^2 - 12x^4y^3$
		\item $D = 4x^3y^2 + 8x^5y^3 : 2x^2y$
	\end{enumerate}
	\loigiai{
		\begin{enumerate}[a)]
			\item Ta có:
			\begin{align*}
				A &= 3x^3y \cdot 4xy^2 \\
				&= (3 \cdot 4) \cdot (x^3 \cdot x) \cdot (y \cdot y^2) \\
				&= 12x^4y^3
			\end{align*}
			\item Ta có:
			\begin{align*}
				B &= 3x^2y^3 \cdot 2x^2y - 11x^4y^4 \\
				&= 6x^4y^4 - 11x^4y^4 \\
				&= (6 - 11)x^4y^4 \\
				&= -5x^4y^4
			\end{align*}
			\item Ta có:
			\begin{align*}
				C &= 18x^6y^5 : 3x^2y^2 - 12x^4y^3 \\
				&= 6x^4y^3 - 12x^4y^3 \\
				&= (6 - 12)x^4y^3 \\
				&= -6x^4y^3
			\end{align*}
			\item Ta có:
			\begin{align*}
				D &= 4x^3y^2 + 8x^5y^3 : 2x^2y \\
				&= 4x^3y^2 + 4x^3y^2 \\
				&= 8x^3y^2
			\end{align*}
		\end{enumerate}
	}
\end{ex}

\newpage
% ----------------------------------------------------------------------
% TRANG 2 (MẶT SAU) -- DÒNG KẺ TIẾP THEO
% ----------------------------------------------------------------------
\noindent
\textbf{Học sinh:} Khánh Ly \quad -- \quad \textbf{Mã đề:} \fbox{\textbf{006}} \hfill \textit{Trang 2/2}\par
\vspace{-0.1cm}
\noindent\hrulefill
\vspace{0.3cm}

% --- CÂU 2 (TRANG 2) ---
\begin{ex}
	Thực hiện phép tính:
	\begin{enumerate}[a)]
		\item $A = -5a^2b^3 \cdot 2ab^2$
		\item $B = 4a^3b \cdot 3ab^2 - 17a^4b^3$
		\item $C = 24a^7b^4 : 6a^4b^2 - 9a^3b^2$
		\item $D = 3a^2b^4 + 14a^4b^5 : 2a^2b$
	\end{enumerate}
	\loigiai{
		\begin{enumerate}[a)]
			\item Ta có:
			\begin{align*}
				A &= -5a^2b^3 \cdot 2ab^2 \\
				&= (-5 \cdot 2) \cdot (a^2 \cdot a) \cdot (b^3 \cdot b^2) \\
				&= -10a^3b^5
			\end{align*}
			\item Ta có:
			\begin{align*}
				B &= 4a^3b \cdot 3ab^2 - 17a^4b^3 \\
				&= 12a^4b^3 - 17a^4b^3 \\
				&= (12 - 17)a^4b^3 \\
				&= -5a^4b^3
			\end{align*}
			\item Ta có:
			\begin{align*}
				C &= 24a^7b^4 : 6a^4b^2 - 9a^3b^2 \\
				&= 4a^3b^2 - 9a^3b^2 \\
				&= (4 - 9)a^3b^2 \\
				&= -5a^3b^2
			\end{align*}
			\item Ta có:
			\begin{align*}
				D &= 3a^2b^4 + 14a^4b^5 : 2a^2b \\
				&= 3a^2b^4 + 7a^2b^4 \\
				&= 10a^2b^4
			\end{align*}
		\end{enumerate}
	}
\end{ex}

\newpage

% ======================================================================
% HỌC SINH: Khánh Ngọc | MÃ: HS024 | MÃ ĐỀ: 004
% ======================================================================
\noindent
\begin{tcolorbox}[colback=white,colframe=black,arc=0mm,boxrule=0.8pt,top=4pt,bottom=4pt,left=6pt,right=6pt]
	\begin{minipage}[c]{0.44\textwidth}
		{\bfseries\large LỚP TOÁN THẦY THÔNG}\\[3pt]
		\textbf{Họ và tên:} \textbf{\large Khánh Ngọc}\\[3pt]
		\textbf{Lớp:} LỚP 10 \quad -- \quad \textbf{Mã HS:} \texttt{HS024}
	\end{minipage}
	\hfill
	\begin{minipage}[c]{0.22\textwidth}
		\begin{tcolorbox}[colback=white,colframe=black,arc=0mm,boxrule=0.6pt,height=1.8cm,halign=center,valign=top,top=2pt]
			{\bfseries\normalsize Điểm}
		\end{tcolorbox}
	\end{minipage}
	\hfill
	\begin{minipage}[c]{0.30\textwidth}
		\raggedleft
		\textit{Thời gian làm bài: 30 phút}\\[4pt]
		\textbf{Mã đề thi:} \fbox{\textbf{\Large\strut\ 004\ }}
	\end{minipage}
\end{tcolorbox}
\vspace{0.25cm}

\noindent\textbf{\large ĐỀ BÀI (MÃ ĐỀ: 004)}
\vspace{0.1cm}

% --- CÂU 1 (TRANG 1) ---
\begin{ex}
	Thực hiện phép tính:
	\begin{enumerate}[a)]
		\item $A = 3x^3y \cdot 4xy^2$
		\item $B = 3x^2y^3 \cdot 2x^2y - 11x^4y^4$
		\item $C = 18x^6y^5 : 3x^2y^2 - 12x^4y^3$
		\item $D = 4x^3y^2 + 8x^5y^3 : 2x^2y$
	\end{enumerate}
	\loigiai{
		\begin{enumerate}[a)]
			\item Ta có:
			\begin{align*}
				A &= 3x^3y \cdot 4xy^2 \\
				&= (3 \cdot 4) \cdot (x^3 \cdot x) \cdot (y \cdot y^2) \\
				&= 12x^4y^3
			\end{align*}
			\item Ta có:
			\begin{align*}
				B &= 3x^2y^3 \cdot 2x^2y - 11x^4y^4 \\
				&= 6x^4y^4 - 11x^4y^4 \\
				&= (6 - 11)x^4y^4 \\
				&= -5x^4y^4
			\end{align*}
			\item Ta có:
			\begin{align*}
				C &= 18x^6y^5 : 3x^2y^2 - 12x^4y^3 \\
				&= 6x^4y^3 - 12x^4y^3 \\
				&= (6 - 12)x^4y^3 \\
				&= -6x^4y^3
			\end{align*}
			\item Ta có:
			\begin{align*}
				D &= 4x^3y^2 + 8x^5y^3 : 2x^2y \\
				&= 4x^3y^2 + 4x^3y^2 \\
				&= 8x^3y^2
			\end{align*}
		\end{enumerate}
	}
\end{ex}

\newpage
% ----------------------------------------------------------------------
% TRANG 2 (MẶT SAU) -- DÒNG KẺ TIẾP THEO
% ----------------------------------------------------------------------
\noindent
\textbf{Học sinh:} Khánh Ngọc \quad -- \quad \textbf{Mã đề:} \fbox{\textbf{004}} \hfill \textit{Trang 2/2}\par
\vspace{-0.1cm}
\noindent\hrulefill
\vspace{0.3cm}

% --- CÂU 2 (TRANG 2) ---
\begin{ex}
	Thực hiện phép tính:
	\begin{enumerate}[a)]
		\item $A = -5a^2b^3 \cdot 2ab^2$
		\item $B = 4a^3b \cdot 3ab^2 - 17a^4b^3$
		\item $C = 24a^7b^4 : 6a^4b^2 - 9a^3b^2$
		\item $D = 3a^2b^4 + 14a^4b^5 : 2a^2b$
	\end{enumerate}
	\loigiai{
		\begin{enumerate}[a)]
			\item Ta có:
			\begin{align*}
				A &= -5a^2b^3 \cdot 2ab^2 \\
				&= (-5 \cdot 2) \cdot (a^2 \cdot a) \cdot (b^3 \cdot b^2) \\
				&= -10a^3b^5
			\end{align*}
			\item Ta có:
			\begin{align*}
				B &= 4a^3b \cdot 3ab^2 - 17a^4b^3 \\
				&= 12a^4b^3 - 17a^4b^3 \\
				&= (12 - 17)a^4b^3 \\
				&= -5a^4b^3
			\end{align*}
			\item Ta có:
			\begin{align*}
				C &= 24a^7b^4 : 6a^4b^2 - 9a^3b^2 \\
				&= 4a^3b^2 - 9a^3b^2 \\
				&= (4 - 9)a^3b^2 \\
				&= -5a^3b^2
			\end{align*}
			\item Ta có:
			\begin{align*}
				D &= 3a^2b^4 + 14a^4b^5 : 2a^2b \\
				&= 3a^2b^4 + 7a^2b^4 \\
				&= 10a^2b^4
			\end{align*}
		\end{enumerate}
	}
\end{ex}

\newpage

% ======================================================================
% HỌC SINH: Khánh Thi | MÃ: HS023 | MÃ ĐỀ: 003
% ======================================================================
\noindent
\begin{tcolorbox}[colback=white,colframe=black,arc=0mm,boxrule=0.8pt,top=4pt,bottom=4pt,left=6pt,right=6pt]
	\begin{minipage}[c]{0.44\textwidth}
		{\bfseries\large LỚP TOÁN THẦY THÔNG}\\[3pt]
		\textbf{Họ và tên:} \textbf{\large Khánh Thi}\\[3pt]
		\textbf{Lớp:} LỚP 10 \quad -- \quad \textbf{Mã HS:} \texttt{HS023}
	\end{minipage}
	\hfill
	\begin{minipage}[c]{0.22\textwidth}
		\begin{tcolorbox}[colback=white,colframe=black,arc=0mm,boxrule=0.6pt,height=1.8cm,halign=center,valign=top,top=2pt]
			{\bfseries\normalsize Điểm}
		\end{tcolorbox}
	\end{minipage}
	\hfill
	\begin{minipage}[c]{0.30\textwidth}
		\raggedleft
		\textit{Thời gian làm bài: 30 phút}\\[4pt]
		\textbf{Mã đề thi:} \fbox{\textbf{\Large\strut\ 003\ }}
	\end{minipage}
\end{tcolorbox}
\vspace{0.25cm}

\noindent\textbf{\large ĐỀ BÀI (MÃ ĐỀ: 003)}
\vspace{0.1cm}

% --- CÂU 1 (TRANG 1) ---
\begin{ex}
	Thực hiện phép tính:
	\begin{enumerate}[a)]
		\item $A = 3x^3y \cdot 4xy^2$
		\item $B = 3x^2y^3 \cdot 2x^2y - 11x^4y^4$
		\item $C = 18x^6y^5 : 3x^2y^2 - 12x^4y^3$
		\item $D = 4x^3y^2 + 8x^5y^3 : 2x^2y$
	\end{enumerate}
	\loigiai{
		\begin{enumerate}[a)]
			\item Ta có:
			\begin{align*}
				A &= 3x^3y \cdot 4xy^2 \\
				&= (3 \cdot 4) \cdot (x^3 \cdot x) \cdot (y \cdot y^2) \\
				&= 12x^4y^3
			\end{align*}
			\item Ta có:
			\begin{align*}
				B &= 3x^2y^3 \cdot 2x^2y - 11x^4y^4 \\
				&= 6x^4y^4 - 11x^4y^4 \\
				&= (6 - 11)x^4y^4 \\
				&= -5x^4y^4
			\end{align*}
			\item Ta có:
			\begin{align*}
				C &= 18x^6y^5 : 3x^2y^2 - 12x^4y^3 \\
				&= 6x^4y^3 - 12x^4y^3 \\
				&= (6 - 12)x^4y^3 \\
				&= -6x^4y^3
			\end{align*}
			\item Ta có:
			\begin{align*}
				D &= 4x^3y^2 + 8x^5y^3 : 2x^2y \\
				&= 4x^3y^2 + 4x^3y^2 \\
				&= 8x^3y^2
			\end{align*}
		\end{enumerate}
	}
\end{ex}

\newpage
% ----------------------------------------------------------------------
% TRANG 2 (MẶT SAU) -- DÒNG KẺ TIẾP THEO
% ----------------------------------------------------------------------
\noindent
\textbf{Học sinh:} Khánh Thi \quad -- \quad \textbf{Mã đề:} \fbox{\textbf{003}} \hfill \textit{Trang 2/2}\par
\vspace{-0.1cm}
\noindent\hrulefill
\vspace{0.3cm}

% --- CÂU 2 (TRANG 2) ---
\begin{ex}
	Thực hiện phép tính:
	\begin{enumerate}[a)]
		\item $A = -5a^2b^3 \cdot 2ab^2$
		\item $B = 4a^3b \cdot 3ab^2 - 17a^4b^3$
		\item $C = 24a^7b^4 : 6a^4b^2 - 9a^3b^2$
		\item $D = 3a^2b^4 + 14a^4b^5 : 2a^2b$
	\end{enumerate}
	\loigiai{
		\begin{enumerate}[a)]
			\item Ta có:
			\begin{align*}
				A &= -5a^2b^3 \cdot 2ab^2 \\
				&= (-5 \cdot 2) \cdot (a^2 \cdot a) \cdot (b^3 \cdot b^2) \\
				&= -10a^3b^5
			\end{align*}
			\item Ta có:
			\begin{align*}
				B &= 4a^3b \cdot 3ab^2 - 17a^4b^3 \\
				&= 12a^4b^3 - 17a^4b^3 \\
				&= (12 - 17)a^4b^3 \\
				&= -5a^4b^3
			\end{align*}
			\item Ta có:
			\begin{align*}
				C &= 24a^7b^4 : 6a^4b^2 - 9a^3b^2 \\
				&= 4a^3b^2 - 9a^3b^2 \\
				&= (4 - 9)a^3b^2 \\
				&= -5a^3b^2
			\end{align*}
			\item Ta có:
			\begin{align*}
				D &= 3a^2b^4 + 14a^4b^5 : 2a^2b \\
				&= 3a^2b^4 + 7a^2b^4 \\
				&= 10a^2b^4
			\end{align*}
		\end{enumerate}
	}
\end{ex}

\newpage

% ======================================================================
% HỌC SINH: Khôi 20/7 | MÃ: HS025 | MÃ ĐỀ: 005
% ======================================================================
\noindent
\begin{tcolorbox}[colback=white,colframe=black,arc=0mm,boxrule=0.8pt,top=4pt,bottom=4pt,left=6pt,right=6pt]
	\begin{minipage}[c]{0.44\textwidth}
		{\bfseries\large LỚP TOÁN THẦY THÔNG}\\[3pt]
		\textbf{Họ và tên:} \textbf{\large Khôi 20/7}\\[3pt]
		\textbf{Lớp:} LỚP 10 \quad -- \quad \textbf{Mã HS:} \texttt{HS025}
	\end{minipage}
	\hfill
	\begin{minipage}[c]{0.22\textwidth}
		\begin{tcolorbox}[colback=white,colframe=black,arc=0mm,boxrule=0.6pt,height=1.8cm,halign=center,valign=top,top=2pt]
			{\bfseries\normalsize Điểm}
		\end{tcolorbox}
	\end{minipage}
	\hfill
	\begin{minipage}[c]{0.30\textwidth}
		\raggedleft
		\textit{Thời gian làm bài: 30 phút}\\[4pt]
		\textbf{Mã đề thi:} \fbox{\textbf{\Large\strut\ 005\ }}
	\end{minipage}
\end{tcolorbox}
\vspace{0.25cm}

\noindent\textbf{\large ĐỀ BÀI (MÃ ĐỀ: 005)}
\vspace{0.1cm}

% --- CÂU 1 (TRANG 1) ---
\begin{ex}
	Thực hiện phép tính:
	\begin{enumerate}[a)]
		\item $A = 3x^3y \cdot 4xy^2$
		\item $B = 3x^2y^3 \cdot 2x^2y - 11x^4y^4$
		\item $C = 18x^6y^5 : 3x^2y^2 - 12x^4y^3$
		\item $D = 4x^3y^2 + 8x^5y^3 : 2x^2y$
	\end{enumerate}
	\loigiai{
		\begin{enumerate}[a)]
			\item Ta có:
			\begin{align*}
				A &= 3x^3y \cdot 4xy^2 \\
				&= (3 \cdot 4) \cdot (x^3 \cdot x) \cdot (y \cdot y^2) \\
				&= 12x^4y^3
			\end{align*}
			\item Ta có:
			\begin{align*}
				B &= 3x^2y^3 \cdot 2x^2y - 11x^4y^4 \\
				&= 6x^4y^4 - 11x^4y^4 \\
				&= (6 - 11)x^4y^4 \\
				&= -5x^4y^4
			\end{align*}
			\item Ta có:
			\begin{align*}
				C &= 18x^6y^5 : 3x^2y^2 - 12x^4y^3 \\
				&= 6x^4y^3 - 12x^4y^3 \\
				&= (6 - 12)x^4y^3 \\
				&= -6x^4y^3
			\end{align*}
			\item Ta có:
			\begin{align*}
				D &= 4x^3y^2 + 8x^5y^3 : 2x^2y \\
				&= 4x^3y^2 + 4x^3y^2 \\
				&= 8x^3y^2
			\end{align*}
		\end{enumerate}
	}
\end{ex}

\newpage
% ----------------------------------------------------------------------
% TRANG 2 (MẶT SAU) -- DÒNG KẺ TIẾP THEO
% ----------------------------------------------------------------------
\noindent
\textbf{Học sinh:} Khôi 20/7 \quad -- \quad \textbf{Mã đề:} \fbox{\textbf{005}} \hfill \textit{Trang 2/2}\par
\vspace{-0.1cm}
\noindent\hrulefill
\vspace{0.3cm}

% --- CÂU 2 (TRANG 2) ---
\begin{ex}
	Thực hiện phép tính:
	\begin{enumerate}[a)]
		\item $A = -5a^2b^3 \cdot 2ab^2$
		\item $B = 4a^3b \cdot 3ab^2 - 17a^4b^3$
		\item $C = 24a^7b^4 : 6a^4b^2 - 9a^3b^2$
		\item $D = 3a^2b^4 + 14a^4b^5 : 2a^2b$
	\end{enumerate}
	\loigiai{
		\begin{enumerate}[a)]
			\item Ta có:
			\begin{align*}
				A &= -5a^2b^3 \cdot 2ab^2 \\
				&= (-5 \cdot 2) \cdot (a^2 \cdot a) \cdot (b^3 \cdot b^2) \\
				&= -10a^3b^5
			\end{align*}
			\item Ta có:
			\begin{align*}
				B &= 4a^3b \cdot 3ab^2 - 17a^4b^3 \\
				&= 12a^4b^3 - 17a^4b^3 \\
				&= (12 - 17)a^4b^3 \\
				&= -5a^4b^3
			\end{align*}
			\item Ta có:
			\begin{align*}
				C &= 24a^7b^4 : 6a^4b^2 - 9a^3b^2 \\
				&= 4a^3b^2 - 9a^3b^2 \\
				&= (4 - 9)a^3b^2 \\
				&= -5a^3b^2
			\end{align*}
			\item Ta có:
			\begin{align*}
				D &= 3a^2b^4 + 14a^4b^5 : 2a^2b \\
				&= 3a^2b^4 + 7a^2b^4 \\
				&= 10a^2b^4
			\end{align*}
		\end{enumerate}
	}
\end{ex}

\newpage

% ======================================================================
% HỌC SINH: L10-phương Vy | MÃ: HS109999 | MÃ ĐỀ: 010
% ======================================================================
\noindent
\begin{tcolorbox}[colback=white,colframe=black,arc=0mm,boxrule=0.8pt,top=4pt,bottom=4pt,left=6pt,right=6pt]
	\begin{minipage}[c]{0.44\textwidth}
		{\bfseries\large LỚP TOÁN THẦY THÔNG}\\[3pt]
		\textbf{Họ và tên:} \textbf{\large L10-phương Vy}\\[3pt]
		\textbf{Lớp:} LỚP 10 \quad -- \quad \textbf{Mã HS:} \texttt{HS109999}
	\end{minipage}
	\hfill
	\begin{minipage}[c]{0.22\textwidth}
		\begin{tcolorbox}[colback=white,colframe=black,arc=0mm,boxrule=0.6pt,height=1.8cm,halign=center,valign=top,top=2pt]
			{\bfseries\normalsize Điểm}
		\end{tcolorbox}
	\end{minipage}
	\hfill
	\begin{minipage}[c]{0.30\textwidth}
		\raggedleft
		\textit{Thời gian làm bài: 30 phút}\\[4pt]
		\textbf{Mã đề thi:} \fbox{\textbf{\Large\strut\ 010\ }}
	\end{minipage}
\end{tcolorbox}
\vspace{0.25cm}

\noindent\textbf{\large ĐỀ BÀI (MÃ ĐỀ: 010)}
\vspace{0.1cm}

% --- CÂU 1 (TRANG 1) ---
\begin{ex}
	Thực hiện phép tính:
	\begin{enumerate}[a)]
		\item $A = 3x^3y \cdot 4xy^2$
		\item $B = 3x^2y^3 \cdot 2x^2y - 11x^4y^4$
		\item $C = 18x^6y^5 : 3x^2y^2 - 12x^4y^3$
		\item $D = 4x^3y^2 + 8x^5y^3 : 2x^2y$
	\end{enumerate}
	\loigiai{
		\begin{enumerate}[a)]
			\item Ta có:
			\begin{align*}
				A &= 3x^3y \cdot 4xy^2 \\
				&= (3 \cdot 4) \cdot (x^3 \cdot x) \cdot (y \cdot y^2) \\
				&= 12x^4y^3
			\end{align*}
			\item Ta có:
			\begin{align*}
				B &= 3x^2y^3 \cdot 2x^2y - 11x^4y^4 \\
				&= 6x^4y^4 - 11x^4y^4 \\
				&= (6 - 11)x^4y^4 \\
				&= -5x^4y^4
			\end{align*}
			\item Ta có:
			\begin{align*}
				C &= 18x^6y^5 : 3x^2y^2 - 12x^4y^3 \\
				&= 6x^4y^3 - 12x^4y^3 \\
				&= (6 - 12)x^4y^3 \\
				&= -6x^4y^3
			\end{align*}
			\item Ta có:
			\begin{align*}
				D &= 4x^3y^2 + 8x^5y^3 : 2x^2y \\
				&= 4x^3y^2 + 4x^3y^2 \\
				&= 8x^3y^2
			\end{align*}
		\end{enumerate}
	}
\end{ex}

\newpage
% ----------------------------------------------------------------------
% TRANG 2 (MẶT SAU) -- DÒNG KẺ TIẾP THEO
% ----------------------------------------------------------------------
\noindent
\textbf{Học sinh:} L10-phương Vy \quad -- \quad \textbf{Mã đề:} \fbox{\textbf{010}} \hfill \textit{Trang 2/2}\par
\vspace{-0.1cm}
\noindent\hrulefill
\vspace{0.3cm}

% --- CÂU 2 (TRANG 2) ---
\begin{ex}
	Thực hiện phép tính:
	\begin{enumerate}[a)]
		\item $A = -5a^2b^3 \cdot 2ab^2$
		\item $B = 4a^3b \cdot 3ab^2 - 17a^4b^3$
		\item $C = 24a^7b^4 : 6a^4b^2 - 9a^3b^2$
		\item $D = 3a^2b^4 + 14a^4b^5 : 2a^2b$
	\end{enumerate}
	\loigiai{
		\begin{enumerate}[a)]
			\item Ta có:
			\begin{align*}
				A &= -5a^2b^3 \cdot 2ab^2 \\
				&= (-5 \cdot 2) \cdot (a^2 \cdot a) \cdot (b^3 \cdot b^2) \\
				&= -10a^3b^5
			\end{align*}
			\item Ta có:
			\begin{align*}
				B &= 4a^3b \cdot 3ab^2 - 17a^4b^3 \\
				&= 12a^4b^3 - 17a^4b^3 \\
				&= (12 - 17)a^4b^3 \\
				&= -5a^4b^3
			\end{align*}
			\item Ta có:
			\begin{align*}
				C &= 24a^7b^4 : 6a^4b^2 - 9a^3b^2 \\
				&= 4a^3b^2 - 9a^3b^2 \\
				&= (4 - 9)a^3b^2 \\
				&= -5a^3b^2
			\end{align*}
			\item Ta có:
			\begin{align*}
				D &= 3a^2b^4 + 14a^4b^5 : 2a^2b \\
				&= 3a^2b^4 + 7a^2b^4 \\
				&= 10a^2b^4
			\end{align*}
		\end{enumerate}
	}
\end{ex}

\newpage

% ======================================================================
% HỌC SINH: L10-uyển Nhi | MÃ: HS1088888 | MÃ ĐỀ: 007
% ======================================================================
\noindent
\begin{tcolorbox}[colback=white,colframe=black,arc=0mm,boxrule=0.8pt,top=4pt,bottom=4pt,left=6pt,right=6pt]
	\begin{minipage}[c]{0.44\textwidth}
		{\bfseries\large LỚP TOÁN THẦY THÔNG}\\[3pt]
		\textbf{Họ và tên:} \textbf{\large L10-uyển Nhi}\\[3pt]
		\textbf{Lớp:} LỚP 10 \quad -- \quad \textbf{Mã HS:} \texttt{HS1088888}
	\end{minipage}
	\hfill
	\begin{minipage}[c]{0.22\textwidth}
		\begin{tcolorbox}[colback=white,colframe=black,arc=0mm,boxrule=0.6pt,height=1.8cm,halign=center,valign=top,top=2pt]
			{\bfseries\normalsize Điểm}
		\end{tcolorbox}
	\end{minipage}
	\hfill
	\begin{minipage}[c]{0.30\textwidth}
		\raggedleft
		\textit{Thời gian làm bài: 30 phút}\\[4pt]
		\textbf{Mã đề thi:} \fbox{\textbf{\Large\strut\ 007\ }}
	\end{minipage}
\end{tcolorbox}
\vspace{0.25cm}

\noindent\textbf{\large ĐỀ BÀI (MÃ ĐỀ: 007)}
\vspace{0.1cm}

% --- CÂU 1 (TRANG 1) ---
\begin{ex}
	Thực hiện phép tính:
	\begin{enumerate}[a)]
		\item $A = 3x^3y \cdot 4xy^2$
		\item $B = 3x^2y^3 \cdot 2x^2y - 11x^4y^4$
		\item $C = 18x^6y^5 : 3x^2y^2 - 12x^4y^3$
		\item $D = 4x^3y^2 + 8x^5y^3 : 2x^2y$
	\end{enumerate}
	\loigiai{
		\begin{enumerate}[a)]
			\item Ta có:
			\begin{align*}
				A &= 3x^3y \cdot 4xy^2 \\
				&= (3 \cdot 4) \cdot (x^3 \cdot x) \cdot (y \cdot y^2) \\
				&= 12x^4y^3
			\end{align*}
			\item Ta có:
			\begin{align*}
				B &= 3x^2y^3 \cdot 2x^2y - 11x^4y^4 \\
				&= 6x^4y^4 - 11x^4y^4 \\
				&= (6 - 11)x^4y^4 \\
				&= -5x^4y^4
			\end{align*}
			\item Ta có:
			\begin{align*}
				C &= 18x^6y^5 : 3x^2y^2 - 12x^4y^3 \\
				&= 6x^4y^3 - 12x^4y^3 \\
				&= (6 - 12)x^4y^3 \\
				&= -6x^4y^3
			\end{align*}
			\item Ta có:
			\begin{align*}
				D &= 4x^3y^2 + 8x^5y^3 : 2x^2y \\
				&= 4x^3y^2 + 4x^3y^2 \\
				&= 8x^3y^2
			\end{align*}
		\end{enumerate}
	}
\end{ex}

\newpage
% ----------------------------------------------------------------------
% TRANG 2 (MẶT SAU) -- DÒNG KẺ TIẾP THEO
% ----------------------------------------------------------------------
\noindent
\textbf{Học sinh:} L10-uyển Nhi \quad -- \quad \textbf{Mã đề:} \fbox{\textbf{007}} \hfill \textit{Trang 2/2}\par
\vspace{-0.1cm}
\noindent\hrulefill
\vspace{0.3cm}

% --- CÂU 2 (TRANG 2) ---
\begin{ex}
	Thực hiện phép tính:
	\begin{enumerate}[a)]
		\item $A = -5a^2b^3 \cdot 2ab^2$
		\item $B = 4a^3b \cdot 3ab^2 - 17a^4b^3$
		\item $C = 24a^7b^4 : 6a^4b^2 - 9a^3b^2$
		\item $D = 3a^2b^4 + 14a^4b^5 : 2a^2b$
	\end{enumerate}
	\loigiai{
		\begin{enumerate}[a)]
			\item Ta có:
			\begin{align*}
				A &= -5a^2b^3 \cdot 2ab^2 \\
				&= (-5 \cdot 2) \cdot (a^2 \cdot a) \cdot (b^3 \cdot b^2) \\
				&= -10a^3b^5
			\end{align*}
			\item Ta có:
			\begin{align*}
				B &= 4a^3b \cdot 3ab^2 - 17a^4b^3 \\
				&= 12a^4b^3 - 17a^4b^3 \\
				&= (12 - 17)a^4b^3 \\
				&= -5a^4b^3
			\end{align*}
			\item Ta có:
			\begin{align*}
				C &= 24a^7b^4 : 6a^4b^2 - 9a^3b^2 \\
				&= 4a^3b^2 - 9a^3b^2 \\
				&= (4 - 9)a^3b^2 \\
				&= -5a^3b^2
			\end{align*}
			\item Ta có:
			\begin{align*}
				D &= 3a^2b^4 + 14a^4b^5 : 2a^2b \\
				&= 3a^2b^4 + 7a^2b^4 \\
				&= 10a^2b^4
			\end{align*}
		\end{enumerate}
	}
\end{ex}

\newpage

% ======================================================================
% HỌC SINH: Mai Nhi | MÃ: HS020 | MÃ ĐỀ: 009
% ======================================================================
\noindent
\begin{tcolorbox}[colback=white,colframe=black,arc=0mm,boxrule=0.8pt,top=4pt,bottom=4pt,left=6pt,right=6pt]
	\begin{minipage}[c]{0.44\textwidth}
		{\bfseries\large LỚP TOÁN THẦY THÔNG}\\[3pt]
		\textbf{Họ và tên:} \textbf{\large Mai Nhi}\\[3pt]
		\textbf{Lớp:} LỚP 10 \quad -- \quad \textbf{Mã HS:} \texttt{HS020}
	\end{minipage}
	\hfill
	\begin{minipage}[c]{0.22\textwidth}
		\begin{tcolorbox}[colback=white,colframe=black,arc=0mm,boxrule=0.6pt,height=1.8cm,halign=center,valign=top,top=2pt]
			{\bfseries\normalsize Điểm}
		\end{tcolorbox}
	\end{minipage}
	\hfill
	\begin{minipage}[c]{0.30\textwidth}
		\raggedleft
		\textit{Thời gian làm bài: 30 phút}\\[4pt]
		\textbf{Mã đề thi:} \fbox{\textbf{\Large\strut\ 009\ }}
	\end{minipage}
\end{tcolorbox}
\vspace{0.25cm}

\noindent\textbf{\large ĐỀ BÀI (MÃ ĐỀ: 009)}
\vspace{0.1cm}

% --- CÂU 1 (TRANG 1) ---
\begin{ex}
	Thực hiện phép tính:
	\begin{enumerate}[a)]
		\item $A = 3x^3y \cdot 4xy^2$
		\item $B = 3x^2y^3 \cdot 2x^2y - 11x^4y^4$
		\item $C = 18x^6y^5 : 3x^2y^2 - 12x^4y^3$
		\item $D = 4x^3y^2 + 8x^5y^3 : 2x^2y$
	\end{enumerate}
	\loigiai{
		\begin{enumerate}[a)]
			\item Ta có:
			\begin{align*}
				A &= 3x^3y \cdot 4xy^2 \\
				&= (3 \cdot 4) \cdot (x^3 \cdot x) \cdot (y \cdot y^2) \\
				&= 12x^4y^3
			\end{align*}
			\item Ta có:
			\begin{align*}
				B &= 3x^2y^3 \cdot 2x^2y - 11x^4y^4 \\
				&= 6x^4y^4 - 11x^4y^4 \\
				&= (6 - 11)x^4y^4 \\
				&= -5x^4y^4
			\end{align*}
			\item Ta có:
			\begin{align*}
				C &= 18x^6y^5 : 3x^2y^2 - 12x^4y^3 \\
				&= 6x^4y^3 - 12x^4y^3 \\
				&= (6 - 12)x^4y^3 \\
				&= -6x^4y^3
			\end{align*}
			\item Ta có:
			\begin{align*}
				D &= 4x^3y^2 + 8x^5y^3 : 2x^2y \\
				&= 4x^3y^2 + 4x^3y^2 \\
				&= 8x^3y^2
			\end{align*}
		\end{enumerate}
	}
\end{ex}

\newpage
% ----------------------------------------------------------------------
% TRANG 2 (MẶT SAU) -- DÒNG KẺ TIẾP THEO
% ----------------------------------------------------------------------
\noindent
\textbf{Học sinh:} Mai Nhi \quad -- \quad \textbf{Mã đề:} \fbox{\textbf{009}} \hfill \textit{Trang 2/2}\par
\vspace{-0.1cm}
\noindent\hrulefill
\vspace{0.3cm}

% --- CÂU 2 (TRANG 2) ---
\begin{ex}
	Thực hiện phép tính:
	\begin{enumerate}[a)]
		\item $A = -5a^2b^3 \cdot 2ab^2$
		\item $B = 4a^3b \cdot 3ab^2 - 17a^4b^3$
		\item $C = 24a^7b^4 : 6a^4b^2 - 9a^3b^2$
		\item $D = 3a^2b^4 + 14a^4b^5 : 2a^2b$
	\end{enumerate}
	\loigiai{
		\begin{enumerate}[a)]
			\item Ta có:
			\begin{align*}
				A &= -5a^2b^3 \cdot 2ab^2 \\
				&= (-5 \cdot 2) \cdot (a^2 \cdot a) \cdot (b^3 \cdot b^2) \\
				&= -10a^3b^5
			\end{align*}
			\item Ta có:
			\begin{align*}
				B &= 4a^3b \cdot 3ab^2 - 17a^4b^3 \\
				&= 12a^4b^3 - 17a^4b^3 \\
				&= (12 - 17)a^4b^3 \\
				&= -5a^4b^3
			\end{align*}
			\item Ta có:
			\begin{align*}
				C &= 24a^7b^4 : 6a^4b^2 - 9a^3b^2 \\
				&= 4a^3b^2 - 9a^3b^2 \\
				&= (4 - 9)a^3b^2 \\
				&= -5a^3b^2
			\end{align*}
			\item Ta có:
			\begin{align*}
				D &= 3a^2b^4 + 14a^4b^5 : 2a^2b \\
				&= 3a^2b^4 + 7a^2b^4 \\
				&= 10a^2b^4
			\end{align*}
		\end{enumerate}
	}
\end{ex}

\newpage

% ======================================================================
% HỌC SINH: Mai Phương | MÃ: HS010 | MÃ ĐỀ: 001
% ======================================================================
\noindent
\begin{tcolorbox}[colback=white,colframe=black,arc=0mm,boxrule=0.8pt,top=4pt,bottom=4pt,left=6pt,right=6pt]
	\begin{minipage}[c]{0.44\textwidth}
		{\bfseries\large LỚP TOÁN THẦY THÔNG}\\[3pt]
		\textbf{Họ và tên:} \textbf{\large Mai Phương}\\[3pt]
		\textbf{Lớp:} LỚP 10 \quad -- \quad \textbf{Mã HS:} \texttt{HS010}
	\end{minipage}
	\hfill
	\begin{minipage}[c]{0.22\textwidth}
		\begin{tcolorbox}[colback=white,colframe=black,arc=0mm,boxrule=0.6pt,height=1.8cm,halign=center,valign=top,top=2pt]
			{\bfseries\normalsize Điểm}
		\end{tcolorbox}
	\end{minipage}
	\hfill
	\begin{minipage}[c]{0.30\textwidth}
		\raggedleft
		\textit{Thời gian làm bài: 30 phút}\\[4pt]
		\textbf{Mã đề thi:} \fbox{\textbf{\Large\strut\ 001\ }}
	\end{minipage}
\end{tcolorbox}
\vspace{0.25cm}

\noindent\textbf{\large ĐỀ BÀI (MÃ ĐỀ: 001)}
\vspace{0.1cm}

% --- CÂU 1 (TRANG 1) ---
\begin{ex}
	Thực hiện phép tính:
	\begin{enumerate}[a)]
		\item $A = 3x^3y \cdot 4xy^2$
		\item $B = 3x^2y^3 \cdot 2x^2y - 11x^4y^4$
		\item $C = 18x^6y^5 : 3x^2y^2 - 12x^4y^3$
		\item $D = 4x^3y^2 + 8x^5y^3 : 2x^2y$
	\end{enumerate}
	\loigiai{
		\begin{enumerate}[a)]
			\item Ta có:
			\begin{align*}
				A &= 3x^3y \cdot 4xy^2 \\
				&= (3 \cdot 4) \cdot (x^3 \cdot x) \cdot (y \cdot y^2) \\
				&= 12x^4y^3
			\end{align*}
			\item Ta có:
			\begin{align*}
				B &= 3x^2y^3 \cdot 2x^2y - 11x^4y^4 \\
				&= 6x^4y^4 - 11x^4y^4 \\
				&= (6 - 11)x^4y^4 \\
				&= -5x^4y^4
			\end{align*}
			\item Ta có:
			\begin{align*}
				C &= 18x^6y^5 : 3x^2y^2 - 12x^4y^3 \\
				&= 6x^4y^3 - 12x^4y^3 \\
				&= (6 - 12)x^4y^3 \\
				&= -6x^4y^3
			\end{align*}
			\item Ta có:
			\begin{align*}
				D &= 4x^3y^2 + 8x^5y^3 : 2x^2y \\
				&= 4x^3y^2 + 4x^3y^2 \\
				&= 8x^3y^2
			\end{align*}
		\end{enumerate}
	}
\end{ex}

\newpage
% ----------------------------------------------------------------------
% TRANG 2 (MẶT SAU) -- DÒNG KẺ TIẾP THEO
% ----------------------------------------------------------------------
\noindent
\textbf{Học sinh:} Mai Phương \quad -- \quad \textbf{Mã đề:} \fbox{\textbf{001}} \hfill \textit{Trang 2/2}\par
\vspace{-0.1cm}
\noindent\hrulefill
\vspace{0.3cm}

% --- CÂU 2 (TRANG 2) ---
\begin{ex}
	Thực hiện phép tính:
	\begin{enumerate}[a)]
		\item $A = -5a^2b^3 \cdot 2ab^2$
		\item $B = 4a^3b \cdot 3ab^2 - 17a^4b^3$
		\item $C = 24a^7b^4 : 6a^4b^2 - 9a^3b^2$
		\item $D = 3a^2b^4 + 14a^4b^5 : 2a^2b$
	\end{enumerate}
	\loigiai{
		\begin{enumerate}[a)]
			\item Ta có:
			\begin{align*}
				A &= -5a^2b^3 \cdot 2ab^2 \\
				&= (-5 \cdot 2) \cdot (a^2 \cdot a) \cdot (b^3 \cdot b^2) \\
				&= -10a^3b^5
			\end{align*}
			\item Ta có:
			\begin{align*}
				B &= 4a^3b \cdot 3ab^2 - 17a^4b^3 \\
				&= 12a^4b^3 - 17a^4b^3 \\
				&= (12 - 17)a^4b^3 \\
				&= -5a^4b^3
			\end{align*}
			\item Ta có:
			\begin{align*}
				C &= 24a^7b^4 : 6a^4b^2 - 9a^3b^2 \\
				&= 4a^3b^2 - 9a^3b^2 \\
				&= (4 - 9)a^3b^2 \\
				&= -5a^3b^2
			\end{align*}
			\item Ta có:
			\begin{align*}
				D &= 3a^2b^4 + 14a^4b^5 : 2a^2b \\
				&= 3a^2b^4 + 7a^2b^4 \\
				&= 10a^2b^4
			\end{align*}
		\end{enumerate}
	}
\end{ex}

\newpage

% ======================================================================
% HỌC SINH: Minh Trí | MÃ: HS040 | MÃ ĐỀ: 004
% ======================================================================
\noindent
\begin{tcolorbox}[colback=white,colframe=black,arc=0mm,boxrule=0.8pt,top=4pt,bottom=4pt,left=6pt,right=6pt]
	\begin{minipage}[c]{0.44\textwidth}
		{\bfseries\large LỚP TOÁN THẦY THÔNG}\\[3pt]
		\textbf{Họ và tên:} \textbf{\large Minh Trí}\\[3pt]
		\textbf{Lớp:} LỚP 10 \quad -- \quad \textbf{Mã HS:} \texttt{HS040}
	\end{minipage}
	\hfill
	\begin{minipage}[c]{0.22\textwidth}
		\begin{tcolorbox}[colback=white,colframe=black,arc=0mm,boxrule=0.6pt,height=1.8cm,halign=center,valign=top,top=2pt]
			{\bfseries\normalsize Điểm}
		\end{tcolorbox}
	\end{minipage}
	\hfill
	\begin{minipage}[c]{0.30\textwidth}
		\raggedleft
		\textit{Thời gian làm bài: 30 phút}\\[4pt]
		\textbf{Mã đề thi:} \fbox{\textbf{\Large\strut\ 004\ }}
	\end{minipage}
\end{tcolorbox}
\vspace{0.25cm}

\noindent\textbf{\large ĐỀ BÀI (MÃ ĐỀ: 004)}
\vspace{0.1cm}

% --- CÂU 1 (TRANG 1) ---
\begin{ex}
	Thực hiện phép tính:
	\begin{enumerate}[a)]
		\item $A = 3x^3y \cdot 4xy^2$
		\item $B = 3x^2y^3 \cdot 2x^2y - 11x^4y^4$
		\item $C = 18x^6y^5 : 3x^2y^2 - 12x^4y^3$
		\item $D = 4x^3y^2 + 8x^5y^3 : 2x^2y$
	\end{enumerate}
	\loigiai{
		\begin{enumerate}[a)]
			\item Ta có:
			\begin{align*}
				A &= 3x^3y \cdot 4xy^2 \\
				&= (3 \cdot 4) \cdot (x^3 \cdot x) \cdot (y \cdot y^2) \\
				&= 12x^4y^3
			\end{align*}
			\item Ta có:
			\begin{align*}
				B &= 3x^2y^3 \cdot 2x^2y - 11x^4y^4 \\
				&= 6x^4y^4 - 11x^4y^4 \\
				&= (6 - 11)x^4y^4 \\
				&= -5x^4y^4
			\end{align*}
			\item Ta có:
			\begin{align*}
				C &= 18x^6y^5 : 3x^2y^2 - 12x^4y^3 \\
				&= 6x^4y^3 - 12x^4y^3 \\
				&= (6 - 12)x^4y^3 \\
				&= -6x^4y^3
			\end{align*}
			\item Ta có:
			\begin{align*}
				D &= 4x^3y^2 + 8x^5y^3 : 2x^2y \\
				&= 4x^3y^2 + 4x^3y^2 \\
				&= 8x^3y^2
			\end{align*}
		\end{enumerate}
	}
\end{ex}

\newpage
% ----------------------------------------------------------------------
% TRANG 2 (MẶT SAU) -- DÒNG KẺ TIẾP THEO
% ----------------------------------------------------------------------
\noindent
\textbf{Học sinh:} Minh Trí \quad -- \quad \textbf{Mã đề:} \fbox{\textbf{004}} \hfill \textit{Trang 2/2}\par
\vspace{-0.1cm}
\noindent\hrulefill
\vspace{0.3cm}

% --- CÂU 2 (TRANG 2) ---
\begin{ex}
	Thực hiện phép tính:
	\begin{enumerate}[a)]
		\item $A = -5a^2b^3 \cdot 2ab^2$
		\item $B = 4a^3b \cdot 3ab^2 - 17a^4b^3$
		\item $C = 24a^7b^4 : 6a^4b^2 - 9a^3b^2$
		\item $D = 3a^2b^4 + 14a^4b^5 : 2a^2b$
	\end{enumerate}
	\loigiai{
		\begin{enumerate}[a)]
			\item Ta có:
			\begin{align*}
				A &= -5a^2b^3 \cdot 2ab^2 \\
				&= (-5 \cdot 2) \cdot (a^2 \cdot a) \cdot (b^3 \cdot b^2) \\
				&= -10a^3b^5
			\end{align*}
			\item Ta có:
			\begin{align*}
				B &= 4a^3b \cdot 3ab^2 - 17a^4b^3 \\
				&= 12a^4b^3 - 17a^4b^3 \\
				&= (12 - 17)a^4b^3 \\
				&= -5a^4b^3
			\end{align*}
			\item Ta có:
			\begin{align*}
				C &= 24a^7b^4 : 6a^4b^2 - 9a^3b^2 \\
				&= 4a^3b^2 - 9a^3b^2 \\
				&= (4 - 9)a^3b^2 \\
				&= -5a^3b^2
			\end{align*}
			\item Ta có:
			\begin{align*}
				D &= 3a^2b^4 + 14a^4b^5 : 2a^2b \\
				&= 3a^2b^4 + 7a^2b^4 \\
				&= 10a^2b^4
			\end{align*}
		\end{enumerate}
	}
\end{ex}

\newpage

% ======================================================================
% HỌC SINH: Ngân Giang | MÃ: HS016 | MÃ ĐỀ: 003
% ======================================================================
\noindent
\begin{tcolorbox}[colback=white,colframe=black,arc=0mm,boxrule=0.8pt,top=4pt,bottom=4pt,left=6pt,right=6pt]
	\begin{minipage}[c]{0.44\textwidth}
		{\bfseries\large LỚP TOÁN THẦY THÔNG}\\[3pt]
		\textbf{Họ và tên:} \textbf{\large Ngân Giang}\\[3pt]
		\textbf{Lớp:} LỚP 10 \quad -- \quad \textbf{Mã HS:} \texttt{HS016}
	\end{minipage}
	\hfill
	\begin{minipage}[c]{0.22\textwidth}
		\begin{tcolorbox}[colback=white,colframe=black,arc=0mm,boxrule=0.6pt,height=1.8cm,halign=center,valign=top,top=2pt]
			{\bfseries\normalsize Điểm}
		\end{tcolorbox}
	\end{minipage}
	\hfill
	\begin{minipage}[c]{0.30\textwidth}
		\raggedleft
		\textit{Thời gian làm bài: 30 phút}\\[4pt]
		\textbf{Mã đề thi:} \fbox{\textbf{\Large\strut\ 003\ }}
	\end{minipage}
\end{tcolorbox}
\vspace{0.25cm}

\noindent\textbf{\large ĐỀ BÀI (MÃ ĐỀ: 003)}
\vspace{0.1cm}

% --- CÂU 1 (TRANG 1) ---
\begin{ex}
	Thực hiện phép tính:
	\begin{enumerate}[a)]
		\item $A = 3x^3y \cdot 4xy^2$
		\item $B = 3x^2y^3 \cdot 2x^2y - 11x^4y^4$
		\item $C = 18x^6y^5 : 3x^2y^2 - 12x^4y^3$
		\item $D = 4x^3y^2 + 8x^5y^3 : 2x^2y$
	\end{enumerate}
	\loigiai{
		\begin{enumerate}[a)]
			\item Ta có:
			\begin{align*}
				A &= 3x^3y \cdot 4xy^2 \\
				&= (3 \cdot 4) \cdot (x^3 \cdot x) \cdot (y \cdot y^2) \\
				&= 12x^4y^3
			\end{align*}
			\item Ta có:
			\begin{align*}
				B &= 3x^2y^3 \cdot 2x^2y - 11x^4y^4 \\
				&= 6x^4y^4 - 11x^4y^4 \\
				&= (6 - 11)x^4y^4 \\
				&= -5x^4y^4
			\end{align*}
			\item Ta có:
			\begin{align*}
				C &= 18x^6y^5 : 3x^2y^2 - 12x^4y^3 \\
				&= 6x^4y^3 - 12x^4y^3 \\
				&= (6 - 12)x^4y^3 \\
				&= -6x^4y^3
			\end{align*}
			\item Ta có:
			\begin{align*}
				D &= 4x^3y^2 + 8x^5y^3 : 2x^2y \\
				&= 4x^3y^2 + 4x^3y^2 \\
				&= 8x^3y^2
			\end{align*}
		\end{enumerate}
	}
\end{ex}

\newpage
% ----------------------------------------------------------------------
% TRANG 2 (MẶT SAU) -- DÒNG KẺ TIẾP THEO
% ----------------------------------------------------------------------
\noindent
\textbf{Học sinh:} Ngân Giang \quad -- \quad \textbf{Mã đề:} \fbox{\textbf{003}} \hfill \textit{Trang 2/2}\par
\vspace{-0.1cm}
\noindent\hrulefill
\vspace{0.3cm}

% --- CÂU 2 (TRANG 2) ---
\begin{ex}
	Thực hiện phép tính:
	\begin{enumerate}[a)]
		\item $A = -5a^2b^3 \cdot 2ab^2$
		\item $B = 4a^3b \cdot 3ab^2 - 17a^4b^3$
		\item $C = 24a^7b^4 : 6a^4b^2 - 9a^3b^2$
		\item $D = 3a^2b^4 + 14a^4b^5 : 2a^2b$
	\end{enumerate}
	\loigiai{
		\begin{enumerate}[a)]
			\item Ta có:
			\begin{align*}
				A &= -5a^2b^3 \cdot 2ab^2 \\
				&= (-5 \cdot 2) \cdot (a^2 \cdot a) \cdot (b^3 \cdot b^2) \\
				&= -10a^3b^5
			\end{align*}
			\item Ta có:
			\begin{align*}
				B &= 4a^3b \cdot 3ab^2 - 17a^4b^3 \\
				&= 12a^4b^3 - 17a^4b^3 \\
				&= (12 - 17)a^4b^3 \\
				&= -5a^4b^3
			\end{align*}
			\item Ta có:
			\begin{align*}
				C &= 24a^7b^4 : 6a^4b^2 - 9a^3b^2 \\
				&= 4a^3b^2 - 9a^3b^2 \\
				&= (4 - 9)a^3b^2 \\
				&= -5a^3b^2
			\end{align*}
			\item Ta có:
			\begin{align*}
				D &= 3a^2b^4 + 14a^4b^5 : 2a^2b \\
				&= 3a^2b^4 + 7a^2b^4 \\
				&= 10a^2b^4
			\end{align*}
		\end{enumerate}
	}
\end{ex}

\newpage

% ======================================================================
% HỌC SINH: Ngọc Diệp | MÃ: HS101 | MÃ ĐỀ: 002
% ======================================================================
\noindent
\begin{tcolorbox}[colback=white,colframe=black,arc=0mm,boxrule=0.8pt,top=4pt,bottom=4pt,left=6pt,right=6pt]
	\begin{minipage}[c]{0.44\textwidth}
		{\bfseries\large LỚP TOÁN THẦY THÔNG}\\[3pt]
		\textbf{Họ và tên:} \textbf{\large Ngọc Diệp}\\[3pt]
		\textbf{Lớp:} LỚP 10 \quad -- \quad \textbf{Mã HS:} \texttt{HS101}
	\end{minipage}
	\hfill
	\begin{minipage}[c]{0.22\textwidth}
		\begin{tcolorbox}[colback=white,colframe=black,arc=0mm,boxrule=0.6pt,height=1.8cm,halign=center,valign=top,top=2pt]
			{\bfseries\normalsize Điểm}
		\end{tcolorbox}
	\end{minipage}
	\hfill
	\begin{minipage}[c]{0.30\textwidth}
		\raggedleft
		\textit{Thời gian làm bài: 30 phút}\\[4pt]
		\textbf{Mã đề thi:} \fbox{\textbf{\Large\strut\ 002\ }}
	\end{minipage}
\end{tcolorbox}
\vspace{0.25cm}

\noindent\textbf{\large ĐỀ BÀI (MÃ ĐỀ: 002)}
\vspace{0.1cm}

% --- CÂU 1 (TRANG 1) ---
\begin{ex}
	Thực hiện phép tính:
	\begin{enumerate}[a)]
		\item $A = 3x^3y \cdot 4xy^2$
		\item $B = 3x^2y^3 \cdot 2x^2y - 11x^4y^4$
		\item $C = 18x^6y^5 : 3x^2y^2 - 12x^4y^3$
		\item $D = 4x^3y^2 + 8x^5y^3 : 2x^2y$
	\end{enumerate}
	\loigiai{
		\begin{enumerate}[a)]
			\item Ta có:
			\begin{align*}
				A &= 3x^3y \cdot 4xy^2 \\
				&= (3 \cdot 4) \cdot (x^3 \cdot x) \cdot (y \cdot y^2) \\
				&= 12x^4y^3
			\end{align*}
			\item Ta có:
			\begin{align*}
				B &= 3x^2y^3 \cdot 2x^2y - 11x^4y^4 \\
				&= 6x^4y^4 - 11x^4y^4 \\
				&= (6 - 11)x^4y^4 \\
				&= -5x^4y^4
			\end{align*}
			\item Ta có:
			\begin{align*}
				C &= 18x^6y^5 : 3x^2y^2 - 12x^4y^3 \\
				&= 6x^4y^3 - 12x^4y^3 \\
				&= (6 - 12)x^4y^3 \\
				&= -6x^4y^3
			\end{align*}
			\item Ta có:
			\begin{align*}
				D &= 4x^3y^2 + 8x^5y^3 : 2x^2y \\
				&= 4x^3y^2 + 4x^3y^2 \\
				&= 8x^3y^2
			\end{align*}
		\end{enumerate}
	}
\end{ex}

\newpage
% ----------------------------------------------------------------------
% TRANG 2 (MẶT SAU) -- DÒNG KẺ TIẾP THEO
% ----------------------------------------------------------------------
\noindent
\textbf{Học sinh:} Ngọc Diệp \quad -- \quad \textbf{Mã đề:} \fbox{\textbf{002}} \hfill \textit{Trang 2/2}\par
\vspace{-0.1cm}
\noindent\hrulefill
\vspace{0.3cm}

% --- CÂU 2 (TRANG 2) ---
\begin{ex}
	Thực hiện phép tính:
	\begin{enumerate}[a)]
		\item $A = -5a^2b^3 \cdot 2ab^2$
		\item $B = 4a^3b \cdot 3ab^2 - 17a^4b^3$
		\item $C = 24a^7b^4 : 6a^4b^2 - 9a^3b^2$
		\item $D = 3a^2b^4 + 14a^4b^5 : 2a^2b$
	\end{enumerate}
	\loigiai{
		\begin{enumerate}[a)]
			\item Ta có:
			\begin{align*}
				A &= -5a^2b^3 \cdot 2ab^2 \\
				&= (-5 \cdot 2) \cdot (a^2 \cdot a) \cdot (b^3 \cdot b^2) \\
				&= -10a^3b^5
			\end{align*}
			\item Ta có:
			\begin{align*}
				B &= 4a^3b \cdot 3ab^2 - 17a^4b^3 \\
				&= 12a^4b^3 - 17a^4b^3 \\
				&= (12 - 17)a^4b^3 \\
				&= -5a^4b^3
			\end{align*}
			\item Ta có:
			\begin{align*}
				C &= 24a^7b^4 : 6a^4b^2 - 9a^3b^2 \\
				&= 4a^3b^2 - 9a^3b^2 \\
				&= (4 - 9)a^3b^2 \\
				&= -5a^3b^2
			\end{align*}
			\item Ta có:
			\begin{align*}
				D &= 3a^2b^4 + 14a^4b^5 : 2a^2b \\
				&= 3a^2b^4 + 7a^2b^4 \\
				&= 10a^2b^4
			\end{align*}
		\end{enumerate}
	}
\end{ex}

\newpage

% ======================================================================
% HỌC SINH: Ngọc Linh | MÃ: HS039 | MÃ ĐỀ: 006
% ======================================================================
\noindent
\begin{tcolorbox}[colback=white,colframe=black,arc=0mm,boxrule=0.8pt,top=4pt,bottom=4pt,left=6pt,right=6pt]
	\begin{minipage}[c]{0.44\textwidth}
		{\bfseries\large LỚP TOÁN THẦY THÔNG}\\[3pt]
		\textbf{Họ và tên:} \textbf{\large Ngọc Linh}\\[3pt]
		\textbf{Lớp:} LỚP 10 \quad -- \quad \textbf{Mã HS:} \texttt{HS039}
	\end{minipage}
	\hfill
	\begin{minipage}[c]{0.22\textwidth}
		\begin{tcolorbox}[colback=white,colframe=black,arc=0mm,boxrule=0.6pt,height=1.8cm,halign=center,valign=top,top=2pt]
			{\bfseries\normalsize Điểm}
		\end{tcolorbox}
	\end{minipage}
	\hfill
	\begin{minipage}[c]{0.30\textwidth}
		\raggedleft
		\textit{Thời gian làm bài: 30 phút}\\[4pt]
		\textbf{Mã đề thi:} \fbox{\textbf{\Large\strut\ 006\ }}
	\end{minipage}
\end{tcolorbox}
\vspace{0.25cm}

\noindent\textbf{\large ĐỀ BÀI (MÃ ĐỀ: 006)}
\vspace{0.1cm}

% --- CÂU 1 (TRANG 1) ---
\begin{ex}
	Thực hiện phép tính:
	\begin{enumerate}[a)]
		\item $A = 3x^3y \cdot 4xy^2$
		\item $B = 3x^2y^3 \cdot 2x^2y - 11x^4y^4$
		\item $C = 18x^6y^5 : 3x^2y^2 - 12x^4y^3$
		\item $D = 4x^3y^2 + 8x^5y^3 : 2x^2y$
	\end{enumerate}
	\loigiai{
		\begin{enumerate}[a)]
			\item Ta có:
			\begin{align*}
				A &= 3x^3y \cdot 4xy^2 \\
				&= (3 \cdot 4) \cdot (x^3 \cdot x) \cdot (y \cdot y^2) \\
				&= 12x^4y^3
			\end{align*}
			\item Ta có:
			\begin{align*}
				B &= 3x^2y^3 \cdot 2x^2y - 11x^4y^4 \\
				&= 6x^4y^4 - 11x^4y^4 \\
				&= (6 - 11)x^4y^4 \\
				&= -5x^4y^4
			\end{align*}
			\item Ta có:
			\begin{align*}
				C &= 18x^6y^5 : 3x^2y^2 - 12x^4y^3 \\
				&= 6x^4y^3 - 12x^4y^3 \\
				&= (6 - 12)x^4y^3 \\
				&= -6x^4y^3
			\end{align*}
			\item Ta có:
			\begin{align*}
				D &= 4x^3y^2 + 8x^5y^3 : 2x^2y \\
				&= 4x^3y^2 + 4x^3y^2 \\
				&= 8x^3y^2
			\end{align*}
		\end{enumerate}
	}
\end{ex}

\newpage
% ----------------------------------------------------------------------
% TRANG 2 (MẶT SAU) -- DÒNG KẺ TIẾP THEO
% ----------------------------------------------------------------------
\noindent
\textbf{Học sinh:} Ngọc Linh \quad -- \quad \textbf{Mã đề:} \fbox{\textbf{006}} \hfill \textit{Trang 2/2}\par
\vspace{-0.1cm}
\noindent\hrulefill
\vspace{0.3cm}

% --- CÂU 2 (TRANG 2) ---
\begin{ex}
	Thực hiện phép tính:
	\begin{enumerate}[a)]
		\item $A = -5a^2b^3 \cdot 2ab^2$
		\item $B = 4a^3b \cdot 3ab^2 - 17a^4b^3$
		\item $C = 24a^7b^4 : 6a^4b^2 - 9a^3b^2$
		\item $D = 3a^2b^4 + 14a^4b^5 : 2a^2b$
	\end{enumerate}
	\loigiai{
		\begin{enumerate}[a)]
			\item Ta có:
			\begin{align*}
				A &= -5a^2b^3 \cdot 2ab^2 \\
				&= (-5 \cdot 2) \cdot (a^2 \cdot a) \cdot (b^3 \cdot b^2) \\
				&= -10a^3b^5
			\end{align*}
			\item Ta có:
			\begin{align*}
				B &= 4a^3b \cdot 3ab^2 - 17a^4b^3 \\
				&= 12a^4b^3 - 17a^4b^3 \\
				&= (12 - 17)a^4b^3 \\
				&= -5a^4b^3
			\end{align*}
			\item Ta có:
			\begin{align*}
				C &= 24a^7b^4 : 6a^4b^2 - 9a^3b^2 \\
				&= 4a^3b^2 - 9a^3b^2 \\
				&= (4 - 9)a^3b^2 \\
				&= -5a^3b^2
			\end{align*}
			\item Ta có:
			\begin{align*}
				D &= 3a^2b^4 + 14a^4b^5 : 2a^2b \\
				&= 3a^2b^4 + 7a^2b^4 \\
				&= 10a^2b^4
			\end{align*}
		\end{enumerate}
	}
\end{ex}

\newpage

% ======================================================================
% HỌC SINH: Ngọc Sương | MÃ: HS029 | MÃ ĐỀ: 008
% ======================================================================
\noindent
\begin{tcolorbox}[colback=white,colframe=black,arc=0mm,boxrule=0.8pt,top=4pt,bottom=4pt,left=6pt,right=6pt]
	\begin{minipage}[c]{0.44\textwidth}
		{\bfseries\large LỚP TOÁN THẦY THÔNG}\\[3pt]
		\textbf{Họ và tên:} \textbf{\large Ngọc Sương}\\[3pt]
		\textbf{Lớp:} LỚP 10 \quad -- \quad \textbf{Mã HS:} \texttt{HS029}
	\end{minipage}
	\hfill
	\begin{minipage}[c]{0.22\textwidth}
		\begin{tcolorbox}[colback=white,colframe=black,arc=0mm,boxrule=0.6pt,height=1.8cm,halign=center,valign=top,top=2pt]
			{\bfseries\normalsize Điểm}
		\end{tcolorbox}
	\end{minipage}
	\hfill
	\begin{minipage}[c]{0.30\textwidth}
		\raggedleft
		\textit{Thời gian làm bài: 30 phút}\\[4pt]
		\textbf{Mã đề thi:} \fbox{\textbf{\Large\strut\ 008\ }}
	\end{minipage}
\end{tcolorbox}
\vspace{0.25cm}

\noindent\textbf{\large ĐỀ BÀI (MÃ ĐỀ: 008)}
\vspace{0.1cm}

% --- CÂU 1 (TRANG 1) ---
\begin{ex}
	Thực hiện phép tính:
	\begin{enumerate}[a)]
		\item $A = 3x^3y \cdot 4xy^2$
		\item $B = 3x^2y^3 \cdot 2x^2y - 11x^4y^4$
		\item $C = 18x^6y^5 : 3x^2y^2 - 12x^4y^3$
		\item $D = 4x^3y^2 + 8x^5y^3 : 2x^2y$
	\end{enumerate}
	\loigiai{
		\begin{enumerate}[a)]
			\item Ta có:
			\begin{align*}
				A &= 3x^3y \cdot 4xy^2 \\
				&= (3 \cdot 4) \cdot (x^3 \cdot x) \cdot (y \cdot y^2) \\
				&= 12x^4y^3
			\end{align*}
			\item Ta có:
			\begin{align*}
				B &= 3x^2y^3 \cdot 2x^2y - 11x^4y^4 \\
				&= 6x^4y^4 - 11x^4y^4 \\
				&= (6 - 11)x^4y^4 \\
				&= -5x^4y^4
			\end{align*}
			\item Ta có:
			\begin{align*}
				C &= 18x^6y^5 : 3x^2y^2 - 12x^4y^3 \\
				&= 6x^4y^3 - 12x^4y^3 \\
				&= (6 - 12)x^4y^3 \\
				&= -6x^4y^3
			\end{align*}
			\item Ta có:
			\begin{align*}
				D &= 4x^3y^2 + 8x^5y^3 : 2x^2y \\
				&= 4x^3y^2 + 4x^3y^2 \\
				&= 8x^3y^2
			\end{align*}
		\end{enumerate}
	}
\end{ex}

\newpage
% ----------------------------------------------------------------------
% TRANG 2 (MẶT SAU) -- DÒNG KẺ TIẾP THEO
% ----------------------------------------------------------------------
\noindent
\textbf{Học sinh:} Ngọc Sương \quad -- \quad \textbf{Mã đề:} \fbox{\textbf{008}} \hfill \textit{Trang 2/2}\par
\vspace{-0.1cm}
\noindent\hrulefill
\vspace{0.3cm}

% --- CÂU 2 (TRANG 2) ---
\begin{ex}
	Thực hiện phép tính:
	\begin{enumerate}[a)]
		\item $A = -5a^2b^3 \cdot 2ab^2$
		\item $B = 4a^3b \cdot 3ab^2 - 17a^4b^3$
		\item $C = 24a^7b^4 : 6a^4b^2 - 9a^3b^2$
		\item $D = 3a^2b^4 + 14a^4b^5 : 2a^2b$
	\end{enumerate}
	\loigiai{
		\begin{enumerate}[a)]
			\item Ta có:
			\begin{align*}
				A &= -5a^2b^3 \cdot 2ab^2 \\
				&= (-5 \cdot 2) \cdot (a^2 \cdot a) \cdot (b^3 \cdot b^2) \\
				&= -10a^3b^5
			\end{align*}
			\item Ta có:
			\begin{align*}
				B &= 4a^3b \cdot 3ab^2 - 17a^4b^3 \\
				&= 12a^4b^3 - 17a^4b^3 \\
				&= (12 - 17)a^4b^3 \\
				&= -5a^4b^3
			\end{align*}
			\item Ta có:
			\begin{align*}
				C &= 24a^7b^4 : 6a^4b^2 - 9a^3b^2 \\
				&= 4a^3b^2 - 9a^3b^2 \\
				&= (4 - 9)a^3b^2 \\
				&= -5a^3b^2
			\end{align*}
			\item Ta có:
			\begin{align*}
				D &= 3a^2b^4 + 14a^4b^5 : 2a^2b \\
				&= 3a^2b^4 + 7a^2b^4 \\
				&= 10a^2b^4
			\end{align*}
		\end{enumerate}
	}
\end{ex}

\newpage

% ======================================================================
% HỌC SINH: Nguyễn Thị Ngọc Ánh | MÃ: HS031 | MÃ ĐỀ: 007
% ======================================================================
\noindent
\begin{tcolorbox}[colback=white,colframe=black,arc=0mm,boxrule=0.8pt,top=4pt,bottom=4pt,left=6pt,right=6pt]
	\begin{minipage}[c]{0.44\textwidth}
		{\bfseries\large LỚP TOÁN THẦY THÔNG}\\[3pt]
		\textbf{Họ và tên:} \textbf{\large Nguyễn Thị Ngọc Ánh}\\[3pt]
		\textbf{Lớp:} LỚP 10 \quad -- \quad \textbf{Mã HS:} \texttt{HS031}
	\end{minipage}
	\hfill
	\begin{minipage}[c]{0.22\textwidth}
		\begin{tcolorbox}[colback=white,colframe=black,arc=0mm,boxrule=0.6pt,height=1.8cm,halign=center,valign=top,top=2pt]
			{\bfseries\normalsize Điểm}
		\end{tcolorbox}
	\end{minipage}
	\hfill
	\begin{minipage}[c]{0.30\textwidth}
		\raggedleft
		\textit{Thời gian làm bài: 30 phút}\\[4pt]
		\textbf{Mã đề thi:} \fbox{\textbf{\Large\strut\ 007\ }}
	\end{minipage}
\end{tcolorbox}
\vspace{0.25cm}

\noindent\textbf{\large ĐỀ BÀI (MÃ ĐỀ: 007)}
\vspace{0.1cm}

% --- CÂU 1 (TRANG 1) ---
\begin{ex}
	Thực hiện phép tính:
	\begin{enumerate}[a)]
		\item $A = 3x^3y \cdot 4xy^2$
		\item $B = 3x^2y^3 \cdot 2x^2y - 11x^4y^4$
		\item $C = 18x^6y^5 : 3x^2y^2 - 12x^4y^3$
		\item $D = 4x^3y^2 + 8x^5y^3 : 2x^2y$
	\end{enumerate}
	\loigiai{
		\begin{enumerate}[a)]
			\item Ta có:
			\begin{align*}
				A &= 3x^3y \cdot 4xy^2 \\
				&= (3 \cdot 4) \cdot (x^3 \cdot x) \cdot (y \cdot y^2) \\
				&= 12x^4y^3
			\end{align*}
			\item Ta có:
			\begin{align*}
				B &= 3x^2y^3 \cdot 2x^2y - 11x^4y^4 \\
				&= 6x^4y^4 - 11x^4y^4 \\
				&= (6 - 11)x^4y^4 \\
				&= -5x^4y^4
			\end{align*}
			\item Ta có:
			\begin{align*}
				C &= 18x^6y^5 : 3x^2y^2 - 12x^4y^3 \\
				&= 6x^4y^3 - 12x^4y^3 \\
				&= (6 - 12)x^4y^3 \\
				&= -6x^4y^3
			\end{align*}
			\item Ta có:
			\begin{align*}
				D &= 4x^3y^2 + 8x^5y^3 : 2x^2y \\
				&= 4x^3y^2 + 4x^3y^2 \\
				&= 8x^3y^2
			\end{align*}
		\end{enumerate}
	}
\end{ex}

\newpage
% ----------------------------------------------------------------------
% TRANG 2 (MẶT SAU) -- DÒNG KẺ TIẾP THEO
% ----------------------------------------------------------------------
\noindent
\textbf{Học sinh:} Nguyễn Thị Ngọc Ánh \quad -- \quad \textbf{Mã đề:} \fbox{\textbf{007}} \hfill \textit{Trang 2/2}\par
\vspace{-0.1cm}
\noindent\hrulefill
\vspace{0.3cm}

% --- CÂU 2 (TRANG 2) ---
\begin{ex}
	Thực hiện phép tính:
	\begin{enumerate}[a)]
		\item $A = -5a^2b^3 \cdot 2ab^2$
		\item $B = 4a^3b \cdot 3ab^2 - 17a^4b^3$
		\item $C = 24a^7b^4 : 6a^4b^2 - 9a^3b^2$
		\item $D = 3a^2b^4 + 14a^4b^5 : 2a^2b$
	\end{enumerate}
	\loigiai{
		\begin{enumerate}[a)]
			\item Ta có:
			\begin{align*}
				A &= -5a^2b^3 \cdot 2ab^2 \\
				&= (-5 \cdot 2) \cdot (a^2 \cdot a) \cdot (b^3 \cdot b^2) \\
				&= -10a^3b^5
			\end{align*}
			\item Ta có:
			\begin{align*}
				B &= 4a^3b \cdot 3ab^2 - 17a^4b^3 \\
				&= 12a^4b^3 - 17a^4b^3 \\
				&= (12 - 17)a^4b^3 \\
				&= -5a^4b^3
			\end{align*}
			\item Ta có:
			\begin{align*}
				C &= 24a^7b^4 : 6a^4b^2 - 9a^3b^2 \\
				&= 4a^3b^2 - 9a^3b^2 \\
				&= (4 - 9)a^3b^2 \\
				&= -5a^3b^2
			\end{align*}
			\item Ta có:
			\begin{align*}
				D &= 3a^2b^4 + 14a^4b^5 : 2a^2b \\
				&= 3a^2b^4 + 7a^2b^4 \\
				&= 10a^2b^4
			\end{align*}
		\end{enumerate}
	}
\end{ex}

\newpage

% ======================================================================
% HỌC SINH: Nguyệt Linh | MÃ: HS017 | MÃ ĐỀ: 003
% ======================================================================
\noindent
\begin{tcolorbox}[colback=white,colframe=black,arc=0mm,boxrule=0.8pt,top=4pt,bottom=4pt,left=6pt,right=6pt]
	\begin{minipage}[c]{0.44\textwidth}
		{\bfseries\large LỚP TOÁN THẦY THÔNG}\\[3pt]
		\textbf{Họ và tên:} \textbf{\large Nguyệt Linh}\\[3pt]
		\textbf{Lớp:} LỚP 10 \quad -- \quad \textbf{Mã HS:} \texttt{HS017}
	\end{minipage}
	\hfill
	\begin{minipage}[c]{0.22\textwidth}
		\begin{tcolorbox}[colback=white,colframe=black,arc=0mm,boxrule=0.6pt,height=1.8cm,halign=center,valign=top,top=2pt]
			{\bfseries\normalsize Điểm}
		\end{tcolorbox}
	\end{minipage}
	\hfill
	\begin{minipage}[c]{0.30\textwidth}
		\raggedleft
		\textit{Thời gian làm bài: 30 phút}\\[4pt]
		\textbf{Mã đề thi:} \fbox{\textbf{\Large\strut\ 003\ }}
	\end{minipage}
\end{tcolorbox}
\vspace{0.25cm}

\noindent\textbf{\large ĐỀ BÀI (MÃ ĐỀ: 003)}
\vspace{0.1cm}

% --- CÂU 1 (TRANG 1) ---
\begin{ex}
	Thực hiện phép tính:
	\begin{enumerate}[a)]
		\item $A = 3x^3y \cdot 4xy^2$
		\item $B = 3x^2y^3 \cdot 2x^2y - 11x^4y^4$
		\item $C = 18x^6y^5 : 3x^2y^2 - 12x^4y^3$
		\item $D = 4x^3y^2 + 8x^5y^3 : 2x^2y$
	\end{enumerate}
	\loigiai{
		\begin{enumerate}[a)]
			\item Ta có:
			\begin{align*}
				A &= 3x^3y \cdot 4xy^2 \\
				&= (3 \cdot 4) \cdot (x^3 \cdot x) \cdot (y \cdot y^2) \\
				&= 12x^4y^3
			\end{align*}
			\item Ta có:
			\begin{align*}
				B &= 3x^2y^3 \cdot 2x^2y - 11x^4y^4 \\
				&= 6x^4y^4 - 11x^4y^4 \\
				&= (6 - 11)x^4y^4 \\
				&= -5x^4y^4
			\end{align*}
			\item Ta có:
			\begin{align*}
				C &= 18x^6y^5 : 3x^2y^2 - 12x^4y^3 \\
				&= 6x^4y^3 - 12x^4y^3 \\
				&= (6 - 12)x^4y^3 \\
				&= -6x^4y^3
			\end{align*}
			\item Ta có:
			\begin{align*}
				D &= 4x^3y^2 + 8x^5y^3 : 2x^2y \\
				&= 4x^3y^2 + 4x^3y^2 \\
				&= 8x^3y^2
			\end{align*}
		\end{enumerate}
	}
\end{ex}

\newpage
% ----------------------------------------------------------------------
% TRANG 2 (MẶT SAU) -- DÒNG KẺ TIẾP THEO
% ----------------------------------------------------------------------
\noindent
\textbf{Học sinh:} Nguyệt Linh \quad -- \quad \textbf{Mã đề:} \fbox{\textbf{003}} \hfill \textit{Trang 2/2}\par
\vspace{-0.1cm}
\noindent\hrulefill
\vspace{0.3cm}

% --- CÂU 2 (TRANG 2) ---
\begin{ex}
	Thực hiện phép tính:
	\begin{enumerate}[a)]
		\item $A = -5a^2b^3 \cdot 2ab^2$
		\item $B = 4a^3b \cdot 3ab^2 - 17a^4b^3$
		\item $C = 24a^7b^4 : 6a^4b^2 - 9a^3b^2$
		\item $D = 3a^2b^4 + 14a^4b^5 : 2a^2b$
	\end{enumerate}
	\loigiai{
		\begin{enumerate}[a)]
			\item Ta có:
			\begin{align*}
				A &= -5a^2b^3 \cdot 2ab^2 \\
				&= (-5 \cdot 2) \cdot (a^2 \cdot a) \cdot (b^3 \cdot b^2) \\
				&= -10a^3b^5
			\end{align*}
			\item Ta có:
			\begin{align*}
				B &= 4a^3b \cdot 3ab^2 - 17a^4b^3 \\
				&= 12a^4b^3 - 17a^4b^3 \\
				&= (12 - 17)a^4b^3 \\
				&= -5a^4b^3
			\end{align*}
			\item Ta có:
			\begin{align*}
				C &= 24a^7b^4 : 6a^4b^2 - 9a^3b^2 \\
				&= 4a^3b^2 - 9a^3b^2 \\
				&= (4 - 9)a^3b^2 \\
				&= -5a^3b^2
			\end{align*}
			\item Ta có:
			\begin{align*}
				D &= 3a^2b^4 + 14a^4b^5 : 2a^2b \\
				&= 3a^2b^4 + 7a^2b^4 \\
				&= 10a^2b^4
			\end{align*}
		\end{enumerate}
	}
\end{ex}

\newpage

% ======================================================================
% HỌC SINH: Nhật Hà On | MÃ: HS044 | MÃ ĐỀ: 004
% ======================================================================
\noindent
\begin{tcolorbox}[colback=white,colframe=black,arc=0mm,boxrule=0.8pt,top=4pt,bottom=4pt,left=6pt,right=6pt]
	\begin{minipage}[c]{0.44\textwidth}
		{\bfseries\large LỚP TOÁN THẦY THÔNG}\\[3pt]
		\textbf{Họ và tên:} \textbf{\large Nhật Hà On}\\[3pt]
		\textbf{Lớp:} LỚP 10 \quad -- \quad \textbf{Mã HS:} \texttt{HS044}
	\end{minipage}
	\hfill
	\begin{minipage}[c]{0.22\textwidth}
		\begin{tcolorbox}[colback=white,colframe=black,arc=0mm,boxrule=0.6pt,height=1.8cm,halign=center,valign=top,top=2pt]
			{\bfseries\normalsize Điểm}
		\end{tcolorbox}
	\end{minipage}
	\hfill
	\begin{minipage}[c]{0.30\textwidth}
		\raggedleft
		\textit{Thời gian làm bài: 30 phút}\\[4pt]
		\textbf{Mã đề thi:} \fbox{\textbf{\Large\strut\ 004\ }}
	\end{minipage}
\end{tcolorbox}
\vspace{0.25cm}

\noindent\textbf{\large ĐỀ BÀI (MÃ ĐỀ: 004)}
\vspace{0.1cm}

% --- CÂU 1 (TRANG 1) ---
\begin{ex}
	Thực hiện phép tính:
	\begin{enumerate}[a)]
		\item $A = 3x^3y \cdot 4xy^2$
		\item $B = 3x^2y^3 \cdot 2x^2y - 11x^4y^4$
		\item $C = 18x^6y^5 : 3x^2y^2 - 12x^4y^3$
		\item $D = 4x^3y^2 + 8x^5y^3 : 2x^2y$
	\end{enumerate}
	\loigiai{
		\begin{enumerate}[a)]
			\item Ta có:
			\begin{align*}
				A &= 3x^3y \cdot 4xy^2 \\
				&= (3 \cdot 4) \cdot (x^3 \cdot x) \cdot (y \cdot y^2) \\
				&= 12x^4y^3
			\end{align*}
			\item Ta có:
			\begin{align*}
				B &= 3x^2y^3 \cdot 2x^2y - 11x^4y^4 \\
				&= 6x^4y^4 - 11x^4y^4 \\
				&= (6 - 11)x^4y^4 \\
				&= -5x^4y^4
			\end{align*}
			\item Ta có:
			\begin{align*}
				C &= 18x^6y^5 : 3x^2y^2 - 12x^4y^3 \\
				&= 6x^4y^3 - 12x^4y^3 \\
				&= (6 - 12)x^4y^3 \\
				&= -6x^4y^3
			\end{align*}
			\item Ta có:
			\begin{align*}
				D &= 4x^3y^2 + 8x^5y^3 : 2x^2y \\
				&= 4x^3y^2 + 4x^3y^2 \\
				&= 8x^3y^2
			\end{align*}
		\end{enumerate}
	}
\end{ex}

\newpage
% ----------------------------------------------------------------------
% TRANG 2 (MẶT SAU) -- DÒNG KẺ TIẾP THEO
% ----------------------------------------------------------------------
\noindent
\textbf{Học sinh:} Nhật Hà On \quad -- \quad \textbf{Mã đề:} \fbox{\textbf{004}} \hfill \textit{Trang 2/2}\par
\vspace{-0.1cm}
\noindent\hrulefill
\vspace{0.3cm}

% --- CÂU 2 (TRANG 2) ---
\begin{ex}
	Thực hiện phép tính:
	\begin{enumerate}[a)]
		\item $A = -5a^2b^3 \cdot 2ab^2$
		\item $B = 4a^3b \cdot 3ab^2 - 17a^4b^3$
		\item $C = 24a^7b^4 : 6a^4b^2 - 9a^3b^2$
		\item $D = 3a^2b^4 + 14a^4b^5 : 2a^2b$
	\end{enumerate}
	\loigiai{
		\begin{enumerate}[a)]
			\item Ta có:
			\begin{align*}
				A &= -5a^2b^3 \cdot 2ab^2 \\
				&= (-5 \cdot 2) \cdot (a^2 \cdot a) \cdot (b^3 \cdot b^2) \\
				&= -10a^3b^5
			\end{align*}
			\item Ta có:
			\begin{align*}
				B &= 4a^3b \cdot 3ab^2 - 17a^4b^3 \\
				&= 12a^4b^3 - 17a^4b^3 \\
				&= (12 - 17)a^4b^3 \\
				&= -5a^4b^3
			\end{align*}
			\item Ta có:
			\begin{align*}
				C &= 24a^7b^4 : 6a^4b^2 - 9a^3b^2 \\
				&= 4a^3b^2 - 9a^3b^2 \\
				&= (4 - 9)a^3b^2 \\
				&= -5a^3b^2
			\end{align*}
			\item Ta có:
			\begin{align*}
				D &= 3a^2b^4 + 14a^4b^5 : 2a^2b \\
				&= 3a^2b^4 + 7a^2b^4 \\
				&= 10a^2b^4
			\end{align*}
		\end{enumerate}
	}
\end{ex}

\newpage

% ======================================================================
% HỌC SINH: Nhật Tuấn On | MÃ: HS045 | MÃ ĐỀ: 002
% ======================================================================
\noindent
\begin{tcolorbox}[colback=white,colframe=black,arc=0mm,boxrule=0.8pt,top=4pt,bottom=4pt,left=6pt,right=6pt]
	\begin{minipage}[c]{0.44\textwidth}
		{\bfseries\large LỚP TOÁN THẦY THÔNG}\\[3pt]
		\textbf{Họ và tên:} \textbf{\large Nhật Tuấn On}\\[3pt]
		\textbf{Lớp:} LỚP 10 \quad -- \quad \textbf{Mã HS:} \texttt{HS045}
	\end{minipage}
	\hfill
	\begin{minipage}[c]{0.22\textwidth}
		\begin{tcolorbox}[colback=white,colframe=black,arc=0mm,boxrule=0.6pt,height=1.8cm,halign=center,valign=top,top=2pt]
			{\bfseries\normalsize Điểm}
		\end{tcolorbox}
	\end{minipage}
	\hfill
	\begin{minipage}[c]{0.30\textwidth}
		\raggedleft
		\textit{Thời gian làm bài: 30 phút}\\[4pt]
		\textbf{Mã đề thi:} \fbox{\textbf{\Large\strut\ 002\ }}
	\end{minipage}
\end{tcolorbox}
\vspace{0.25cm}

\noindent\textbf{\large ĐỀ BÀI (MÃ ĐỀ: 002)}
\vspace{0.1cm}

% --- CÂU 1 (TRANG 1) ---
\begin{ex}
	Thực hiện phép tính:
	\begin{enumerate}[a)]
		\item $A = 3x^3y \cdot 4xy^2$
		\item $B = 3x^2y^3 \cdot 2x^2y - 11x^4y^4$
		\item $C = 18x^6y^5 : 3x^2y^2 - 12x^4y^3$
		\item $D = 4x^3y^2 + 8x^5y^3 : 2x^2y$
	\end{enumerate}
	\loigiai{
		\begin{enumerate}[a)]
			\item Ta có:
			\begin{align*}
				A &= 3x^3y \cdot 4xy^2 \\
				&= (3 \cdot 4) \cdot (x^3 \cdot x) \cdot (y \cdot y^2) \\
				&= 12x^4y^3
			\end{align*}
			\item Ta có:
			\begin{align*}
				B &= 3x^2y^3 \cdot 2x^2y - 11x^4y^4 \\
				&= 6x^4y^4 - 11x^4y^4 \\
				&= (6 - 11)x^4y^4 \\
				&= -5x^4y^4
			\end{align*}
			\item Ta có:
			\begin{align*}
				C &= 18x^6y^5 : 3x^2y^2 - 12x^4y^3 \\
				&= 6x^4y^3 - 12x^4y^3 \\
				&= (6 - 12)x^4y^3 \\
				&= -6x^4y^3
			\end{align*}
			\item Ta có:
			\begin{align*}
				D &= 4x^3y^2 + 8x^5y^3 : 2x^2y \\
				&= 4x^3y^2 + 4x^3y^2 \\
				&= 8x^3y^2
			\end{align*}
		\end{enumerate}
	}
\end{ex}

\newpage
% ----------------------------------------------------------------------
% TRANG 2 (MẶT SAU) -- DÒNG KẺ TIẾP THEO
% ----------------------------------------------------------------------
\noindent
\textbf{Học sinh:} Nhật Tuấn On \quad -- \quad \textbf{Mã đề:} \fbox{\textbf{002}} \hfill \textit{Trang 2/2}\par
\vspace{-0.1cm}
\noindent\hrulefill
\vspace{0.3cm}

% --- CÂU 2 (TRANG 2) ---
\begin{ex}
	Thực hiện phép tính:
	\begin{enumerate}[a)]
		\item $A = -5a^2b^3 \cdot 2ab^2$
		\item $B = 4a^3b \cdot 3ab^2 - 17a^4b^3$
		\item $C = 24a^7b^4 : 6a^4b^2 - 9a^3b^2$
		\item $D = 3a^2b^4 + 14a^4b^5 : 2a^2b$
	\end{enumerate}
	\loigiai{
		\begin{enumerate}[a)]
			\item Ta có:
			\begin{align*}
				A &= -5a^2b^3 \cdot 2ab^2 \\
				&= (-5 \cdot 2) \cdot (a^2 \cdot a) \cdot (b^3 \cdot b^2) \\
				&= -10a^3b^5
			\end{align*}
			\item Ta có:
			\begin{align*}
				B &= 4a^3b \cdot 3ab^2 - 17a^4b^3 \\
				&= 12a^4b^3 - 17a^4b^3 \\
				&= (12 - 17)a^4b^3 \\
				&= -5a^4b^3
			\end{align*}
			\item Ta có:
			\begin{align*}
				C &= 24a^7b^4 : 6a^4b^2 - 9a^3b^2 \\
				&= 4a^3b^2 - 9a^3b^2 \\
				&= (4 - 9)a^3b^2 \\
				&= -5a^3b^2
			\end{align*}
			\item Ta có:
			\begin{align*}
				D &= 3a^2b^4 + 14a^4b^5 : 2a^2b \\
				&= 3a^2b^4 + 7a^2b^4 \\
				&= 10a^2b^4
			\end{align*}
		\end{enumerate}
	}
\end{ex}

\newpage

% ======================================================================
% HỌC SINH: Phậm Phan Nhật On | MÃ: HS047 | MÃ ĐỀ: 005
% ======================================================================
\noindent
\begin{tcolorbox}[colback=white,colframe=black,arc=0mm,boxrule=0.8pt,top=4pt,bottom=4pt,left=6pt,right=6pt]
	\begin{minipage}[c]{0.44\textwidth}
		{\bfseries\large LỚP TOÁN THẦY THÔNG}\\[3pt]
		\textbf{Họ và tên:} \textbf{\large Phậm Phan Nhật On}\\[3pt]
		\textbf{Lớp:} LỚP 10 \quad -- \quad \textbf{Mã HS:} \texttt{HS047}
	\end{minipage}
	\hfill
	\begin{minipage}[c]{0.22\textwidth}
		\begin{tcolorbox}[colback=white,colframe=black,arc=0mm,boxrule=0.6pt,height=1.8cm,halign=center,valign=top,top=2pt]
			{\bfseries\normalsize Điểm}
		\end{tcolorbox}
	\end{minipage}
	\hfill
	\begin{minipage}[c]{0.30\textwidth}
		\raggedleft
		\textit{Thời gian làm bài: 30 phút}\\[4pt]
		\textbf{Mã đề thi:} \fbox{\textbf{\Large\strut\ 005\ }}
	\end{minipage}
\end{tcolorbox}
\vspace{0.25cm}

\noindent\textbf{\large ĐỀ BÀI (MÃ ĐỀ: 005)}
\vspace{0.1cm}

% --- CÂU 1 (TRANG 1) ---
\begin{ex}
	Thực hiện phép tính:
	\begin{enumerate}[a)]
		\item $A = 3x^3y \cdot 4xy^2$
		\item $B = 3x^2y^3 \cdot 2x^2y - 11x^4y^4$
		\item $C = 18x^6y^5 : 3x^2y^2 - 12x^4y^3$
		\item $D = 4x^3y^2 + 8x^5y^3 : 2x^2y$
	\end{enumerate}
	\loigiai{
		\begin{enumerate}[a)]
			\item Ta có:
			\begin{align*}
				A &= 3x^3y \cdot 4xy^2 \\
				&= (3 \cdot 4) \cdot (x^3 \cdot x) \cdot (y \cdot y^2) \\
				&= 12x^4y^3
			\end{align*}
			\item Ta có:
			\begin{align*}
				B &= 3x^2y^3 \cdot 2x^2y - 11x^4y^4 \\
				&= 6x^4y^4 - 11x^4y^4 \\
				&= (6 - 11)x^4y^4 \\
				&= -5x^4y^4
			\end{align*}
			\item Ta có:
			\begin{align*}
				C &= 18x^6y^5 : 3x^2y^2 - 12x^4y^3 \\
				&= 6x^4y^3 - 12x^4y^3 \\
				&= (6 - 12)x^4y^3 \\
				&= -6x^4y^3
			\end{align*}
			\item Ta có:
			\begin{align*}
				D &= 4x^3y^2 + 8x^5y^3 : 2x^2y \\
				&= 4x^3y^2 + 4x^3y^2 \\
				&= 8x^3y^2
			\end{align*}
		\end{enumerate}
	}
\end{ex}

\newpage
% ----------------------------------------------------------------------
% TRANG 2 (MẶT SAU) -- DÒNG KẺ TIẾP THEO
% ----------------------------------------------------------------------
\noindent
\textbf{Học sinh:} Phậm Phan Nhật On \quad -- \quad \textbf{Mã đề:} \fbox{\textbf{005}} \hfill \textit{Trang 2/2}\par
\vspace{-0.1cm}
\noindent\hrulefill
\vspace{0.3cm}

% --- CÂU 2 (TRANG 2) ---
\begin{ex}
	Thực hiện phép tính:
	\begin{enumerate}[a)]
		\item $A = -5a^2b^3 \cdot 2ab^2$
		\item $B = 4a^3b \cdot 3ab^2 - 17a^4b^3$
		\item $C = 24a^7b^4 : 6a^4b^2 - 9a^3b^2$
		\item $D = 3a^2b^4 + 14a^4b^5 : 2a^2b$
	\end{enumerate}
	\loigiai{
		\begin{enumerate}[a)]
			\item Ta có:
			\begin{align*}
				A &= -5a^2b^3 \cdot 2ab^2 \\
				&= (-5 \cdot 2) \cdot (a^2 \cdot a) \cdot (b^3 \cdot b^2) \\
				&= -10a^3b^5
			\end{align*}
			\item Ta có:
			\begin{align*}
				B &= 4a^3b \cdot 3ab^2 - 17a^4b^3 \\
				&= 12a^4b^3 - 17a^4b^3 \\
				&= (12 - 17)a^4b^3 \\
				&= -5a^4b^3
			\end{align*}
			\item Ta có:
			\begin{align*}
				C &= 24a^7b^4 : 6a^4b^2 - 9a^3b^2 \\
				&= 4a^3b^2 - 9a^3b^2 \\
				&= (4 - 9)a^3b^2 \\
				&= -5a^3b^2
			\end{align*}
			\item Ta có:
			\begin{align*}
				D &= 3a^2b^4 + 14a^4b^5 : 2a^2b \\
				&= 3a^2b^4 + 7a^2b^4 \\
				&= 10a^2b^4
			\end{align*}
		\end{enumerate}
	}
\end{ex}

\newpage

% ======================================================================
% HỌC SINH: Phúc Hưng | MÃ: HS037 | MÃ ĐỀ: 008
% ======================================================================
\noindent
\begin{tcolorbox}[colback=white,colframe=black,arc=0mm,boxrule=0.8pt,top=4pt,bottom=4pt,left=6pt,right=6pt]
	\begin{minipage}[c]{0.44\textwidth}
		{\bfseries\large LỚP TOÁN THẦY THÔNG}\\[3pt]
		\textbf{Họ và tên:} \textbf{\large Phúc Hưng}\\[3pt]
		\textbf{Lớp:} LỚP 10 \quad -- \quad \textbf{Mã HS:} \texttt{HS037}
	\end{minipage}
	\hfill
	\begin{minipage}[c]{0.22\textwidth}
		\begin{tcolorbox}[colback=white,colframe=black,arc=0mm,boxrule=0.6pt,height=1.8cm,halign=center,valign=top,top=2pt]
			{\bfseries\normalsize Điểm}
		\end{tcolorbox}
	\end{minipage}
	\hfill
	\begin{minipage}[c]{0.30\textwidth}
		\raggedleft
		\textit{Thời gian làm bài: 30 phút}\\[4pt]
		\textbf{Mã đề thi:} \fbox{\textbf{\Large\strut\ 008\ }}
	\end{minipage}
\end{tcolorbox}
\vspace{0.25cm}

\noindent\textbf{\large ĐỀ BÀI (MÃ ĐỀ: 008)}
\vspace{0.1cm}

% --- CÂU 1 (TRANG 1) ---
\begin{ex}
	Thực hiện phép tính:
	\begin{enumerate}[a)]
		\item $A = 3x^3y \cdot 4xy^2$
		\item $B = 3x^2y^3 \cdot 2x^2y - 11x^4y^4$
		\item $C = 18x^6y^5 : 3x^2y^2 - 12x^4y^3$
		\item $D = 4x^3y^2 + 8x^5y^3 : 2x^2y$
	\end{enumerate}
	\loigiai{
		\begin{enumerate}[a)]
			\item Ta có:
			\begin{align*}
				A &= 3x^3y \cdot 4xy^2 \\
				&= (3 \cdot 4) \cdot (x^3 \cdot x) \cdot (y \cdot y^2) \\
				&= 12x^4y^3
			\end{align*}
			\item Ta có:
			\begin{align*}
				B &= 3x^2y^3 \cdot 2x^2y - 11x^4y^4 \\
				&= 6x^4y^4 - 11x^4y^4 \\
				&= (6 - 11)x^4y^4 \\
				&= -5x^4y^4
			\end{align*}
			\item Ta có:
			\begin{align*}
				C &= 18x^6y^5 : 3x^2y^2 - 12x^4y^3 \\
				&= 6x^4y^3 - 12x^4y^3 \\
				&= (6 - 12)x^4y^3 \\
				&= -6x^4y^3
			\end{align*}
			\item Ta có:
			\begin{align*}
				D &= 4x^3y^2 + 8x^5y^3 : 2x^2y \\
				&= 4x^3y^2 + 4x^3y^2 \\
				&= 8x^3y^2
			\end{align*}
		\end{enumerate}
	}
\end{ex}

\newpage
% ----------------------------------------------------------------------
% TRANG 2 (MẶT SAU) -- DÒNG KẺ TIẾP THEO
% ----------------------------------------------------------------------
\noindent
\textbf{Học sinh:} Phúc Hưng \quad -- \quad \textbf{Mã đề:} \fbox{\textbf{008}} \hfill \textit{Trang 2/2}\par
\vspace{-0.1cm}
\noindent\hrulefill
\vspace{0.3cm}

% --- CÂU 2 (TRANG 2) ---
\begin{ex}
	Thực hiện phép tính:
	\begin{enumerate}[a)]
		\item $A = -5a^2b^3 \cdot 2ab^2$
		\item $B = 4a^3b \cdot 3ab^2 - 17a^4b^3$
		\item $C = 24a^7b^4 : 6a^4b^2 - 9a^3b^2$
		\item $D = 3a^2b^4 + 14a^4b^5 : 2a^2b$
	\end{enumerate}
	\loigiai{
		\begin{enumerate}[a)]
			\item Ta có:
			\begin{align*}
				A &= -5a^2b^3 \cdot 2ab^2 \\
				&= (-5 \cdot 2) \cdot (a^2 \cdot a) \cdot (b^3 \cdot b^2) \\
				&= -10a^3b^5
			\end{align*}
			\item Ta có:
			\begin{align*}
				B &= 4a^3b \cdot 3ab^2 - 17a^4b^3 \\
				&= 12a^4b^3 - 17a^4b^3 \\
				&= (12 - 17)a^4b^3 \\
				&= -5a^4b^3
			\end{align*}
			\item Ta có:
			\begin{align*}
				C &= 24a^7b^4 : 6a^4b^2 - 9a^3b^2 \\
				&= 4a^3b^2 - 9a^3b^2 \\
				&= (4 - 9)a^3b^2 \\
				&= -5a^3b^2
			\end{align*}
			\item Ta có:
			\begin{align*}
				D &= 3a^2b^4 + 14a^4b^5 : 2a^2b \\
				&= 3a^2b^4 + 7a^2b^4 \\
				&= 10a^2b^4
			\end{align*}
		\end{enumerate}
	}
\end{ex}

\newpage

% ======================================================================
% HỌC SINH: Phương Anh | MÃ: HS022 | MÃ ĐỀ: 010
% ======================================================================
\noindent
\begin{tcolorbox}[colback=white,colframe=black,arc=0mm,boxrule=0.8pt,top=4pt,bottom=4pt,left=6pt,right=6pt]
	\begin{minipage}[c]{0.44\textwidth}
		{\bfseries\large LỚP TOÁN THẦY THÔNG}\\[3pt]
		\textbf{Họ và tên:} \textbf{\large Phương Anh}\\[3pt]
		\textbf{Lớp:} LỚP 10 \quad -- \quad \textbf{Mã HS:} \texttt{HS022}
	\end{minipage}
	\hfill
	\begin{minipage}[c]{0.22\textwidth}
		\begin{tcolorbox}[colback=white,colframe=black,arc=0mm,boxrule=0.6pt,height=1.8cm,halign=center,valign=top,top=2pt]
			{\bfseries\normalsize Điểm}
		\end{tcolorbox}
	\end{minipage}
	\hfill
	\begin{minipage}[c]{0.30\textwidth}
		\raggedleft
		\textit{Thời gian làm bài: 30 phút}\\[4pt]
		\textbf{Mã đề thi:} \fbox{\textbf{\Large\strut\ 010\ }}
	\end{minipage}
\end{tcolorbox}
\vspace{0.25cm}

\noindent\textbf{\large ĐỀ BÀI (MÃ ĐỀ: 010)}
\vspace{0.1cm}

% --- CÂU 1 (TRANG 1) ---
\begin{ex}
	Thực hiện phép tính:
	\begin{enumerate}[a)]
		\item $A = 3x^3y \cdot 4xy^2$
		\item $B = 3x^2y^3 \cdot 2x^2y - 11x^4y^4$
		\item $C = 18x^6y^5 : 3x^2y^2 - 12x^4y^3$
		\item $D = 4x^3y^2 + 8x^5y^3 : 2x^2y$
	\end{enumerate}
	\loigiai{
		\begin{enumerate}[a)]
			\item Ta có:
			\begin{align*}
				A &= 3x^3y \cdot 4xy^2 \\
				&= (3 \cdot 4) \cdot (x^3 \cdot x) \cdot (y \cdot y^2) \\
				&= 12x^4y^3
			\end{align*}
			\item Ta có:
			\begin{align*}
				B &= 3x^2y^3 \cdot 2x^2y - 11x^4y^4 \\
				&= 6x^4y^4 - 11x^4y^4 \\
				&= (6 - 11)x^4y^4 \\
				&= -5x^4y^4
			\end{align*}
			\item Ta có:
			\begin{align*}
				C &= 18x^6y^5 : 3x^2y^2 - 12x^4y^3 \\
				&= 6x^4y^3 - 12x^4y^3 \\
				&= (6 - 12)x^4y^3 \\
				&= -6x^4y^3
			\end{align*}
			\item Ta có:
			\begin{align*}
				D &= 4x^3y^2 + 8x^5y^3 : 2x^2y \\
				&= 4x^3y^2 + 4x^3y^2 \\
				&= 8x^3y^2
			\end{align*}
		\end{enumerate}
	}
\end{ex}

\newpage
% ----------------------------------------------------------------------
% TRANG 2 (MẶT SAU) -- DÒNG KẺ TIẾP THEO
% ----------------------------------------------------------------------
\noindent
\textbf{Học sinh:} Phương Anh \quad -- \quad \textbf{Mã đề:} \fbox{\textbf{010}} \hfill \textit{Trang 2/2}\par
\vspace{-0.1cm}
\noindent\hrulefill
\vspace{0.3cm}

% --- CÂU 2 (TRANG 2) ---
\begin{ex}
	Thực hiện phép tính:
	\begin{enumerate}[a)]
		\item $A = -5a^2b^3 \cdot 2ab^2$
		\item $B = 4a^3b \cdot 3ab^2 - 17a^4b^3$
		\item $C = 24a^7b^4 : 6a^4b^2 - 9a^3b^2$
		\item $D = 3a^2b^4 + 14a^4b^5 : 2a^2b$
	\end{enumerate}
	\loigiai{
		\begin{enumerate}[a)]
			\item Ta có:
			\begin{align*}
				A &= -5a^2b^3 \cdot 2ab^2 \\
				&= (-5 \cdot 2) \cdot (a^2 \cdot a) \cdot (b^3 \cdot b^2) \\
				&= -10a^3b^5
			\end{align*}
			\item Ta có:
			\begin{align*}
				B &= 4a^3b \cdot 3ab^2 - 17a^4b^3 \\
				&= 12a^4b^3 - 17a^4b^3 \\
				&= (12 - 17)a^4b^3 \\
				&= -5a^4b^3
			\end{align*}
			\item Ta có:
			\begin{align*}
				C &= 24a^7b^4 : 6a^4b^2 - 9a^3b^2 \\
				&= 4a^3b^2 - 9a^3b^2 \\
				&= (4 - 9)a^3b^2 \\
				&= -5a^3b^2
			\end{align*}
			\item Ta có:
			\begin{align*}
				D &= 3a^2b^4 + 14a^4b^5 : 2a^2b \\
				&= 3a^2b^4 + 7a^2b^4 \\
				&= 10a^2b^4
			\end{align*}
		\end{enumerate}
	}
\end{ex}

\newpage

% ======================================================================
% HỌC SINH: Phương Thảo | MÃ: HS038 | MÃ ĐỀ: 009
% ======================================================================
\noindent
\begin{tcolorbox}[colback=white,colframe=black,arc=0mm,boxrule=0.8pt,top=4pt,bottom=4pt,left=6pt,right=6pt]
	\begin{minipage}[c]{0.44\textwidth}
		{\bfseries\large LỚP TOÁN THẦY THÔNG}\\[3pt]
		\textbf{Họ và tên:} \textbf{\large Phương Thảo}\\[3pt]
		\textbf{Lớp:} LỚP 10 \quad -- \quad \textbf{Mã HS:} \texttt{HS038}
	\end{minipage}
	\hfill
	\begin{minipage}[c]{0.22\textwidth}
		\begin{tcolorbox}[colback=white,colframe=black,arc=0mm,boxrule=0.6pt,height=1.8cm,halign=center,valign=top,top=2pt]
			{\bfseries\normalsize Điểm}
		\end{tcolorbox}
	\end{minipage}
	\hfill
	\begin{minipage}[c]{0.30\textwidth}
		\raggedleft
		\textit{Thời gian làm bài: 30 phút}\\[4pt]
		\textbf{Mã đề thi:} \fbox{\textbf{\Large\strut\ 009\ }}
	\end{minipage}
\end{tcolorbox}
\vspace{0.25cm}

\noindent\textbf{\large ĐỀ BÀI (MÃ ĐỀ: 009)}
\vspace{0.1cm}

% --- CÂU 1 (TRANG 1) ---
\begin{ex}
	Thực hiện phép tính:
	\begin{enumerate}[a)]
		\item $A = 3x^3y \cdot 4xy^2$
		\item $B = 3x^2y^3 \cdot 2x^2y - 11x^4y^4$
		\item $C = 18x^6y^5 : 3x^2y^2 - 12x^4y^3$
		\item $D = 4x^3y^2 + 8x^5y^3 : 2x^2y$
	\end{enumerate}
	\loigiai{
		\begin{enumerate}[a)]
			\item Ta có:
			\begin{align*}
				A &= 3x^3y \cdot 4xy^2 \\
				&= (3 \cdot 4) \cdot (x^3 \cdot x) \cdot (y \cdot y^2) \\
				&= 12x^4y^3
			\end{align*}
			\item Ta có:
			\begin{align*}
				B &= 3x^2y^3 \cdot 2x^2y - 11x^4y^4 \\
				&= 6x^4y^4 - 11x^4y^4 \\
				&= (6 - 11)x^4y^4 \\
				&= -5x^4y^4
			\end{align*}
			\item Ta có:
			\begin{align*}
				C &= 18x^6y^5 : 3x^2y^2 - 12x^4y^3 \\
				&= 6x^4y^3 - 12x^4y^3 \\
				&= (6 - 12)x^4y^3 \\
				&= -6x^4y^3
			\end{align*}
			\item Ta có:
			\begin{align*}
				D &= 4x^3y^2 + 8x^5y^3 : 2x^2y \\
				&= 4x^3y^2 + 4x^3y^2 \\
				&= 8x^3y^2
			\end{align*}
		\end{enumerate}
	}
\end{ex}

\newpage
% ----------------------------------------------------------------------
% TRANG 2 (MẶT SAU) -- DÒNG KẺ TIẾP THEO
% ----------------------------------------------------------------------
\noindent
\textbf{Học sinh:} Phương Thảo \quad -- \quad \textbf{Mã đề:} \fbox{\textbf{009}} \hfill \textit{Trang 2/2}\par
\vspace{-0.1cm}
\noindent\hrulefill
\vspace{0.3cm}

% --- CÂU 2 (TRANG 2) ---
\begin{ex}
	Thực hiện phép tính:
	\begin{enumerate}[a)]
		\item $A = -5a^2b^3 \cdot 2ab^2$
		\item $B = 4a^3b \cdot 3ab^2 - 17a^4b^3$
		\item $C = 24a^7b^4 : 6a^4b^2 - 9a^3b^2$
		\item $D = 3a^2b^4 + 14a^4b^5 : 2a^2b$
	\end{enumerate}
	\loigiai{
		\begin{enumerate}[a)]
			\item Ta có:
			\begin{align*}
				A &= -5a^2b^3 \cdot 2ab^2 \\
				&= (-5 \cdot 2) \cdot (a^2 \cdot a) \cdot (b^3 \cdot b^2) \\
				&= -10a^3b^5
			\end{align*}
			\item Ta có:
			\begin{align*}
				B &= 4a^3b \cdot 3ab^2 - 17a^4b^3 \\
				&= 12a^4b^3 - 17a^4b^3 \\
				&= (12 - 17)a^4b^3 \\
				&= -5a^4b^3
			\end{align*}
			\item Ta có:
			\begin{align*}
				C &= 24a^7b^4 : 6a^4b^2 - 9a^3b^2 \\
				&= 4a^3b^2 - 9a^3b^2 \\
				&= (4 - 9)a^3b^2 \\
				&= -5a^3b^2
			\end{align*}
			\item Ta có:
			\begin{align*}
				D &= 3a^2b^4 + 14a^4b^5 : 2a^2b \\
				&= 3a^2b^4 + 7a^2b^4 \\
				&= 10a^2b^4
			\end{align*}
		\end{enumerate}
	}
\end{ex}

\newpage

% ======================================================================
% HỌC SINH: Quang Huy-01-2-2001 | MÃ: HS050 | MÃ ĐỀ: 006
% ======================================================================
\noindent
\begin{tcolorbox}[colback=white,colframe=black,arc=0mm,boxrule=0.8pt,top=4pt,bottom=4pt,left=6pt,right=6pt]
	\begin{minipage}[c]{0.44\textwidth}
		{\bfseries\large LỚP TOÁN THẦY THÔNG}\\[3pt]
		\textbf{Họ và tên:} \textbf{\large Quang Huy-01-2-2001}\\[3pt]
		\textbf{Lớp:} LỚP 10 \quad -- \quad \textbf{Mã HS:} \texttt{HS050}
	\end{minipage}
	\hfill
	\begin{minipage}[c]{0.22\textwidth}
		\begin{tcolorbox}[colback=white,colframe=black,arc=0mm,boxrule=0.6pt,height=1.8cm,halign=center,valign=top,top=2pt]
			{\bfseries\normalsize Điểm}
		\end{tcolorbox}
	\end{minipage}
	\hfill
	\begin{minipage}[c]{0.30\textwidth}
		\raggedleft
		\textit{Thời gian làm bài: 30 phút}\\[4pt]
		\textbf{Mã đề thi:} \fbox{\textbf{\Large\strut\ 006\ }}
	\end{minipage}
\end{tcolorbox}
\vspace{0.25cm}

\noindent\textbf{\large ĐỀ BÀI (MÃ ĐỀ: 006)}
\vspace{0.1cm}

% --- CÂU 1 (TRANG 1) ---
\begin{ex}
	Thực hiện phép tính:
	\begin{enumerate}[a)]
		\item $A = 3x^3y \cdot 4xy^2$
		\item $B = 3x^2y^3 \cdot 2x^2y - 11x^4y^4$
		\item $C = 18x^6y^5 : 3x^2y^2 - 12x^4y^3$
		\item $D = 4x^3y^2 + 8x^5y^3 : 2x^2y$
	\end{enumerate}
	\loigiai{
		\begin{enumerate}[a)]
			\item Ta có:
			\begin{align*}
				A &= 3x^3y \cdot 4xy^2 \\
				&= (3 \cdot 4) \cdot (x^3 \cdot x) \cdot (y \cdot y^2) \\
				&= 12x^4y^3
			\end{align*}
			\item Ta có:
			\begin{align*}
				B &= 3x^2y^3 \cdot 2x^2y - 11x^4y^4 \\
				&= 6x^4y^4 - 11x^4y^4 \\
				&= (6 - 11)x^4y^4 \\
				&= -5x^4y^4
			\end{align*}
			\item Ta có:
			\begin{align*}
				C &= 18x^6y^5 : 3x^2y^2 - 12x^4y^3 \\
				&= 6x^4y^3 - 12x^4y^3 \\
				&= (6 - 12)x^4y^3 \\
				&= -6x^4y^3
			\end{align*}
			\item Ta có:
			\begin{align*}
				D &= 4x^3y^2 + 8x^5y^3 : 2x^2y \\
				&= 4x^3y^2 + 4x^3y^2 \\
				&= 8x^3y^2
			\end{align*}
		\end{enumerate}
	}
\end{ex}

\newpage
% ----------------------------------------------------------------------
% TRANG 2 (MẶT SAU) -- DÒNG KẺ TIẾP THEO
% ----------------------------------------------------------------------
\noindent
\textbf{Học sinh:} Quang Huy-01-2-2001 \quad -- \quad \textbf{Mã đề:} \fbox{\textbf{006}} \hfill \textit{Trang 2/2}\par
\vspace{-0.1cm}
\noindent\hrulefill
\vspace{0.3cm}

% --- CÂU 2 (TRANG 2) ---
\begin{ex}
	Thực hiện phép tính:
	\begin{enumerate}[a)]
		\item $A = -5a^2b^3 \cdot 2ab^2$
		\item $B = 4a^3b \cdot 3ab^2 - 17a^4b^3$
		\item $C = 24a^7b^4 : 6a^4b^2 - 9a^3b^2$
		\item $D = 3a^2b^4 + 14a^4b^5 : 2a^2b$
	\end{enumerate}
	\loigiai{
		\begin{enumerate}[a)]
			\item Ta có:
			\begin{align*}
				A &= -5a^2b^3 \cdot 2ab^2 \\
				&= (-5 \cdot 2) \cdot (a^2 \cdot a) \cdot (b^3 \cdot b^2) \\
				&= -10a^3b^5
			\end{align*}
			\item Ta có:
			\begin{align*}
				B &= 4a^3b \cdot 3ab^2 - 17a^4b^3 \\
				&= 12a^4b^3 - 17a^4b^3 \\
				&= (12 - 17)a^4b^3 \\
				&= -5a^4b^3
			\end{align*}
			\item Ta có:
			\begin{align*}
				C &= 24a^7b^4 : 6a^4b^2 - 9a^3b^2 \\
				&= 4a^3b^2 - 9a^3b^2 \\
				&= (4 - 9)a^3b^2 \\
				&= -5a^3b^2
			\end{align*}
			\item Ta có:
			\begin{align*}
				D &= 3a^2b^4 + 14a^4b^5 : 2a^2b \\
				&= 3a^2b^4 + 7a^2b^4 \\
				&= 10a^2b^4
			\end{align*}
		\end{enumerate}
	}
\end{ex}

\newpage

% ======================================================================
% HỌC SINH: Quốc Khánh | MÃ: HS002 | MÃ ĐỀ: 001
% ======================================================================
\noindent
\begin{tcolorbox}[colback=white,colframe=black,arc=0mm,boxrule=0.8pt,top=4pt,bottom=4pt,left=6pt,right=6pt]
	\begin{minipage}[c]{0.44\textwidth}
		{\bfseries\large LỚP TOÁN THẦY THÔNG}\\[3pt]
		\textbf{Họ và tên:} \textbf{\large Quốc Khánh}\\[3pt]
		\textbf{Lớp:} LỚP 10 \quad -- \quad \textbf{Mã HS:} \texttt{HS002}
	\end{minipage}
	\hfill
	\begin{minipage}[c]{0.22\textwidth}
		\begin{tcolorbox}[colback=white,colframe=black,arc=0mm,boxrule=0.6pt,height=1.8cm,halign=center,valign=top,top=2pt]
			{\bfseries\normalsize Điểm}
		\end{tcolorbox}
	\end{minipage}
	\hfill
	\begin{minipage}[c]{0.30\textwidth}
		\raggedleft
		\textit{Thời gian làm bài: 30 phút}\\[4pt]
		\textbf{Mã đề thi:} \fbox{\textbf{\Large\strut\ 001\ }}
	\end{minipage}
\end{tcolorbox}
\vspace{0.25cm}

\noindent\textbf{\large ĐỀ BÀI (MÃ ĐỀ: 001)}
\vspace{0.1cm}

% --- CÂU 1 (TRANG 1) ---
\begin{ex}
	Thực hiện phép tính:
	\begin{enumerate}[a)]
		\item $A = 3x^3y \cdot 4xy^2$
		\item $B = 3x^2y^3 \cdot 2x^2y - 11x^4y^4$
		\item $C = 18x^6y^5 : 3x^2y^2 - 12x^4y^3$
		\item $D = 4x^3y^2 + 8x^5y^3 : 2x^2y$
	\end{enumerate}
	\loigiai{
		\begin{enumerate}[a)]
			\item Ta có:
			\begin{align*}
				A &= 3x^3y \cdot 4xy^2 \\
				&= (3 \cdot 4) \cdot (x^3 \cdot x) \cdot (y \cdot y^2) \\
				&= 12x^4y^3
			\end{align*}
			\item Ta có:
			\begin{align*}
				B &= 3x^2y^3 \cdot 2x^2y - 11x^4y^4 \\
				&= 6x^4y^4 - 11x^4y^4 \\
				&= (6 - 11)x^4y^4 \\
				&= -5x^4y^4
			\end{align*}
			\item Ta có:
			\begin{align*}
				C &= 18x^6y^5 : 3x^2y^2 - 12x^4y^3 \\
				&= 6x^4y^3 - 12x^4y^3 \\
				&= (6 - 12)x^4y^3 \\
				&= -6x^4y^3
			\end{align*}
			\item Ta có:
			\begin{align*}
				D &= 4x^3y^2 + 8x^5y^3 : 2x^2y \\
				&= 4x^3y^2 + 4x^3y^2 \\
				&= 8x^3y^2
			\end{align*}
		\end{enumerate}
	}
\end{ex}

\newpage
% ----------------------------------------------------------------------
% TRANG 2 (MẶT SAU) -- DÒNG KẺ TIẾP THEO
% ----------------------------------------------------------------------
\noindent
\textbf{Học sinh:} Quốc Khánh \quad -- \quad \textbf{Mã đề:} \fbox{\textbf{001}} \hfill \textit{Trang 2/2}\par
\vspace{-0.1cm}
\noindent\hrulefill
\vspace{0.3cm}

% --- CÂU 2 (TRANG 2) ---
\begin{ex}
	Thực hiện phép tính:
	\begin{enumerate}[a)]
		\item $A = -5a^2b^3 \cdot 2ab^2$
		\item $B = 4a^3b \cdot 3ab^2 - 17a^4b^3$
		\item $C = 24a^7b^4 : 6a^4b^2 - 9a^3b^2$
		\item $D = 3a^2b^4 + 14a^4b^5 : 2a^2b$
	\end{enumerate}
	\loigiai{
		\begin{enumerate}[a)]
			\item Ta có:
			\begin{align*}
				A &= -5a^2b^3 \cdot 2ab^2 \\
				&= (-5 \cdot 2) \cdot (a^2 \cdot a) \cdot (b^3 \cdot b^2) \\
				&= -10a^3b^5
			\end{align*}
			\item Ta có:
			\begin{align*}
				B &= 4a^3b \cdot 3ab^2 - 17a^4b^3 \\
				&= 12a^4b^3 - 17a^4b^3 \\
				&= (12 - 17)a^4b^3 \\
				&= -5a^4b^3
			\end{align*}
			\item Ta có:
			\begin{align*}
				C &= 24a^7b^4 : 6a^4b^2 - 9a^3b^2 \\
				&= 4a^3b^2 - 9a^3b^2 \\
				&= (4 - 9)a^3b^2 \\
				&= -5a^3b^2
			\end{align*}
			\item Ta có:
			\begin{align*}
				D &= 3a^2b^4 + 14a^4b^5 : 2a^2b \\
				&= 3a^2b^4 + 7a^2b^4 \\
				&= 10a^2b^4
			\end{align*}
		\end{enumerate}
	}
\end{ex}

\newpage

% ======================================================================
% HỌC SINH: Quý Công On | MÃ: HS046 | MÃ ĐỀ: 006
% ======================================================================
\noindent
\begin{tcolorbox}[colback=white,colframe=black,arc=0mm,boxrule=0.8pt,top=4pt,bottom=4pt,left=6pt,right=6pt]
	\begin{minipage}[c]{0.44\textwidth}
		{\bfseries\large LỚP TOÁN THẦY THÔNG}\\[3pt]
		\textbf{Họ và tên:} \textbf{\large Quý Công On}\\[3pt]
		\textbf{Lớp:} LỚP 10 \quad -- \quad \textbf{Mã HS:} \texttt{HS046}
	\end{minipage}
	\hfill
	\begin{minipage}[c]{0.22\textwidth}
		\begin{tcolorbox}[colback=white,colframe=black,arc=0mm,boxrule=0.6pt,height=1.8cm,halign=center,valign=top,top=2pt]
			{\bfseries\normalsize Điểm}
		\end{tcolorbox}
	\end{minipage}
	\hfill
	\begin{minipage}[c]{0.30\textwidth}
		\raggedleft
		\textit{Thời gian làm bài: 30 phút}\\[4pt]
		\textbf{Mã đề thi:} \fbox{\textbf{\Large\strut\ 006\ }}
	\end{minipage}
\end{tcolorbox}
\vspace{0.25cm}

\noindent\textbf{\large ĐỀ BÀI (MÃ ĐỀ: 006)}
\vspace{0.1cm}

% --- CÂU 1 (TRANG 1) ---
\begin{ex}
	Thực hiện phép tính:
	\begin{enumerate}[a)]
		\item $A = 3x^3y \cdot 4xy^2$
		\item $B = 3x^2y^3 \cdot 2x^2y - 11x^4y^4$
		\item $C = 18x^6y^5 : 3x^2y^2 - 12x^4y^3$
		\item $D = 4x^3y^2 + 8x^5y^3 : 2x^2y$
	\end{enumerate}
	\loigiai{
		\begin{enumerate}[a)]
			\item Ta có:
			\begin{align*}
				A &= 3x^3y \cdot 4xy^2 \\
				&= (3 \cdot 4) \cdot (x^3 \cdot x) \cdot (y \cdot y^2) \\
				&= 12x^4y^3
			\end{align*}
			\item Ta có:
			\begin{align*}
				B &= 3x^2y^3 \cdot 2x^2y - 11x^4y^4 \\
				&= 6x^4y^4 - 11x^4y^4 \\
				&= (6 - 11)x^4y^4 \\
				&= -5x^4y^4
			\end{align*}
			\item Ta có:
			\begin{align*}
				C &= 18x^6y^5 : 3x^2y^2 - 12x^4y^3 \\
				&= 6x^4y^3 - 12x^4y^3 \\
				&= (6 - 12)x^4y^3 \\
				&= -6x^4y^3
			\end{align*}
			\item Ta có:
			\begin{align*}
				D &= 4x^3y^2 + 8x^5y^3 : 2x^2y \\
				&= 4x^3y^2 + 4x^3y^2 \\
				&= 8x^3y^2
			\end{align*}
		\end{enumerate}
	}
\end{ex}

\newpage
% ----------------------------------------------------------------------
% TRANG 2 (MẶT SAU) -- DÒNG KẺ TIẾP THEO
% ----------------------------------------------------------------------
\noindent
\textbf{Học sinh:} Quý Công On \quad -- \quad \textbf{Mã đề:} \fbox{\textbf{006}} \hfill \textit{Trang 2/2}\par
\vspace{-0.1cm}
\noindent\hrulefill
\vspace{0.3cm}

% --- CÂU 2 (TRANG 2) ---
\begin{ex}
	Thực hiện phép tính:
	\begin{enumerate}[a)]
		\item $A = -5a^2b^3 \cdot 2ab^2$
		\item $B = 4a^3b \cdot 3ab^2 - 17a^4b^3$
		\item $C = 24a^7b^4 : 6a^4b^2 - 9a^3b^2$
		\item $D = 3a^2b^4 + 14a^4b^5 : 2a^2b$
	\end{enumerate}
	\loigiai{
		\begin{enumerate}[a)]
			\item Ta có:
			\begin{align*}
				A &= -5a^2b^3 \cdot 2ab^2 \\
				&= (-5 \cdot 2) \cdot (a^2 \cdot a) \cdot (b^3 \cdot b^2) \\
				&= -10a^3b^5
			\end{align*}
			\item Ta có:
			\begin{align*}
				B &= 4a^3b \cdot 3ab^2 - 17a^4b^3 \\
				&= 12a^4b^3 - 17a^4b^3 \\
				&= (12 - 17)a^4b^3 \\
				&= -5a^4b^3
			\end{align*}
			\item Ta có:
			\begin{align*}
				C &= 24a^7b^4 : 6a^4b^2 - 9a^3b^2 \\
				&= 4a^3b^2 - 9a^3b^2 \\
				&= (4 - 9)a^3b^2 \\
				&= -5a^3b^2
			\end{align*}
			\item Ta có:
			\begin{align*}
				D &= 3a^2b^4 + 14a^4b^5 : 2a^2b \\
				&= 3a^2b^4 + 7a^2b^4 \\
				&= 10a^2b^4
			\end{align*}
		\end{enumerate}
	}
\end{ex}

\newpage

% ======================================================================
% HỌC SINH: Quỳnh Trâm | MÃ: HS011 | MÃ ĐỀ: 001
% ======================================================================
\noindent
\begin{tcolorbox}[colback=white,colframe=black,arc=0mm,boxrule=0.8pt,top=4pt,bottom=4pt,left=6pt,right=6pt]
	\begin{minipage}[c]{0.44\textwidth}
		{\bfseries\large LỚP TOÁN THẦY THÔNG}\\[3pt]
		\textbf{Họ và tên:} \textbf{\large Quỳnh Trâm}\\[3pt]
		\textbf{Lớp:} LỚP 10 \quad -- \quad \textbf{Mã HS:} \texttt{HS011}
	\end{minipage}
	\hfill
	\begin{minipage}[c]{0.22\textwidth}
		\begin{tcolorbox}[colback=white,colframe=black,arc=0mm,boxrule=0.6pt,height=1.8cm,halign=center,valign=top,top=2pt]
			{\bfseries\normalsize Điểm}
		\end{tcolorbox}
	\end{minipage}
	\hfill
	\begin{minipage}[c]{0.30\textwidth}
		\raggedleft
		\textit{Thời gian làm bài: 30 phút}\\[4pt]
		\textbf{Mã đề thi:} \fbox{\textbf{\Large\strut\ 001\ }}
	\end{minipage}
\end{tcolorbox}
\vspace{0.25cm}

\noindent\textbf{\large ĐỀ BÀI (MÃ ĐỀ: 001)}
\vspace{0.1cm}

% --- CÂU 1 (TRANG 1) ---
\begin{ex}
	Thực hiện phép tính:
	\begin{enumerate}[a)]
		\item $A = 3x^3y \cdot 4xy^2$
		\item $B = 3x^2y^3 \cdot 2x^2y - 11x^4y^4$
		\item $C = 18x^6y^5 : 3x^2y^2 - 12x^4y^3$
		\item $D = 4x^3y^2 + 8x^5y^3 : 2x^2y$
	\end{enumerate}
	\loigiai{
		\begin{enumerate}[a)]
			\item Ta có:
			\begin{align*}
				A &= 3x^3y \cdot 4xy^2 \\
				&= (3 \cdot 4) \cdot (x^3 \cdot x) \cdot (y \cdot y^2) \\
				&= 12x^4y^3
			\end{align*}
			\item Ta có:
			\begin{align*}
				B &= 3x^2y^3 \cdot 2x^2y - 11x^4y^4 \\
				&= 6x^4y^4 - 11x^4y^4 \\
				&= (6 - 11)x^4y^4 \\
				&= -5x^4y^4
			\end{align*}
			\item Ta có:
			\begin{align*}
				C &= 18x^6y^5 : 3x^2y^2 - 12x^4y^3 \\
				&= 6x^4y^3 - 12x^4y^3 \\
				&= (6 - 12)x^4y^3 \\
				&= -6x^4y^3
			\end{align*}
			\item Ta có:
			\begin{align*}
				D &= 4x^3y^2 + 8x^5y^3 : 2x^2y \\
				&= 4x^3y^2 + 4x^3y^2 \\
				&= 8x^3y^2
			\end{align*}
		\end{enumerate}
	}
\end{ex}

\newpage
% ----------------------------------------------------------------------
% TRANG 2 (MẶT SAU) -- DÒNG KẺ TIẾP THEO
% ----------------------------------------------------------------------
\noindent
\textbf{Học sinh:} Quỳnh Trâm \quad -- \quad \textbf{Mã đề:} \fbox{\textbf{001}} \hfill \textit{Trang 2/2}\par
\vspace{-0.1cm}
\noindent\hrulefill
\vspace{0.3cm}

% --- CÂU 2 (TRANG 2) ---
\begin{ex}
	Thực hiện phép tính:
	\begin{enumerate}[a)]
		\item $A = -5a^2b^3 \cdot 2ab^2$
		\item $B = 4a^3b \cdot 3ab^2 - 17a^4b^3$
		\item $C = 24a^7b^4 : 6a^4b^2 - 9a^3b^2$
		\item $D = 3a^2b^4 + 14a^4b^5 : 2a^2b$
	\end{enumerate}
	\loigiai{
		\begin{enumerate}[a)]
			\item Ta có:
			\begin{align*}
				A &= -5a^2b^3 \cdot 2ab^2 \\
				&= (-5 \cdot 2) \cdot (a^2 \cdot a) \cdot (b^3 \cdot b^2) \\
				&= -10a^3b^5
			\end{align*}
			\item Ta có:
			\begin{align*}
				B &= 4a^3b \cdot 3ab^2 - 17a^4b^3 \\
				&= 12a^4b^3 - 17a^4b^3 \\
				&= (12 - 17)a^4b^3 \\
				&= -5a^4b^3
			\end{align*}
			\item Ta có:
			\begin{align*}
				C &= 24a^7b^4 : 6a^4b^2 - 9a^3b^2 \\
				&= 4a^3b^2 - 9a^3b^2 \\
				&= (4 - 9)a^3b^2 \\
				&= -5a^3b^2
			\end{align*}
			\item Ta có:
			\begin{align*}
				D &= 3a^2b^4 + 14a^4b^5 : 2a^2b \\
				&= 3a^2b^4 + 7a^2b^4 \\
				&= 10a^2b^4
			\end{align*}
		\end{enumerate}
	}
\end{ex}

\newpage

% ======================================================================
% HỌC SINH: Thế Anh | MÃ: HS028 | MÃ ĐỀ: 009
% ======================================================================
\noindent
\begin{tcolorbox}[colback=white,colframe=black,arc=0mm,boxrule=0.8pt,top=4pt,bottom=4pt,left=6pt,right=6pt]
	\begin{minipage}[c]{0.44\textwidth}
		{\bfseries\large LỚP TOÁN THẦY THÔNG}\\[3pt]
		\textbf{Họ và tên:} \textbf{\large Thế Anh}\\[3pt]
		\textbf{Lớp:} LỚP 10 \quad -- \quad \textbf{Mã HS:} \texttt{HS028}
	\end{minipage}
	\hfill
	\begin{minipage}[c]{0.22\textwidth}
		\begin{tcolorbox}[colback=white,colframe=black,arc=0mm,boxrule=0.6pt,height=1.8cm,halign=center,valign=top,top=2pt]
			{\bfseries\normalsize Điểm}
		\end{tcolorbox}
	\end{minipage}
	\hfill
	\begin{minipage}[c]{0.30\textwidth}
		\raggedleft
		\textit{Thời gian làm bài: 30 phút}\\[4pt]
		\textbf{Mã đề thi:} \fbox{\textbf{\Large\strut\ 009\ }}
	\end{minipage}
\end{tcolorbox}
\vspace{0.25cm}

\noindent\textbf{\large ĐỀ BÀI (MÃ ĐỀ: 009)}
\vspace{0.1cm}

% --- CÂU 1 (TRANG 1) ---
\begin{ex}
	Thực hiện phép tính:
	\begin{enumerate}[a)]
		\item $A = 3x^3y \cdot 4xy^2$
		\item $B = 3x^2y^3 \cdot 2x^2y - 11x^4y^4$
		\item $C = 18x^6y^5 : 3x^2y^2 - 12x^4y^3$
		\item $D = 4x^3y^2 + 8x^5y^3 : 2x^2y$
	\end{enumerate}
	\loigiai{
		\begin{enumerate}[a)]
			\item Ta có:
			\begin{align*}
				A &= 3x^3y \cdot 4xy^2 \\
				&= (3 \cdot 4) \cdot (x^3 \cdot x) \cdot (y \cdot y^2) \\
				&= 12x^4y^3
			\end{align*}
			\item Ta có:
			\begin{align*}
				B &= 3x^2y^3 \cdot 2x^2y - 11x^4y^4 \\
				&= 6x^4y^4 - 11x^4y^4 \\
				&= (6 - 11)x^4y^4 \\
				&= -5x^4y^4
			\end{align*}
			\item Ta có:
			\begin{align*}
				C &= 18x^6y^5 : 3x^2y^2 - 12x^4y^3 \\
				&= 6x^4y^3 - 12x^4y^3 \\
				&= (6 - 12)x^4y^3 \\
				&= -6x^4y^3
			\end{align*}
			\item Ta có:
			\begin{align*}
				D &= 4x^3y^2 + 8x^5y^3 : 2x^2y \\
				&= 4x^3y^2 + 4x^3y^2 \\
				&= 8x^3y^2
			\end{align*}
		\end{enumerate}
	}
\end{ex}

\newpage
% ----------------------------------------------------------------------
% TRANG 2 (MẶT SAU) -- DÒNG KẺ TIẾP THEO
% ----------------------------------------------------------------------
\noindent
\textbf{Học sinh:} Thế Anh \quad -- \quad \textbf{Mã đề:} \fbox{\textbf{009}} \hfill \textit{Trang 2/2}\par
\vspace{-0.1cm}
\noindent\hrulefill
\vspace{0.3cm}

% --- CÂU 2 (TRANG 2) ---
\begin{ex}
	Thực hiện phép tính:
	\begin{enumerate}[a)]
		\item $A = -5a^2b^3 \cdot 2ab^2$
		\item $B = 4a^3b \cdot 3ab^2 - 17a^4b^3$
		\item $C = 24a^7b^4 : 6a^4b^2 - 9a^3b^2$
		\item $D = 3a^2b^4 + 14a^4b^5 : 2a^2b$
	\end{enumerate}
	\loigiai{
		\begin{enumerate}[a)]
			\item Ta có:
			\begin{align*}
				A &= -5a^2b^3 \cdot 2ab^2 \\
				&= (-5 \cdot 2) \cdot (a^2 \cdot a) \cdot (b^3 \cdot b^2) \\
				&= -10a^3b^5
			\end{align*}
			\item Ta có:
			\begin{align*}
				B &= 4a^3b \cdot 3ab^2 - 17a^4b^3 \\
				&= 12a^4b^3 - 17a^4b^3 \\
				&= (12 - 17)a^4b^3 \\
				&= -5a^4b^3
			\end{align*}
			\item Ta có:
			\begin{align*}
				C &= 24a^7b^4 : 6a^4b^2 - 9a^3b^2 \\
				&= 4a^3b^2 - 9a^3b^2 \\
				&= (4 - 9)a^3b^2 \\
				&= -5a^3b^2
			\end{align*}
			\item Ta có:
			\begin{align*}
				D &= 3a^2b^4 + 14a^4b^5 : 2a^2b \\
				&= 3a^2b^4 + 7a^2b^4 \\
				&= 10a^2b^4
			\end{align*}
		\end{enumerate}
	}
\end{ex}

\newpage

% ======================================================================
% HỌC SINH: Thiện | MÃ: HS019 | MÃ ĐỀ: 002
% ======================================================================
\noindent
\begin{tcolorbox}[colback=white,colframe=black,arc=0mm,boxrule=0.8pt,top=4pt,bottom=4pt,left=6pt,right=6pt]
	\begin{minipage}[c]{0.44\textwidth}
		{\bfseries\large LỚP TOÁN THẦY THÔNG}\\[3pt]
		\textbf{Họ và tên:} \textbf{\large Thiện}\\[3pt]
		\textbf{Lớp:} LỚP 10 \quad -- \quad \textbf{Mã HS:} \texttt{HS019}
	\end{minipage}
	\hfill
	\begin{minipage}[c]{0.22\textwidth}
		\begin{tcolorbox}[colback=white,colframe=black,arc=0mm,boxrule=0.6pt,height=1.8cm,halign=center,valign=top,top=2pt]
			{\bfseries\normalsize Điểm}
		\end{tcolorbox}
	\end{minipage}
	\hfill
	\begin{minipage}[c]{0.30\textwidth}
		\raggedleft
		\textit{Thời gian làm bài: 30 phút}\\[4pt]
		\textbf{Mã đề thi:} \fbox{\textbf{\Large\strut\ 002\ }}
	\end{minipage}
\end{tcolorbox}
\vspace{0.25cm}

\noindent\textbf{\large ĐỀ BÀI (MÃ ĐỀ: 002)}
\vspace{0.1cm}

% --- CÂU 1 (TRANG 1) ---
\begin{ex}
	Thực hiện phép tính:
	\begin{enumerate}[a)]
		\item $A = 3x^3y \cdot 4xy^2$
		\item $B = 3x^2y^3 \cdot 2x^2y - 11x^4y^4$
		\item $C = 18x^6y^5 : 3x^2y^2 - 12x^4y^3$
		\item $D = 4x^3y^2 + 8x^5y^3 : 2x^2y$
	\end{enumerate}
	\loigiai{
		\begin{enumerate}[a)]
			\item Ta có:
			\begin{align*}
				A &= 3x^3y \cdot 4xy^2 \\
				&= (3 \cdot 4) \cdot (x^3 \cdot x) \cdot (y \cdot y^2) \\
				&= 12x^4y^3
			\end{align*}
			\item Ta có:
			\begin{align*}
				B &= 3x^2y^3 \cdot 2x^2y - 11x^4y^4 \\
				&= 6x^4y^4 - 11x^4y^4 \\
				&= (6 - 11)x^4y^4 \\
				&= -5x^4y^4
			\end{align*}
			\item Ta có:
			\begin{align*}
				C &= 18x^6y^5 : 3x^2y^2 - 12x^4y^3 \\
				&= 6x^4y^3 - 12x^4y^3 \\
				&= (6 - 12)x^4y^3 \\
				&= -6x^4y^3
			\end{align*}
			\item Ta có:
			\begin{align*}
				D &= 4x^3y^2 + 8x^5y^3 : 2x^2y \\
				&= 4x^3y^2 + 4x^3y^2 \\
				&= 8x^3y^2
			\end{align*}
		\end{enumerate}
	}
\end{ex}

\newpage
% ----------------------------------------------------------------------
% TRANG 2 (MẶT SAU) -- DÒNG KẺ TIẾP THEO
% ----------------------------------------------------------------------
\noindent
\textbf{Học sinh:} Thiện \quad -- \quad \textbf{Mã đề:} \fbox{\textbf{002}} \hfill \textit{Trang 2/2}\par
\vspace{-0.1cm}
\noindent\hrulefill
\vspace{0.3cm}

% --- CÂU 2 (TRANG 2) ---
\begin{ex}
	Thực hiện phép tính:
	\begin{enumerate}[a)]
		\item $A = -5a^2b^3 \cdot 2ab^2$
		\item $B = 4a^3b \cdot 3ab^2 - 17a^4b^3$
		\item $C = 24a^7b^4 : 6a^4b^2 - 9a^3b^2$
		\item $D = 3a^2b^4 + 14a^4b^5 : 2a^2b$
	\end{enumerate}
	\loigiai{
		\begin{enumerate}[a)]
			\item Ta có:
			\begin{align*}
				A &= -5a^2b^3 \cdot 2ab^2 \\
				&= (-5 \cdot 2) \cdot (a^2 \cdot a) \cdot (b^3 \cdot b^2) \\
				&= -10a^3b^5
			\end{align*}
			\item Ta có:
			\begin{align*}
				B &= 4a^3b \cdot 3ab^2 - 17a^4b^3 \\
				&= 12a^4b^3 - 17a^4b^3 \\
				&= (12 - 17)a^4b^3 \\
				&= -5a^4b^3
			\end{align*}
			\item Ta có:
			\begin{align*}
				C &= 24a^7b^4 : 6a^4b^2 - 9a^3b^2 \\
				&= 4a^3b^2 - 9a^3b^2 \\
				&= (4 - 9)a^3b^2 \\
				&= -5a^3b^2
			\end{align*}
			\item Ta có:
			\begin{align*}
				D &= 3a^2b^4 + 14a^4b^5 : 2a^2b \\
				&= 3a^2b^4 + 7a^2b^4 \\
				&= 10a^2b^4
			\end{align*}
		\end{enumerate}
	}
\end{ex}

\newpage

% ======================================================================
% HỌC SINH: Thục Linh | MÃ: HS004 | MÃ ĐỀ: 004
% ======================================================================
\noindent
\begin{tcolorbox}[colback=white,colframe=black,arc=0mm,boxrule=0.8pt,top=4pt,bottom=4pt,left=6pt,right=6pt]
	\begin{minipage}[c]{0.44\textwidth}
		{\bfseries\large LỚP TOÁN THẦY THÔNG}\\[3pt]
		\textbf{Họ và tên:} \textbf{\large Thục Linh}\\[3pt]
		\textbf{Lớp:} LỚP 10 \quad -- \quad \textbf{Mã HS:} \texttt{HS004}
	\end{minipage}
	\hfill
	\begin{minipage}[c]{0.22\textwidth}
		\begin{tcolorbox}[colback=white,colframe=black,arc=0mm,boxrule=0.6pt,height=1.8cm,halign=center,valign=top,top=2pt]
			{\bfseries\normalsize Điểm}
		\end{tcolorbox}
	\end{minipage}
	\hfill
	\begin{minipage}[c]{0.30\textwidth}
		\raggedleft
		\textit{Thời gian làm bài: 30 phút}\\[4pt]
		\textbf{Mã đề thi:} \fbox{\textbf{\Large\strut\ 004\ }}
	\end{minipage}
\end{tcolorbox}
\vspace{0.25cm}

\noindent\textbf{\large ĐỀ BÀI (MÃ ĐỀ: 004)}
\vspace{0.1cm}

% --- CÂU 1 (TRANG 1) ---
\begin{ex}
	Thực hiện phép tính:
	\begin{enumerate}[a)]
		\item $A = 3x^3y \cdot 4xy^2$
		\item $B = 3x^2y^3 \cdot 2x^2y - 11x^4y^4$
		\item $C = 18x^6y^5 : 3x^2y^2 - 12x^4y^3$
		\item $D = 4x^3y^2 + 8x^5y^3 : 2x^2y$
	\end{enumerate}
	\loigiai{
		\begin{enumerate}[a)]
			\item Ta có:
			\begin{align*}
				A &= 3x^3y \cdot 4xy^2 \\
				&= (3 \cdot 4) \cdot (x^3 \cdot x) \cdot (y \cdot y^2) \\
				&= 12x^4y^3
			\end{align*}
			\item Ta có:
			\begin{align*}
				B &= 3x^2y^3 \cdot 2x^2y - 11x^4y^4 \\
				&= 6x^4y^4 - 11x^4y^4 \\
				&= (6 - 11)x^4y^4 \\
				&= -5x^4y^4
			\end{align*}
			\item Ta có:
			\begin{align*}
				C &= 18x^6y^5 : 3x^2y^2 - 12x^4y^3 \\
				&= 6x^4y^3 - 12x^4y^3 \\
				&= (6 - 12)x^4y^3 \\
				&= -6x^4y^3
			\end{align*}
			\item Ta có:
			\begin{align*}
				D &= 4x^3y^2 + 8x^5y^3 : 2x^2y \\
				&= 4x^3y^2 + 4x^3y^2 \\
				&= 8x^3y^2
			\end{align*}
		\end{enumerate}
	}
\end{ex}

\newpage
% ----------------------------------------------------------------------
% TRANG 2 (MẶT SAU) -- DÒNG KẺ TIẾP THEO
% ----------------------------------------------------------------------
\noindent
\textbf{Học sinh:} Thục Linh \quad -- \quad \textbf{Mã đề:} \fbox{\textbf{004}} \hfill \textit{Trang 2/2}\par
\vspace{-0.1cm}
\noindent\hrulefill
\vspace{0.3cm}

% --- CÂU 2 (TRANG 2) ---
\begin{ex}
	Thực hiện phép tính:
	\begin{enumerate}[a)]
		\item $A = -5a^2b^3 \cdot 2ab^2$
		\item $B = 4a^3b \cdot 3ab^2 - 17a^4b^3$
		\item $C = 24a^7b^4 : 6a^4b^2 - 9a^3b^2$
		\item $D = 3a^2b^4 + 14a^4b^5 : 2a^2b$
	\end{enumerate}
	\loigiai{
		\begin{enumerate}[a)]
			\item Ta có:
			\begin{align*}
				A &= -5a^2b^3 \cdot 2ab^2 \\
				&= (-5 \cdot 2) \cdot (a^2 \cdot a) \cdot (b^3 \cdot b^2) \\
				&= -10a^3b^5
			\end{align*}
			\item Ta có:
			\begin{align*}
				B &= 4a^3b \cdot 3ab^2 - 17a^4b^3 \\
				&= 12a^4b^3 - 17a^4b^3 \\
				&= (12 - 17)a^4b^3 \\
				&= -5a^4b^3
			\end{align*}
			\item Ta có:
			\begin{align*}
				C &= 24a^7b^4 : 6a^4b^2 - 9a^3b^2 \\
				&= 4a^3b^2 - 9a^3b^2 \\
				&= (4 - 9)a^3b^2 \\
				&= -5a^3b^2
			\end{align*}
			\item Ta có:
			\begin{align*}
				D &= 3a^2b^4 + 14a^4b^5 : 2a^2b \\
				&= 3a^2b^4 + 7a^2b^4 \\
				&= 10a^2b^4
			\end{align*}
		\end{enumerate}
	}
\end{ex}

\newpage

% ======================================================================
% HỌC SINH: Thử | MÃ: HS1707 | MÃ ĐỀ: 003
% ======================================================================
\noindent
\begin{tcolorbox}[colback=white,colframe=black,arc=0mm,boxrule=0.8pt,top=4pt,bottom=4pt,left=6pt,right=6pt]
	\begin{minipage}[c]{0.44\textwidth}
		{\bfseries\large LỚP TOÁN THẦY THÔNG}\\[3pt]
		\textbf{Họ và tên:} \textbf{\large Thử}\\[3pt]
		\textbf{Lớp:} LỚP 10 \quad -- \quad \textbf{Mã HS:} \texttt{HS1707}
	\end{minipage}
	\hfill
	\begin{minipage}[c]{0.22\textwidth}
		\begin{tcolorbox}[colback=white,colframe=black,arc=0mm,boxrule=0.6pt,height=1.8cm,halign=center,valign=top,top=2pt]
			{\bfseries\normalsize Điểm}
		\end{tcolorbox}
	\end{minipage}
	\hfill
	\begin{minipage}[c]{0.30\textwidth}
		\raggedleft
		\textit{Thời gian làm bài: 30 phút}\\[4pt]
		\textbf{Mã đề thi:} \fbox{\textbf{\Large\strut\ 003\ }}
	\end{minipage}
\end{tcolorbox}
\vspace{0.25cm}

\noindent\textbf{\large ĐỀ BÀI (MÃ ĐỀ: 003)}
\vspace{0.1cm}

% --- CÂU 1 (TRANG 1) ---
\begin{ex}
	Thực hiện phép tính:
	\begin{enumerate}[a)]
		\item $A = 3x^3y \cdot 4xy^2$
		\item $B = 3x^2y^3 \cdot 2x^2y - 11x^4y^4$
		\item $C = 18x^6y^5 : 3x^2y^2 - 12x^4y^3$
		\item $D = 4x^3y^2 + 8x^5y^3 : 2x^2y$
	\end{enumerate}
	\loigiai{
		\begin{enumerate}[a)]
			\item Ta có:
			\begin{align*}
				A &= 3x^3y \cdot 4xy^2 \\
				&= (3 \cdot 4) \cdot (x^3 \cdot x) \cdot (y \cdot y^2) \\
				&= 12x^4y^3
			\end{align*}
			\item Ta có:
			\begin{align*}
				B &= 3x^2y^3 \cdot 2x^2y - 11x^4y^4 \\
				&= 6x^4y^4 - 11x^4y^4 \\
				&= (6 - 11)x^4y^4 \\
				&= -5x^4y^4
			\end{align*}
			\item Ta có:
			\begin{align*}
				C &= 18x^6y^5 : 3x^2y^2 - 12x^4y^3 \\
				&= 6x^4y^3 - 12x^4y^3 \\
				&= (6 - 12)x^4y^3 \\
				&= -6x^4y^3
			\end{align*}
			\item Ta có:
			\begin{align*}
				D &= 4x^3y^2 + 8x^5y^3 : 2x^2y \\
				&= 4x^3y^2 + 4x^3y^2 \\
				&= 8x^3y^2
			\end{align*}
		\end{enumerate}
	}
\end{ex}

\newpage
% ----------------------------------------------------------------------
% TRANG 2 (MẶT SAU) -- DÒNG KẺ TIẾP THEO
% ----------------------------------------------------------------------
\noindent
\textbf{Học sinh:} Thử \quad -- \quad \textbf{Mã đề:} \fbox{\textbf{003}} \hfill \textit{Trang 2/2}\par
\vspace{-0.1cm}
\noindent\hrulefill
\vspace{0.3cm}

% --- CÂU 2 (TRANG 2) ---
\begin{ex}
	Thực hiện phép tính:
	\begin{enumerate}[a)]
		\item $A = -5a^2b^3 \cdot 2ab^2$
		\item $B = 4a^3b \cdot 3ab^2 - 17a^4b^3$
		\item $C = 24a^7b^4 : 6a^4b^2 - 9a^3b^2$
		\item $D = 3a^2b^4 + 14a^4b^5 : 2a^2b$
	\end{enumerate}
	\loigiai{
		\begin{enumerate}[a)]
			\item Ta có:
			\begin{align*}
				A &= -5a^2b^3 \cdot 2ab^2 \\
				&= (-5 \cdot 2) \cdot (a^2 \cdot a) \cdot (b^3 \cdot b^2) \\
				&= -10a^3b^5
			\end{align*}
			\item Ta có:
			\begin{align*}
				B &= 4a^3b \cdot 3ab^2 - 17a^4b^3 \\
				&= 12a^4b^3 - 17a^4b^3 \\
				&= (12 - 17)a^4b^3 \\
				&= -5a^4b^3
			\end{align*}
			\item Ta có:
			\begin{align*}
				C &= 24a^7b^4 : 6a^4b^2 - 9a^3b^2 \\
				&= 4a^3b^2 - 9a^3b^2 \\
				&= (4 - 9)a^3b^2 \\
				&= -5a^3b^2
			\end{align*}
			\item Ta có:
			\begin{align*}
				D &= 3a^2b^4 + 14a^4b^5 : 2a^2b \\
				&= 3a^2b^4 + 7a^2b^4 \\
				&= 10a^2b^4
			\end{align*}
		\end{enumerate}
	}
\end{ex}

\newpage

% ======================================================================
% HỌC SINH: Tiến Đạt | MÃ: HS015 | MÃ ĐỀ: 008
% ======================================================================
\noindent
\begin{tcolorbox}[colback=white,colframe=black,arc=0mm,boxrule=0.8pt,top=4pt,bottom=4pt,left=6pt,right=6pt]
	\begin{minipage}[c]{0.44\textwidth}
		{\bfseries\large LỚP TOÁN THẦY THÔNG}\\[3pt]
		\textbf{Họ và tên:} \textbf{\large Tiến Đạt}\\[3pt]
		\textbf{Lớp:} LỚP 10 \quad -- \quad \textbf{Mã HS:} \texttt{HS015}
	\end{minipage}
	\hfill
	\begin{minipage}[c]{0.22\textwidth}
		\begin{tcolorbox}[colback=white,colframe=black,arc=0mm,boxrule=0.6pt,height=1.8cm,halign=center,valign=top,top=2pt]
			{\bfseries\normalsize Điểm}
		\end{tcolorbox}
	\end{minipage}
	\hfill
	\begin{minipage}[c]{0.30\textwidth}
		\raggedleft
		\textit{Thời gian làm bài: 30 phút}\\[4pt]
		\textbf{Mã đề thi:} \fbox{\textbf{\Large\strut\ 008\ }}
	\end{minipage}
\end{tcolorbox}
\vspace{0.25cm}

\noindent\textbf{\large ĐỀ BÀI (MÃ ĐỀ: 008)}
\vspace{0.1cm}

% --- CÂU 1 (TRANG 1) ---
\begin{ex}
	Thực hiện phép tính:
	\begin{enumerate}[a)]
		\item $A = 3x^3y \cdot 4xy^2$
		\item $B = 3x^2y^3 \cdot 2x^2y - 11x^4y^4$
		\item $C = 18x^6y^5 : 3x^2y^2 - 12x^4y^3$
		\item $D = 4x^3y^2 + 8x^5y^3 : 2x^2y$
	\end{enumerate}
	\loigiai{
		\begin{enumerate}[a)]
			\item Ta có:
			\begin{align*}
				A &= 3x^3y \cdot 4xy^2 \\
				&= (3 \cdot 4) \cdot (x^3 \cdot x) \cdot (y \cdot y^2) \\
				&= 12x^4y^3
			\end{align*}
			\item Ta có:
			\begin{align*}
				B &= 3x^2y^3 \cdot 2x^2y - 11x^4y^4 \\
				&= 6x^4y^4 - 11x^4y^4 \\
				&= (6 - 11)x^4y^4 \\
				&= -5x^4y^4
			\end{align*}
			\item Ta có:
			\begin{align*}
				C &= 18x^6y^5 : 3x^2y^2 - 12x^4y^3 \\
				&= 6x^4y^3 - 12x^4y^3 \\
				&= (6 - 12)x^4y^3 \\
				&= -6x^4y^3
			\end{align*}
			\item Ta có:
			\begin{align*}
				D &= 4x^3y^2 + 8x^5y^3 : 2x^2y \\
				&= 4x^3y^2 + 4x^3y^2 \\
				&= 8x^3y^2
			\end{align*}
		\end{enumerate}
	}
\end{ex}

\newpage
% ----------------------------------------------------------------------
% TRANG 2 (MẶT SAU) -- DÒNG KẺ TIẾP THEO
% ----------------------------------------------------------------------
\noindent
\textbf{Học sinh:} Tiến Đạt \quad -- \quad \textbf{Mã đề:} \fbox{\textbf{008}} \hfill \textit{Trang 2/2}\par
\vspace{-0.1cm}
\noindent\hrulefill
\vspace{0.3cm}

% --- CÂU 2 (TRANG 2) ---
\begin{ex}
	Thực hiện phép tính:
	\begin{enumerate}[a)]
		\item $A = -5a^2b^3 \cdot 2ab^2$
		\item $B = 4a^3b \cdot 3ab^2 - 17a^4b^3$
		\item $C = 24a^7b^4 : 6a^4b^2 - 9a^3b^2$
		\item $D = 3a^2b^4 + 14a^4b^5 : 2a^2b$
	\end{enumerate}
	\loigiai{
		\begin{enumerate}[a)]
			\item Ta có:
			\begin{align*}
				A &= -5a^2b^3 \cdot 2ab^2 \\
				&= (-5 \cdot 2) \cdot (a^2 \cdot a) \cdot (b^3 \cdot b^2) \\
				&= -10a^3b^5
			\end{align*}
			\item Ta có:
			\begin{align*}
				B &= 4a^3b \cdot 3ab^2 - 17a^4b^3 \\
				&= 12a^4b^3 - 17a^4b^3 \\
				&= (12 - 17)a^4b^3 \\
				&= -5a^4b^3
			\end{align*}
			\item Ta có:
			\begin{align*}
				C &= 24a^7b^4 : 6a^4b^2 - 9a^3b^2 \\
				&= 4a^3b^2 - 9a^3b^2 \\
				&= (4 - 9)a^3b^2 \\
				&= -5a^3b^2
			\end{align*}
			\item Ta có:
			\begin{align*}
				D &= 3a^2b^4 + 14a^4b^5 : 2a^2b \\
				&= 3a^2b^4 + 7a^2b^4 \\
				&= 10a^2b^4
			\end{align*}
		\end{enumerate}
	}
\end{ex}

\newpage

% ======================================================================
% HỌC SINH: Tiến Quang | MÃ: HS018 | MÃ ĐỀ: 007
% ======================================================================
\noindent
\begin{tcolorbox}[colback=white,colframe=black,arc=0mm,boxrule=0.8pt,top=4pt,bottom=4pt,left=6pt,right=6pt]
	\begin{minipage}[c]{0.44\textwidth}
		{\bfseries\large LỚP TOÁN THẦY THÔNG}\\[3pt]
		\textbf{Họ và tên:} \textbf{\large Tiến Quang}\\[3pt]
		\textbf{Lớp:} LỚP 10 \quad -- \quad \textbf{Mã HS:} \texttt{HS018}
	\end{minipage}
	\hfill
	\begin{minipage}[c]{0.22\textwidth}
		\begin{tcolorbox}[colback=white,colframe=black,arc=0mm,boxrule=0.6pt,height=1.8cm,halign=center,valign=top,top=2pt]
			{\bfseries\normalsize Điểm}
		\end{tcolorbox}
	\end{minipage}
	\hfill
	\begin{minipage}[c]{0.30\textwidth}
		\raggedleft
		\textit{Thời gian làm bài: 30 phút}\\[4pt]
		\textbf{Mã đề thi:} \fbox{\textbf{\Large\strut\ 007\ }}
	\end{minipage}
\end{tcolorbox}
\vspace{0.25cm}

\noindent\textbf{\large ĐỀ BÀI (MÃ ĐỀ: 007)}
\vspace{0.1cm}

% --- CÂU 1 (TRANG 1) ---
\begin{ex}
	Thực hiện phép tính:
	\begin{enumerate}[a)]
		\item $A = 3x^3y \cdot 4xy^2$
		\item $B = 3x^2y^3 \cdot 2x^2y - 11x^4y^4$
		\item $C = 18x^6y^5 : 3x^2y^2 - 12x^4y^3$
		\item $D = 4x^3y^2 + 8x^5y^3 : 2x^2y$
	\end{enumerate}
	\loigiai{
		\begin{enumerate}[a)]
			\item Ta có:
			\begin{align*}
				A &= 3x^3y \cdot 4xy^2 \\
				&= (3 \cdot 4) \cdot (x^3 \cdot x) \cdot (y \cdot y^2) \\
				&= 12x^4y^3
			\end{align*}
			\item Ta có:
			\begin{align*}
				B &= 3x^2y^3 \cdot 2x^2y - 11x^4y^4 \\
				&= 6x^4y^4 - 11x^4y^4 \\
				&= (6 - 11)x^4y^4 \\
				&= -5x^4y^4
			\end{align*}
			\item Ta có:
			\begin{align*}
				C &= 18x^6y^5 : 3x^2y^2 - 12x^4y^3 \\
				&= 6x^4y^3 - 12x^4y^3 \\
				&= (6 - 12)x^4y^3 \\
				&= -6x^4y^3
			\end{align*}
			\item Ta có:
			\begin{align*}
				D &= 4x^3y^2 + 8x^5y^3 : 2x^2y \\
				&= 4x^3y^2 + 4x^3y^2 \\
				&= 8x^3y^2
			\end{align*}
		\end{enumerate}
	}
\end{ex}

\newpage
% ----------------------------------------------------------------------
% TRANG 2 (MẶT SAU) -- DÒNG KẺ TIẾP THEO
% ----------------------------------------------------------------------
\noindent
\textbf{Học sinh:} Tiến Quang \quad -- \quad \textbf{Mã đề:} \fbox{\textbf{007}} \hfill \textit{Trang 2/2}\par
\vspace{-0.1cm}
\noindent\hrulefill
\vspace{0.3cm}

% --- CÂU 2 (TRANG 2) ---
\begin{ex}
	Thực hiện phép tính:
	\begin{enumerate}[a)]
		\item $A = -5a^2b^3 \cdot 2ab^2$
		\item $B = 4a^3b \cdot 3ab^2 - 17a^4b^3$
		\item $C = 24a^7b^4 : 6a^4b^2 - 9a^3b^2$
		\item $D = 3a^2b^4 + 14a^4b^5 : 2a^2b$
	\end{enumerate}
	\loigiai{
		\begin{enumerate}[a)]
			\item Ta có:
			\begin{align*}
				A &= -5a^2b^3 \cdot 2ab^2 \\
				&= (-5 \cdot 2) \cdot (a^2 \cdot a) \cdot (b^3 \cdot b^2) \\
				&= -10a^3b^5
			\end{align*}
			\item Ta có:
			\begin{align*}
				B &= 4a^3b \cdot 3ab^2 - 17a^4b^3 \\
				&= 12a^4b^3 - 17a^4b^3 \\
				&= (12 - 17)a^4b^3 \\
				&= -5a^4b^3
			\end{align*}
			\item Ta có:
			\begin{align*}
				C &= 24a^7b^4 : 6a^4b^2 - 9a^3b^2 \\
				&= 4a^3b^2 - 9a^3b^2 \\
				&= (4 - 9)a^3b^2 \\
				&= -5a^3b^2
			\end{align*}
			\item Ta có:
			\begin{align*}
				D &= 3a^2b^4 + 14a^4b^5 : 2a^2b \\
				&= 3a^2b^4 + 7a^2b^4 \\
				&= 10a^2b^4
			\end{align*}
		\end{enumerate}
	}
\end{ex}

\newpage

% ======================================================================
% HỌC SINH: Trần Khánh Chí | MÃ: HS030 | MÃ ĐỀ: 005
% ======================================================================
\noindent
\begin{tcolorbox}[colback=white,colframe=black,arc=0mm,boxrule=0.8pt,top=4pt,bottom=4pt,left=6pt,right=6pt]
	\begin{minipage}[c]{0.44\textwidth}
		{\bfseries\large LỚP TOÁN THẦY THÔNG}\\[3pt]
		\textbf{Họ và tên:} \textbf{\large Trần Khánh Chí}\\[3pt]
		\textbf{Lớp:} LỚP 10 \quad -- \quad \textbf{Mã HS:} \texttt{HS030}
	\end{minipage}
	\hfill
	\begin{minipage}[c]{0.22\textwidth}
		\begin{tcolorbox}[colback=white,colframe=black,arc=0mm,boxrule=0.6pt,height=1.8cm,halign=center,valign=top,top=2pt]
			{\bfseries\normalsize Điểm}
		\end{tcolorbox}
	\end{minipage}
	\hfill
	\begin{minipage}[c]{0.30\textwidth}
		\raggedleft
		\textit{Thời gian làm bài: 30 phút}\\[4pt]
		\textbf{Mã đề thi:} \fbox{\textbf{\Large\strut\ 005\ }}
	\end{minipage}
\end{tcolorbox}
\vspace{0.25cm}

\noindent\textbf{\large ĐỀ BÀI (MÃ ĐỀ: 005)}
\vspace{0.1cm}

% --- CÂU 1 (TRANG 1) ---
\begin{ex}
	Thực hiện phép tính:
	\begin{enumerate}[a)]
		\item $A = 3x^3y \cdot 4xy^2$
		\item $B = 3x^2y^3 \cdot 2x^2y - 11x^4y^4$
		\item $C = 18x^6y^5 : 3x^2y^2 - 12x^4y^3$
		\item $D = 4x^3y^2 + 8x^5y^3 : 2x^2y$
	\end{enumerate}
	\loigiai{
		\begin{enumerate}[a)]
			\item Ta có:
			\begin{align*}
				A &= 3x^3y \cdot 4xy^2 \\
				&= (3 \cdot 4) \cdot (x^3 \cdot x) \cdot (y \cdot y^2) \\
				&= 12x^4y^3
			\end{align*}
			\item Ta có:
			\begin{align*}
				B &= 3x^2y^3 \cdot 2x^2y - 11x^4y^4 \\
				&= 6x^4y^4 - 11x^4y^4 \\
				&= (6 - 11)x^4y^4 \\
				&= -5x^4y^4
			\end{align*}
			\item Ta có:
			\begin{align*}
				C &= 18x^6y^5 : 3x^2y^2 - 12x^4y^3 \\
				&= 6x^4y^3 - 12x^4y^3 \\
				&= (6 - 12)x^4y^3 \\
				&= -6x^4y^3
			\end{align*}
			\item Ta có:
			\begin{align*}
				D &= 4x^3y^2 + 8x^5y^3 : 2x^2y \\
				&= 4x^3y^2 + 4x^3y^2 \\
				&= 8x^3y^2
			\end{align*}
		\end{enumerate}
	}
\end{ex}

\newpage
% ----------------------------------------------------------------------
% TRANG 2 (MẶT SAU) -- DÒNG KẺ TIẾP THEO
% ----------------------------------------------------------------------
\noindent
\textbf{Học sinh:} Trần Khánh Chí \quad -- \quad \textbf{Mã đề:} \fbox{\textbf{005}} \hfill \textit{Trang 2/2}\par
\vspace{-0.1cm}
\noindent\hrulefill
\vspace{0.3cm}

% --- CÂU 2 (TRANG 2) ---
\begin{ex}
	Thực hiện phép tính:
	\begin{enumerate}[a)]
		\item $A = -5a^2b^3 \cdot 2ab^2$
		\item $B = 4a^3b \cdot 3ab^2 - 17a^4b^3$
		\item $C = 24a^7b^4 : 6a^4b^2 - 9a^3b^2$
		\item $D = 3a^2b^4 + 14a^4b^5 : 2a^2b$
	\end{enumerate}
	\loigiai{
		\begin{enumerate}[a)]
			\item Ta có:
			\begin{align*}
				A &= -5a^2b^3 \cdot 2ab^2 \\
				&= (-5 \cdot 2) \cdot (a^2 \cdot a) \cdot (b^3 \cdot b^2) \\
				&= -10a^3b^5
			\end{align*}
			\item Ta có:
			\begin{align*}
				B &= 4a^3b \cdot 3ab^2 - 17a^4b^3 \\
				&= 12a^4b^3 - 17a^4b^3 \\
				&= (12 - 17)a^4b^3 \\
				&= -5a^4b^3
			\end{align*}
			\item Ta có:
			\begin{align*}
				C &= 24a^7b^4 : 6a^4b^2 - 9a^3b^2 \\
				&= 4a^3b^2 - 9a^3b^2 \\
				&= (4 - 9)a^3b^2 \\
				&= -5a^3b^2
			\end{align*}
			\item Ta có:
			\begin{align*}
				D &= 3a^2b^4 + 14a^4b^5 : 2a^2b \\
				&= 3a^2b^4 + 7a^2b^4 \\
				&= 10a^2b^4
			\end{align*}
		\end{enumerate}
	}
\end{ex}

\newpage

% ======================================================================
% HỌC SINH: Trọng Huy | MÃ: HS009 | MÃ ĐỀ: 010
% ======================================================================
\noindent
\begin{tcolorbox}[colback=white,colframe=black,arc=0mm,boxrule=0.8pt,top=4pt,bottom=4pt,left=6pt,right=6pt]
	\begin{minipage}[c]{0.44\textwidth}
		{\bfseries\large LỚP TOÁN THẦY THÔNG}\\[3pt]
		\textbf{Họ và tên:} \textbf{\large Trọng Huy}\\[3pt]
		\textbf{Lớp:} LỚP 10 \quad -- \quad \textbf{Mã HS:} \texttt{HS009}
	\end{minipage}
	\hfill
	\begin{minipage}[c]{0.22\textwidth}
		\begin{tcolorbox}[colback=white,colframe=black,arc=0mm,boxrule=0.6pt,height=1.8cm,halign=center,valign=top,top=2pt]
			{\bfseries\normalsize Điểm}
		\end{tcolorbox}
	\end{minipage}
	\hfill
	\begin{minipage}[c]{0.30\textwidth}
		\raggedleft
		\textit{Thời gian làm bài: 30 phút}\\[4pt]
		\textbf{Mã đề thi:} \fbox{\textbf{\Large\strut\ 010\ }}
	\end{minipage}
\end{tcolorbox}
\vspace{0.25cm}

\noindent\textbf{\large ĐỀ BÀI (MÃ ĐỀ: 010)}
\vspace{0.1cm}

% --- CÂU 1 (TRANG 1) ---
\begin{ex}
	Thực hiện phép tính:
	\begin{enumerate}[a)]
		\item $A = 3x^3y \cdot 4xy^2$
		\item $B = 3x^2y^3 \cdot 2x^2y - 11x^4y^4$
		\item $C = 18x^6y^5 : 3x^2y^2 - 12x^4y^3$
		\item $D = 4x^3y^2 + 8x^5y^3 : 2x^2y$
	\end{enumerate}
	\loigiai{
		\begin{enumerate}[a)]
			\item Ta có:
			\begin{align*}
				A &= 3x^3y \cdot 4xy^2 \\
				&= (3 \cdot 4) \cdot (x^3 \cdot x) \cdot (y \cdot y^2) \\
				&= 12x^4y^3
			\end{align*}
			\item Ta có:
			\begin{align*}
				B &= 3x^2y^3 \cdot 2x^2y - 11x^4y^4 \\
				&= 6x^4y^4 - 11x^4y^4 \\
				&= (6 - 11)x^4y^4 \\
				&= -5x^4y^4
			\end{align*}
			\item Ta có:
			\begin{align*}
				C &= 18x^6y^5 : 3x^2y^2 - 12x^4y^3 \\
				&= 6x^4y^3 - 12x^4y^3 \\
				&= (6 - 12)x^4y^3 \\
				&= -6x^4y^3
			\end{align*}
			\item Ta có:
			\begin{align*}
				D &= 4x^3y^2 + 8x^5y^3 : 2x^2y \\
				&= 4x^3y^2 + 4x^3y^2 \\
				&= 8x^3y^2
			\end{align*}
		\end{enumerate}
	}
\end{ex}

\newpage
% ----------------------------------------------------------------------
% TRANG 2 (MẶT SAU) -- DÒNG KẺ TIẾP THEO
% ----------------------------------------------------------------------
\noindent
\textbf{Học sinh:} Trọng Huy \quad -- \quad \textbf{Mã đề:} \fbox{\textbf{010}} \hfill \textit{Trang 2/2}\par
\vspace{-0.1cm}
\noindent\hrulefill
\vspace{0.3cm}

% --- CÂU 2 (TRANG 2) ---
\begin{ex}
	Thực hiện phép tính:
	\begin{enumerate}[a)]
		\item $A = -5a^2b^3 \cdot 2ab^2$
		\item $B = 4a^3b \cdot 3ab^2 - 17a^4b^3$
		\item $C = 24a^7b^4 : 6a^4b^2 - 9a^3b^2$
		\item $D = 3a^2b^4 + 14a^4b^5 : 2a^2b$
	\end{enumerate}
	\loigiai{
		\begin{enumerate}[a)]
			\item Ta có:
			\begin{align*}
				A &= -5a^2b^3 \cdot 2ab^2 \\
				&= (-5 \cdot 2) \cdot (a^2 \cdot a) \cdot (b^3 \cdot b^2) \\
				&= -10a^3b^5
			\end{align*}
			\item Ta có:
			\begin{align*}
				B &= 4a^3b \cdot 3ab^2 - 17a^4b^3 \\
				&= 12a^4b^3 - 17a^4b^3 \\
				&= (12 - 17)a^4b^3 \\
				&= -5a^4b^3
			\end{align*}
			\item Ta có:
			\begin{align*}
				C &= 24a^7b^4 : 6a^4b^2 - 9a^3b^2 \\
				&= 4a^3b^2 - 9a^3b^2 \\
				&= (4 - 9)a^3b^2 \\
				&= -5a^3b^2
			\end{align*}
			\item Ta có:
			\begin{align*}
				D &= 3a^2b^4 + 14a^4b^5 : 2a^2b \\
				&= 3a^2b^4 + 7a^2b^4 \\
				&= 10a^2b^4
			\end{align*}
		\end{enumerate}
	}
\end{ex}

\newpage

% ======================================================================
% HỌC SINH: Trung Nguyên | MÃ: HS032 | MÃ ĐỀ: 008
% ======================================================================
\noindent
\begin{tcolorbox}[colback=white,colframe=black,arc=0mm,boxrule=0.8pt,top=4pt,bottom=4pt,left=6pt,right=6pt]
	\begin{minipage}[c]{0.44\textwidth}
		{\bfseries\large LỚP TOÁN THẦY THÔNG}\\[3pt]
		\textbf{Họ và tên:} \textbf{\large Trung Nguyên}\\[3pt]
		\textbf{Lớp:} LỚP 10 \quad -- \quad \textbf{Mã HS:} \texttt{HS032}
	\end{minipage}
	\hfill
	\begin{minipage}[c]{0.22\textwidth}
		\begin{tcolorbox}[colback=white,colframe=black,arc=0mm,boxrule=0.6pt,height=1.8cm,halign=center,valign=top,top=2pt]
			{\bfseries\normalsize Điểm}
		\end{tcolorbox}
	\end{minipage}
	\hfill
	\begin{minipage}[c]{0.30\textwidth}
		\raggedleft
		\textit{Thời gian làm bài: 30 phút}\\[4pt]
		\textbf{Mã đề thi:} \fbox{\textbf{\Large\strut\ 008\ }}
	\end{minipage}
\end{tcolorbox}
\vspace{0.25cm}

\noindent\textbf{\large ĐỀ BÀI (MÃ ĐỀ: 008)}
\vspace{0.1cm}

% --- CÂU 1 (TRANG 1) ---
\begin{ex}
	Thực hiện phép tính:
	\begin{enumerate}[a)]
		\item $A = 3x^3y \cdot 4xy^2$
		\item $B = 3x^2y^3 \cdot 2x^2y - 11x^4y^4$
		\item $C = 18x^6y^5 : 3x^2y^2 - 12x^4y^3$
		\item $D = 4x^3y^2 + 8x^5y^3 : 2x^2y$
	\end{enumerate}
	\loigiai{
		\begin{enumerate}[a)]
			\item Ta có:
			\begin{align*}
				A &= 3x^3y \cdot 4xy^2 \\
				&= (3 \cdot 4) \cdot (x^3 \cdot x) \cdot (y \cdot y^2) \\
				&= 12x^4y^3
			\end{align*}
			\item Ta có:
			\begin{align*}
				B &= 3x^2y^3 \cdot 2x^2y - 11x^4y^4 \\
				&= 6x^4y^4 - 11x^4y^4 \\
				&= (6 - 11)x^4y^4 \\
				&= -5x^4y^4
			\end{align*}
			\item Ta có:
			\begin{align*}
				C &= 18x^6y^5 : 3x^2y^2 - 12x^4y^3 \\
				&= 6x^4y^3 - 12x^4y^3 \\
				&= (6 - 12)x^4y^3 \\
				&= -6x^4y^3
			\end{align*}
			\item Ta có:
			\begin{align*}
				D &= 4x^3y^2 + 8x^5y^3 : 2x^2y \\
				&= 4x^3y^2 + 4x^3y^2 \\
				&= 8x^3y^2
			\end{align*}
		\end{enumerate}
	}
\end{ex}

\newpage
% ----------------------------------------------------------------------
% TRANG 2 (MẶT SAU) -- DÒNG KẺ TIẾP THEO
% ----------------------------------------------------------------------
\noindent
\textbf{Học sinh:} Trung Nguyên \quad -- \quad \textbf{Mã đề:} \fbox{\textbf{008}} \hfill \textit{Trang 2/2}\par
\vspace{-0.1cm}
\noindent\hrulefill
\vspace{0.3cm}

% --- CÂU 2 (TRANG 2) ---
\begin{ex}
	Thực hiện phép tính:
	\begin{enumerate}[a)]
		\item $A = -5a^2b^3 \cdot 2ab^2$
		\item $B = 4a^3b \cdot 3ab^2 - 17a^4b^3$
		\item $C = 24a^7b^4 : 6a^4b^2 - 9a^3b^2$
		\item $D = 3a^2b^4 + 14a^4b^5 : 2a^2b$
	\end{enumerate}
	\loigiai{
		\begin{enumerate}[a)]
			\item Ta có:
			\begin{align*}
				A &= -5a^2b^3 \cdot 2ab^2 \\
				&= (-5 \cdot 2) \cdot (a^2 \cdot a) \cdot (b^3 \cdot b^2) \\
				&= -10a^3b^5
			\end{align*}
			\item Ta có:
			\begin{align*}
				B &= 4a^3b \cdot 3ab^2 - 17a^4b^3 \\
				&= 12a^4b^3 - 17a^4b^3 \\
				&= (12 - 17)a^4b^3 \\
				&= -5a^4b^3
			\end{align*}
			\item Ta có:
			\begin{align*}
				C &= 24a^7b^4 : 6a^4b^2 - 9a^3b^2 \\
				&= 4a^3b^2 - 9a^3b^2 \\
				&= (4 - 9)a^3b^2 \\
				&= -5a^3b^2
			\end{align*}
			\item Ta có:
			\begin{align*}
				D &= 3a^2b^4 + 14a^4b^5 : 2a^2b \\
				&= 3a^2b^4 + 7a^2b^4 \\
				&= 10a^2b^4
			\end{align*}
		\end{enumerate}
	}
\end{ex}

\newpage

% ======================================================================
% HỌC SINH: Việt | MÃ: HS013 | MÃ ĐỀ: 001
% ======================================================================
\noindent
\begin{tcolorbox}[colback=white,colframe=black,arc=0mm,boxrule=0.8pt,top=4pt,bottom=4pt,left=6pt,right=6pt]
	\begin{minipage}[c]{0.44\textwidth}
		{\bfseries\large LỚP TOÁN THẦY THÔNG}\\[3pt]
		\textbf{Họ và tên:} \textbf{\large Việt}\\[3pt]
		\textbf{Lớp:} LỚP 10 \quad -- \quad \textbf{Mã HS:} \texttt{HS013}
	\end{minipage}
	\hfill
	\begin{minipage}[c]{0.22\textwidth}
		\begin{tcolorbox}[colback=white,colframe=black,arc=0mm,boxrule=0.6pt,height=1.8cm,halign=center,valign=top,top=2pt]
			{\bfseries\normalsize Điểm}
		\end{tcolorbox}
	\end{minipage}
	\hfill
	\begin{minipage}[c]{0.30\textwidth}
		\raggedleft
		\textit{Thời gian làm bài: 30 phút}\\[4pt]
		\textbf{Mã đề thi:} \fbox{\textbf{\Large\strut\ 001\ }}
	\end{minipage}
\end{tcolorbox}
\vspace{0.25cm}

\noindent\textbf{\large ĐỀ BÀI (MÃ ĐỀ: 001)}
\vspace{0.1cm}

% --- CÂU 1 (TRANG 1) ---
\begin{ex}
	Thực hiện phép tính:
	\begin{enumerate}[a)]
		\item $A = 3x^3y \cdot 4xy^2$
		\item $B = 3x^2y^3 \cdot 2x^2y - 11x^4y^4$
		\item $C = 18x^6y^5 : 3x^2y^2 - 12x^4y^3$
		\item $D = 4x^3y^2 + 8x^5y^3 : 2x^2y$
	\end{enumerate}
	\loigiai{
		\begin{enumerate}[a)]
			\item Ta có:
			\begin{align*}
				A &= 3x^3y \cdot 4xy^2 \\
				&= (3 \cdot 4) \cdot (x^3 \cdot x) \cdot (y \cdot y^2) \\
				&= 12x^4y^3
			\end{align*}
			\item Ta có:
			\begin{align*}
				B &= 3x^2y^3 \cdot 2x^2y - 11x^4y^4 \\
				&= 6x^4y^4 - 11x^4y^4 \\
				&= (6 - 11)x^4y^4 \\
				&= -5x^4y^4
			\end{align*}
			\item Ta có:
			\begin{align*}
				C &= 18x^6y^5 : 3x^2y^2 - 12x^4y^3 \\
				&= 6x^4y^3 - 12x^4y^3 \\
				&= (6 - 12)x^4y^3 \\
				&= -6x^4y^3
			\end{align*}
			\item Ta có:
			\begin{align*}
				D &= 4x^3y^2 + 8x^5y^3 : 2x^2y \\
				&= 4x^3y^2 + 4x^3y^2 \\
				&= 8x^3y^2
			\end{align*}
		\end{enumerate}
	}
\end{ex}

\newpage
% ----------------------------------------------------------------------
% TRANG 2 (MẶT SAU) -- DÒNG KẺ TIẾP THEO
% ----------------------------------------------------------------------
\noindent
\textbf{Học sinh:} Việt \quad -- \quad \textbf{Mã đề:} \fbox{\textbf{001}} \hfill \textit{Trang 2/2}\par
\vspace{-0.1cm}
\noindent\hrulefill
\vspace{0.3cm}

% --- CÂU 2 (TRANG 2) ---
\begin{ex}
	Thực hiện phép tính:
	\begin{enumerate}[a)]
		\item $A = -5a^2b^3 \cdot 2ab^2$
		\item $B = 4a^3b \cdot 3ab^2 - 17a^4b^3$
		\item $C = 24a^7b^4 : 6a^4b^2 - 9a^3b^2$
		\item $D = 3a^2b^4 + 14a^4b^5 : 2a^2b$
	\end{enumerate}
	\loigiai{
		\begin{enumerate}[a)]
			\item Ta có:
			\begin{align*}
				A &= -5a^2b^3 \cdot 2ab^2 \\
				&= (-5 \cdot 2) \cdot (a^2 \cdot a) \cdot (b^3 \cdot b^2) \\
				&= -10a^3b^5
			\end{align*}
			\item Ta có:
			\begin{align*}
				B &= 4a^3b \cdot 3ab^2 - 17a^4b^3 \\
				&= 12a^4b^3 - 17a^4b^3 \\
				&= (12 - 17)a^4b^3 \\
				&= -5a^4b^3
			\end{align*}
			\item Ta có:
			\begin{align*}
				C &= 24a^7b^4 : 6a^4b^2 - 9a^3b^2 \\
				&= 4a^3b^2 - 9a^3b^2 \\
				&= (4 - 9)a^3b^2 \\
				&= -5a^3b^2
			\end{align*}
			\item Ta có:
			\begin{align*}
				D &= 3a^2b^4 + 14a^4b^5 : 2a^2b \\
				&= 3a^2b^4 + 7a^2b^4 \\
				&= 10a^2b^4
			\end{align*}
		\end{enumerate}
	}
\end{ex}

\newpage

\end{document}
